% Generated mechanically from frozen input p12/ja/build/ch111/exercises_full.tex.
% Do not edit this generated file; edit the bound locale source instead.

% OK, start here.
%

\setcounter{chapter}{110}
\chapter{演習}



\phantomsection
\label{exercises-section-phantom}


\section{代数}
\label{exercises-section-algebra}

\noindent
この最初の節には，雑多な問題をいくつか集める．

\begin{exercise}
\label{exercises-exercise-isomorphism-localizations}
$A$ を環，${\mathfrak m}$ を極大イデアルとする．$A[X]$ において
$\tilde {\mathfrak m}_1 = ({\mathfrak m}, X)$ および
$\tilde {\mathfrak m}_2 = ({\mathfrak m}, X-1)$ とおく．次を示せ：
$$
A[X]_{\tilde {\mathfrak m}_1} \cong A[X]_{\tilde {\mathfrak m}_2}.
$$
\end{exercise}

\begin{exercise}
\label{exercises-exercise-coherent}
% P12-ERR-0415: frozen English compound-modifier defect; see ../control/ERRATA_CANDIDATES.jsonl.
非 Noether 環 $R$ であって，$R$ の任意の有限生成イデアルが
$R$-加群として有限表示となるものの例を挙げよ．
（最後の性質を満たす環を {\it コヒーレント} という．）
\end{exercise}

\begin{exercise}
\label{exercises-exercise-flat-ideals-pid}
$(A, {\mathfrak m}, k)$ を Noether 局所環とする．任意の有限 $A$-加群 $M$ に対し，
$r(M)$ を $M$ の $A$-加群としての最小生成元数と定める．NAK により，この数は
$\dim_k M/{\mathfrak m} M = \dim_k M \otimes_A k$ に等しい．
\begin{enumerate}
\item $r(M \otimes_A N) = r(M)r(N)$ を示せ．
\item $I\subset A $ を $r(I) > 1$ を満たすイデアルとする．
$r(I^2) < r(I)^2$ を示せ．
\item $A$ のすべてのイデアルが平坦加群ならば，$A$ は主イデアル整域
（または体）であることを導け．
\end{enumerate}
\end{exercise}

\begin{exercise}
\label{exercises-exercise-not-isomorphic}
$k$ を体とする．次の各組の $k$-代数が同型でないことを示せ：
\begin{enumerate}
\item 任意の $n\geq 1$ に対する $k[x_1, \ldots, x_n]$ と
$k[x_1, \ldots, x_{n + 1}]$．
\item $k[a, b, c, d, e, f]/(ab + cd + ef)$ と $k[x_1, \ldots, x_n]$
（$n = 5$）．
\item $k[a, b, c, d, e, f]/(ab + cd + ef)$ と $k[x_1, \ldots, x_n]$
（$n = 6$）．
\end{enumerate}
\end{exercise}

\begin{remark}
\label{exercises-remark-simple-geometric}
もちろん，この演習の狙いは各場合に「大きな」定理を適用するのではなく，
簡明な議論を見つけることにある．それでも，一般原理を指針とするのは有益である．
\end{remark}

\begin{exercise}
\label{exercises-exercise-silly}
代数．（些末で，容易なはずである．）
\begin{enumerate}
\item 環 $A$ と，$A$-加群の非分裂短完全列
$$
0 \to M_1 \to M_2 \to M_3 \to 0.
$$
の例を挙げよ．
\item 上のような $A$-加群の非分裂列と，忠実平坦な $A \to B$ であって，
$$
0 \to M_1\otimes_A B \to M_2\otimes_A B \to M_3\otimes_A B \to 0.
$$
が $B$-加群の列として分裂するものの例を挙げよ．
\end{enumerate}
\end{exercise}

\begin{exercise}
\label{exercises-exercise-field-kummer}
$k$ を，1 の原始 $n$ 乗根 $\zeta$ をもつ体とする．すなわち，
$\zeta^n = 1$ であるが，$0 < m < n$ に対して $\zeta^m\not = 1$ である．
\begin{enumerate}
\item $k$ の標数が $n$ と互いに素であることを示せ．
\item
% P12-ERR-0416: frozen English indefinite-article defect; see ../control/ERRATA_CANDIDATES.jsonl.
$a \in k$ が $k$ の元であり，$n \geq d > 1$ を満たす $n$ のどの約数 $d$ に対しても
$k$ における $d$ 乗でないと仮定する．$k[x]/(x^n-a)$ が体であることを示せ．
（ヒント：$x^n-a$ の分解体を考え，Galois 理論を用いよ．）
\end{enumerate}
\end{exercise}

\begin{exercise}
\label{exercises-exercise-valuation}
% P12-ERR-0417: frozen bare min operators at all three loci; see ../control/ERRATA_CANDIDATES.jsonl.
$\nu : k[x]\setminus \{0\}  \to {\mathbf Z}$ を，次の性質をもつ写像とする：
$f$，$g$ が零でないとき $\nu(fg) = \nu(f) + \nu(g)$，
$f$，$g$，$f + g$ が零でないとき $\nu(f + g) \geq min(\nu(f), \nu(g))$，
またすべての $c\in k^*$ に対して $\nu(c) = 0$ である．
\begin{enumerate}
\item $f$，$g$，$f + g$ が零でなく，$\nu(f) \not = \nu(g)$ ならば，
$\nu(f + g) = min(\nu(f), \nu(g))$ が成り立つことを示せ．
\item $f = \sum a_i x^i$，$f\not = 0$ ならば，
$\nu(f) \geq min(\{i\nu(x)\}_{a_i\not = 0})$ であることを示せ．
等号はいつ成り立つか．
\item $\nu$ が負の値をとるならば，ある $n\in {\mathbf N}$ に対して
$\nu(f) = -n \deg(f)$ となることを示せ．
\item $\nu(x) \geq 0$ と仮定する．
$\{f \mid f = 0, \ or\ \nu(f) > 0\}$ が $k[x]$ の素イデアルであることを示せ．
\item 可能なすべての $\nu$ を記述せよ．
\end{enumerate}
\end{exercise}


\noindent
$A$ を環とする．{\it 冪等元} とは，$e^2 = e$ を満たす元 $e \in A$ のことである．
$1$ と $0$ は常に冪等元である．
% P12-ERR-0418: frozen zero-only nontrivial-idempotent definition; see ../control/ERRATA_CANDIDATES.jsonl.
{\it 非自明な冪等元} とは，零に等しくない冪等元のことである．
二つの冪等元 $e, e' \in A$ が {\it 直交する} とは，$ee' = 0$ となることをいう．

\begin{exercise}
\label{exercises-exercise-product}
$A$ を環とする．$A$ が二つの零でない環の積であるための必要十分条件は，
$A$ が非自明な冪等元をもつことであることを示せ．
\end{exercise}

\begin{exercise}
\label{exercises-exercise-lift-idempotents}
$A$ を環，$I \subset A$ を局所冪零イデアルとする．
写像 $A \to A/I$ が冪等元の集合上の全単射を誘導することを示せ．
（ヒント：$I$ が冪零の場合を先に証明するとよい．その後，「絶対 Noether 簡約」を用いて
冪零の場合へ帰着せよ．）
\end{exercise}

\section{余極限}
\label{exercises-section-colimits}


\begin{definition}
\label{exercises-definition-directed-poset}
{\it 有向集合} とは，前順序 $\leq$ を備えた空でない集合 $I$ であって，
任意の組 $i, j \in I$ に対し，$i \leq k$ かつ $j \leq k$ を満たす
$k \in I$ が存在するものをいう．$I$ 上の {\it 環の系} とは，
各 $i \in I$ に対する環 $A_i$ と，$i \leq j$ のときの環準同型
$\varphi_{ij} : A_i \to A_j$ であって，$i \leq j \leq k$ のとき合成
$A_i \to A_j \to A_k$ が $A_i \to A_k$ に等しくなるものからなる．
\end{definition}

\noindent
同様にして，群の系，固定した環上の加群の系，体上のベクトル空間の系などを定義する．

\begin{exercise}
\label{exercises-exercise-directed-colimit}
$I$ を有向集合，$(A_i, \varphi_{ij})$ を $I$ 上の環の系とする．
環 $A$ と写像 $\varphi_i : A_i \to A$ であって，すべての $i \leq j$ に対し
$\varphi_j \circ \varphi_{ij} = \varphi_i$ となり，さらに次の普遍性をもつものが
存在することを示せ：任意の環 $B$ と，すべての $i \leq j$ に対して
$\psi_j \circ \varphi_{ij} = \psi_i$ を満たす写像 $\psi_i : A_i \to B$ が
与えられたとき，$\psi_i = \psi \circ \varphi_i$ を満たす一意な環準同型
$\psi : A \to B$ が存在する．
\end{exercise}

\begin{definition}
\label{exercises-definition-colimit}
演習 \ref{exercises-exercise-directed-colimit} で構成した環 $A$ を，この系の
{\it 余極限} という．記法は $\colim A_i$ とする．
\end{definition}

\begin{exercise}
\label{exercises-exercise-prime-in-colimit}
% P12-ERR-0419: first frozen reversed-order declaration; see ../control/ERRATA_CANDIDATES.jsonl.
$(I, \geq)$ を有向集合，$(A_i, \varphi_{ij})$ を $I$ 上の環の系とし，
その余極限を $A$ とする．次の全単射が存在することを証明せよ：
$$
\Spec(A) = \{(\mathfrak p_i)_{i \in I} \mid
\mathfrak p_i \subset A_i \text{ かつ }
\mathfrak p_i = \varphi_{ij}^{-1}(\mathfrak p_j)\ \forall i \leq j\}
\subset \prod\nolimits_{i \in I} \Spec(A_i)
$$
等号の右辺の集合は，集合 $\Spec(A_i)$ の逆極限である．
記法は $\lim \Spec(A_i)$ とする．
\end{exercise}

\begin{exercise}
\label{exercises-exercise-colimit-surjective}
% P12-ERR-0419: second frozen reversed-order declaration; see ../control/ERRATA_CANDIDATES.jsonl.
$(I, \geq)$ を有向集合，$(A_i, \varphi_{ij})$ を $I$ 上の環の系とし，
その余極限を $A$ とする．すべての $i \leq j$ に対して
$\Spec(A_j) \to \Spec(A_i)$ が全射であると仮定する．
このとき，すべての $i$ に対して $\Spec(A) \to \Spec(A_i)$ が
全射であることを示せ．
（ヒント：Tychonoff の定理を用いてもよいが，代数，補題
\ref{algebra-lemma-in-image} に基づく本質的に自明な直接の代数的証明もある．）
\end{exercise}

\begin{exercise}
\label{exercises-exercise-integral-colimit-finite}
% P12-ERR-0420: frozen broken English parallel syntax; see ../control/ERRATA_CANDIDATES.jsonl.
$A \subset B$ を整環拡大とする．$\Spec(B) \to \Spec(A)$ が全射であることを証明せよ．
上の演習，有限環拡大の場合にはこの主張が成り立つという事実
（講義で証明した），および $B = \colim B_i$ が有限拡大 $A \subset B_i$ の
有向余極限であることの証明を用いよ．
\end{exercise}

\begin{exercise}
\label{exercises-exercise-colimit-tensor}
$(I, \geq)$ を有向集合とする．$A$ を環とし，$(N_i, \varphi_{i, i'})$ を
$I$ で添字づけられた $A$-加群の有向系とする．$M$ を別の $A$-加群とする．
次を証明せよ：
$$
\colim_{i\in I} M \otimes_A N_i\cong
M \otimes_A \Big( \colim_{i\in I} N_i\Big).
$$
\end{exercise}

\begin{definition}
\label{exercises-definition-finite-presentation}
$R$ 上の加群 $M$ が $R$ 上 {\it 有限表示} であるとは，有限自由加群の写像
$ R^{\oplus n} \to R^{\oplus m}$ の余核に同型であることをいう．
\end{definition}

\begin{exercise}
\label{exercises-exercise-colimit-modules}
任意の環上の任意の加群が次のように表されることを証明せよ：
\begin{enumerate}
\item その有限生成部分加群の余極限，
\item 何らかの仕方による有限表示加群の余極限．
\end{enumerate}
\end{exercise}

\section{加法圏とアーベル圏}
\label{exercises-section-additive}

\begin{exercise}
\label{exercises-exercise-filtered-vector-spaces}
$k$ を体とする．$\mathcal{C}$ を $k$ 上のフィルター付きベクトル空間の圏とする．
任意の圏のフィルター付き対象の定義についてはホモロジー，定義
\ref{homology-definition-filtered} を参照せよ．
\begin{enumerate}
\item これが加法圏であることを示せ（二つの対象の直和が何であるかを
注意深く説明せよ）．
\item $f : (V, F) \to (W, F)$ を $\mathcal{C}$ の射とする．
$f$ が核と余核をもつことを示せ（$f$ の核と余核が何であるかを
正確に説明せよ）．
\item 標準写像 $\Coim(f) \to \Im(f)$ が同型でないような
$\mathcal{C}$ の写像の例を挙げよ．
\end{enumerate}
\end{exercise}

\begin{exercise}
\label{exercises-exercise-torsion-free}
$R$ を Noether 整域とする．$\mathcal{C}$ を有限生成捩れなし $R$-加群の圏とする．
\begin{enumerate}
\item これが加法圏であることを示せ．
\item $f : N \to M$ を $\mathcal{C}$ の射とする．
$f$ が核と余核をもつことを示せ（$f$ の核と余核を必ず正確に定義せよ）．
\item 標準写像 $\Coim(f) \to \Im(f)$ が同型でないような Noether 整域 $R$ と
$\mathcal{C}$ の写像の例を挙げよ．
\end{enumerate}
\end{exercise}

\begin{exercise}
\label{exercises-exercise-other}
加法的で核と余核をもつが，演習
\ref{exercises-exercise-filtered-vector-spaces} および
\ref{exercises-exercise-torsion-free} のものとは異なる圏の例を挙げよ．
\end{exercise}

\section{テンソル積}
\label{exercises-section-tensor-product}

\noindent
テンソル積は代数，節
\ref{algebra-section-tensor-product} で導入した．
$R$ を環とし，$\text{Mod}_R$ を $R$-加群の圏とする．
関手 $F : \text{Mod}_R \to \text{Mod}_R$ について，次のようにいう：
\begin{enumerate}
\item 任意の $R$-加群 $M, N$ に対して
$F : \Hom_R(M, N) \to \Hom_R(F(M), F(N))$
がアーベル群の準同型であるとき，加法的であるという．
ホモロジー，定義 \ref{homology-definition-preadditive} を参照せよ．
% P12-ERR-0421: frozen missing-predicate defect; see ../control/ERRATA_CANDIDATES.jsonl.
\item 任意の $R$-加群 $M, N$ に対して
$F : \Hom_R(M, N) \to \Hom_R(F(M), F(N))$ が $R$-線形であるとき，
$R$-線形であるという．
\item 任意の短完全列
$0 \to M_1 \to M_2 \to M_3 \to 0$ に対して列
$F(M_1) \to F(M_2) \to F(M_3) \to 0$ が完全であるとき，
右完全であるという．
\item 任意の短完全列
$0 \to M_1 \to M_2 \to M_3 \to 0$ に対して列
$0 \to F(M_1) \to F(M_2) \to F(M_3)$ が完全であるとき，
左完全であるという．
\item 集合 $I$ と $R$-加群 $M_i$ が与えられたとき，写像
$F(M_i) \to F(\bigoplus M_i)$ が同型
$\bigoplus F(M_i) = F(\bigoplus M_i)$ を誘導するならば，
直和と可換するという．
\end{enumerate}

\begin{exercise}
\label{exercises-exercise-characterize-tensor-functor}
$R$ を環とし，上の記法を用いる．
\begin{enumerate}
\item $R$-線形でない加法的関手
$F : \text{Mod}_R \to \text{Mod}_R$ をもつ環 $R$ の例を挙げよ．
\item $N$ を $R$-加群とする．関手
$F(M) = M \otimes_R N$ が $R$-線形かつ右完全であり，
直和と可換することを示せ．
\item 逆に，$R$-線形かつ右完全で，直和と可換する任意の関手
$F : \text{Mod}_R \to \text{Mod}_R$ は，ある $R$-加群 $N$ によって
$F(M) = M \otimes_R N$ の形に表されることを示せ．
\item (3) において，$F$ が直和と可換するという仮定を外すと，
結論はもはや成り立たないことを示せ．
\end{enumerate}
\end{exercise}





\section{平坦な環準同型}
\label{exercises-section-flat}

\begin{exercise}
\label{exercises-exercise-localization-flat}
$S$ を環 $A$ の乗法的部分集合とする．
\begin{enumerate}
\item $A$-加群 $M$ に対して $S^{-1}M = S^{-1}A \otimes_A M$ を示せ．
\item $S^{-1}A$ が $A$ 上平坦であることを示せ．
\end{enumerate}
\end{exercise}

\begin{exercise}
\label{exercises-exercise-examples-not-flat}
次の各場合に，$A$-加群の単射 $M_1 \to M_2$ であって
$M_1\otimes N \to M_2 \otimes N$ が単射でないものを見つけよ：
\begin{enumerate}
\item $A = k[x, y]$ かつ $N = (x, y) \subset A$．（ここでも以下でも $k$ は体とする．）
\item $A = k[x, y]$ かつ $N = A/(x, y)$．
\end{enumerate}
\end{exercise}

\begin{exercise}
\label{exercises-exercise-flat-not-projective}
% P12-ERR-0422: frozen mismatched-article coordination; see ../control/ERRATA_CANDIDATES.jsonl.
環 $A$ と，有限 $A$-加群 $M$ であって，
平坦だが射影的でない $A$-加群であるものの例を挙げよ．
\end{exercise}

\begin{remark}
\label{exercises-remark-flat-not-projective}
$M$ が有限表示であり，$A$ 上平坦ならば，$M$ は $A$ 上射影的である．
したがって，求める例では環 $A$ は Noether 環でない必要がある．
私は，${\mathbf R}$ 上の ${\mathcal C}^\infty$-函数の環を $A$ とする例を知っている．
\end{remark}

\begin{exercise}
\label{exercises-exercise-flat-not-free-dvr}
${\mathbf Z}_{(2)}$ 上の平坦だが自由でない加群を見つけよ．
\end{exercise}

\begin{exercise}
\label{exercises-exercise-flat-deformations}
平坦変形．
\begin{enumerate}
\item $k$ を体とし，$k[\epsilon]$ を
双対数環 $k[\epsilon] = k[x]/(x^2)$ とし，$\epsilon = \bar x$ とする．
任意の $k$-代数 $A$ に対して，$A$ が $B/\epsilon B$ と同型になるような
平坦 $k[\epsilon]$-代数 $B$ が存在することを示せ．
\item $k = {\mathbf F}_p = {\mathbf Z}/p{\mathbf Z}$ とし，
$$
A = k[x_1, x_2, x_3, x_4, x_5, x_6]/
(x_1^p, x_2^p, x_3^p, x_4^p, x_5^p, x_6^p).
$$
$B/pB$ が $A$ と同型になるような平坦
${\mathbf Z}/p^2{\mathbf Z}$-代数 $B$ が存在することを示せ．
（ここで $p$ は $\epsilon$ の役割を果たす．）
\item 今度は $p = 2$ とし，
$k = {\mathbf F}_2 = {\mathbf Z}/2{\mathbf Z}$ および
$$
A = k[x_1, x_2, x_3, x_4, x_5, x_6]/
(x_1^2, x_2^2, x_3^2, x_4^2, x_5^2, x_6^2, x_1x_2 + x_3x_4 + x_5x_6).
$$
に対して同じ問題を考える．しかし，この場合には，$B/2B$ が $A$ と同型になるような
平坦 ${\mathbf Z}/4{\mathbf Z}$-代数 $B$ は {\it 存在しない} ことを示せ．
（仕掛けを見抜け！ 同じ例は任意の標数 $p > 0$ で成り立つが，
計算はより難しくなる．）
\end{enumerate}
\end{exercise}

\begin{exercise}
\label{exercises-exercise-flat-given-residue-field-extension}
$(A, {\mathfrak m}, k)$ を局所環とし，$k'/k$ を有限次体拡大とする．
局所環の平坦な局所準同型 $A \to B$ であって，
${\mathfrak m}_B = {\mathfrak m} B$ となり，かつ
$B/{\mathfrak m} B$ が $k$-代数として $k'$ と同型になるものが存在することを示せ．
（ヒント：まず $k \subset k'$ が一つの元で生成される場合を扱え．）
\end{exercise}

\begin{remark}
\label{exercises-remark-flat-given-residue-field-extension-general}
同じ結果は任意の体拡大 $K/k$ に対して成り立つ．
\end{remark}








\section{環のスペクトル}
\label{exercises-section-spectrum-ring}

\begin{exercise}
\label{exercises-exercise-spec-Z}
$\Spec(\mathbf{Z})$ を集合として計算し，その位相を記述せよ．
\end{exercise}

\begin{exercise}
\label{exercises-exercise-basis-opens-standard}
$A$ を任意の環とする．$f\in A$ に対して
$D(f):= \{\mathfrak p \subset A \mid f \not \in \mathfrak p\}$ と定める．
開集合 $D(f)$ が $\Spec(A)$ の位相の基底をなすことを証明せよ．
\end{exercise}

\begin{exercise}
\label{exercises-exercise-radical-ideals-closed}
写像 $I\mapsto V(I)$ が自然な全単射
$$
\{I\subset A\text{ で }I = \sqrt{I}\}
\longrightarrow
\{T\subset \Spec(A)\text{ は閉集合}\}
$$
を定めることを証明せよ．
\end{exercise}

\begin{definition}
\label{exercises-definition-quasi-compact}
位相空間 $X$ が {\it 準コンパクト} であるとは，任意の開被覆
$X = \bigcup_{i\in I} U_i$ に対して，有限部分集合
$\{i_1, \ldots, i_n\}\subset I$ であって
$X = U_{i_1}\cup\ldots U_{i_n}$ となるものが存在することをいう．
\end{definition}

\begin{exercise}
\label{exercises-exercise-spec-quasi-compact}
任意の環 $A$ に対して $\Spec(A)$ が準コンパクトであることを証明せよ．
\end{exercise}

\begin{definition}
\label{exercises-definition-Hausdorff}
位相空間 $X$ が分離公理 $T_0$ を満たすとは，任意の相異なる二点
$x, y\in X$，$x\not = y$ に対して，その一方を含み他方を含まない
$X$ の開部分集合が存在することをいう．
$X$ が {\it Hausdorff} であるとは，任意の相異なる二点
$x, y\in X$，$x\not = y$ に対して，$x\in U$ および
$y\in V$ を満たす互いに素な開部分集合 $U, V$ が存在することをいう．
\end{definition}

\begin{exercise}
\label{exercises-exercise-not-hausdorff}
$\Spec(A)$ が一般には {\bf Hausdorff でない} ことを示せ．
$\Spec(A)$ が $T_0$ であることを証明せよ．$T_0$ でない位相空間
$X$ の例を挙げよ．
\end{exercise}

\begin{remark}
\label{exercises-remark-not-hausdorff}
通常，コンパクトという語は，準コンパクトかつ Hausdorff な空間に対して用いる．
\end{remark}

\begin{definition}
\label{exercises-definition-irreducible}
位相空間 $X$ が {\it 既約} であるとは，$X$ が空でなく，閉部分集合
$Z_1, Z_2\subset X$ によって $X = Z_1\cup Z_2$ と書けるとき，
$Z_1 = X$ または $Z_2 = X$ のいずれかが成り立つことをいう．
位相空間の部分集合 $T\subset X$ が {\it 既約} であるとは，
$X$ から誘導される位相に関して既約位相空間であることをいう．
この定義から，$T$ が既約であるための必要十分条件は，
$X$ における $T$ の閉包 $\bar T$ が既約であることである．
\end{definition}

\begin{exercise}
\label{exercises-exercise-irreducible-spec}
$\Spec(A)$ が既約であるための必要十分条件は
$Nil(A)$ が素イデアルであることであり，この場合それが
$A$ の一意な極小素イデアルであることを証明せよ．
\end{exercise}

\begin{exercise}
\label{exercises-exercise-irreducible-prime}
閉部分集合 $T\subset \Spec(A)$ が既約であるための必要十分条件は，
ある素イデアル ${\mathfrak p}\subset A$ に対して
$T = V({\mathfrak p})$ と書けることであることを証明せよ．
\end{exercise}

\begin{definition}
\label{exercises-definition-generic-point}
既約位相空間 $X$ の点 $x$ が $X$ の {\it 生成点} であるとは，
$X$ が部分集合 $\{x\}$ の閉包に等しいことをいう．
\end{definition}

\begin{exercise}
\label{exercises-exercise-irreducible-T0-at-most-one-generic}
$T_0$ 空間 $X$ において，任意の既約閉部分集合が高々一つの
生成点しかもたないことを示せ．
\end{exercise}

\begin{exercise}
\label{exercises-exercise-spec-sober}
% P12-ERR-0423: frozen undefined-X codomain; see ../control/ERRATA_CANDIDATES.jsonl.
$\Spec(A)$ では，任意の既約閉部分集合が生成点を {\it 実際にもつ} ことを証明せよ．
実際，写像
${\mathfrak p} \mapsto \overline{\{{\mathfrak p}\}}$ が，
$\Spec(A)$ と $X$ の既約閉部分集合全体との間の全単射であることを示せ．
\end{exercise}

\begin{exercise}
\label{exercises-exercise-irreducible-subset-not-generic}
$\Spec(\mathbf{Z})$ の既約部分集合が必ずしも生成点をもたないことを示す
例を挙げよ．
\end{exercise}

\begin{definition}
\label{exercises-definition-Noetherian-space}
位相空間 $X$ が {\it Noether} であるとは，$X$ の閉部分集合の任意の減少列
$Z_1\supset Z_2 \supset Z_3\supset \ldots$ が停留することをいう．
（閉部分集合の任意の増加列が停留するとき，{\it Artin} であるという．）
\end{definition}

\begin{exercise}
\label{exercises-exercise-Noetherian-spec}
環 $A$ が Noether ならば，位相空間 $\Spec(A)$ が Noether であることを示せ．
逆が偽であることを示す例を挙げよ．（望むなら Artin の場合についても同様にせよ．）
\end{exercise}

\begin{definition}
\label{exercises-definition-irreducible-component}
極大な既約部分集合 $T\subset X$ を，空間 $X$ の
{\it 既約成分} という．$X$ のそのような既約成分は自動的に
$X$ の閉部分集合となる．
\end{definition}

\begin{exercise}
\label{exercises-exercise-irreducible-in-irreducible}
$X$ の任意の既約部分集合が $X$ のある既約成分に含まれることを証明せよ．
\end{exercise}

\begin{exercise}
\label{exercises-exercise-Noetherian-finite-nr-irreducible}
Noether 位相空間 $X$ が有限個の既約成分，たとえば
$X_1, \ldots, X_n$ のみをもち，
$X = X_1\cup X_2\cup\ldots\cup X_n$ となることを証明せよ．
（任意の $X$ は常にその既約成分の合併であることに注意せよ．
しかし，たとえば $X = {\mathbf R}$ に通常の位相を入れると，
$X$ の既約成分は一点部分集合である．これはあまり興味深くない．）
\end{exercise}

\begin{exercise}
\label{exercises-exercise-irreducible-components-minimal-primes}
$\Spec(A)$ の既約成分が $A$ の極小素イデアルに対応することを示せ．
\end{exercise}

\begin{definition}
\label{exercises-definition-closed}
点 $x\in X$ が {\it 閉点} であるとは，$\overline{\{x\}} = \{ x\}$ となることをいう．
$x, y$ を $X$ の点とする．$x\in \overline{\{y\}}$ であるとき，
$x$ は $y$ の {\it 特殊化} である，または $y$ は $x$ の
{\it 一般化} であるという．
\end{definition}

\begin{exercise}
\label{exercises-exercise-closed-maximal}
$\Spec(A)$ の閉点が $A$ の極大イデアルに対応することを示せ．
\end{exercise}

\begin{exercise}
\label{exercises-exercise-generalization}
${\mathfrak p}$ が $\Spec(A)$ における ${\mathfrak q}$ の一般化であるための
必要十分条件は，${\mathfrak p}\subset {\mathfrak q}$ であることを示せ．
閉点，極大イデアル，生成点，極小素イデアルを，一般化と特殊化を用いて特徴づけよ．
% P12-ERR-0424: frozen singular/plural generic-point mismatch; see ../control/ERRATA_CANDIDATES.jsonl.
（ここでは，既約とは限らない位相空間 $X$ の点が，$X$ の既約成分の一つの
生成点であるとき，それを生成点と呼ぶという用語法を採用する．）
\end{exercise}

\begin{exercise}
\label{exercises-exercise-disjoint-closed-spec}
$I$ と $J$ を $A$ のイデアルとする．
$V(I)$ と $V(J)$ が互いに素であるための条件は何か．
\end{exercise}

\begin{definition}
\label{exercises-definition-connected-component}
位相空間 $X$ が {\it 連結} であるとは，空でなく，二つの空でない互いに素な
開部分集合の合併でないことをいう．$X$ の極大な連結部分集合を
{\it 連結成分} という．$X$ の任意の点は $X$ のある連結成分に含まれ，
$X$ の任意の連結成分は $X$ において閉じている．
（しかし，一般に連結成分は $X$ において開であるとは限らない．）
\end{definition}

\begin{exercise}
\label{exercises-exercise-disconnected-spec}
$A$ を零でない環とする．
$\Spec(A)$ が非連結であるための必要十分条件は，ある零でない環 $B, C$ に対して
$A\cong B \times C$ となることであることを示せ．
\end{exercise}

\begin{exercise}
\label{exercises-exercise-connected-component-stable-generalization}
$T$ を $\Spec(A)$ の連結成分とする．$T$ が一般化のもとで安定であることを証明せよ．
$A$ が Noether ならば，$T$ が $\Spec(A)$ の開部分集合であることを証明せよ．
（注意：たとえば $A$ が ${\mathbf F}_2$ のコピーの無限積であるとき，
これは誤りである．この環のスペクトルは無限個の閉点からなる．）
\end{exercise}

\begin{exercise}
\label{exercises-exercise-primes-kx}
$\Spec(k[x])$ を計算せよ．すなわち，この環の素イデアル，可能な特殊化，
および位相を記述せよ．（$k$ が代数閉体の場合だけでなく，
$k$ が代数閉体でない場合についても求めよ．）
\end{exercise}

\begin{exercise}
\label{exercises-exercise-primes-kxy}
$k$ を代数閉体とし，$\Spec(k[x, y])$ を計算せよ．
［ヒント：射
$\varphi : \Spec(k[x, y]) \to \Spec(k[x])$ を用いよ．
$\varphi({\mathfrak p}) = (0)$ ならば，
$S = \{f\in k[x] \mid f \not = 0\}$ に関して局所化し，
局所化と $\Spec$ に関する講義の結果を用いよ．］
（なぜ代数幾何学者はこれをアフィン 2-空間と呼ぶのだろうか．）
\end{exercise}

\begin{exercise}
\label{exercises-exercise-primes-Zy}
$\Spec(\mathbf{Z}[y])$ を計算せよ．
［ヒント：上と同様である．］（$\mathbf{Z}$ 上のアフィン 1-空間．）
\end{exercise}






\section{局所化}
\label{exercises-section-localization}


\begin{exercise}
\label{exercises-exercise-submodule-localization}
$A$ を環，$S \subset A$ を乗法的部分集合とする．
$M$ を $A$-加群，$N \subset S^{-1}M$ を $S^{-1}A$-部分加群とする．
$N = S^{-1}N'$ となる $A$-部分加群 $N' \subset M$ が存在することを示せ．
（この有用な結果は，特に $S^{-1}A$ のイデアルに適用できる．）
\end{exercise}

\begin{exercise}
\label{exercises-exercise-localize-zero}
$A$ を環，$M$ を $A$-加群，$m \in M$ とする．
\begin{enumerate}
\item $I = \{a \in A \mid am = 0\}$ が $A$ のイデアルであることを示せ．
\item $A$ の素イデアル $\mathfrak p$ に対して，$M_\mathfrak p$ における
$m$ の像が零であるための必要十分条件は $I \not \subset \mathfrak p$ であることを示せ．
\item $m$ が零であるための必要十分条件は，$A$ のすべての素イデアル
$\mathfrak p$ に対して $M_\mathfrak p$ における $m$ の像が零であることを示せ．
\item $m$ が零であるための必要十分条件は，$A$ のすべての極大イデアル
$\mathfrak m$ に対して $M_\mathfrak m$ における $m$ の像が零であることを示せ．
\item $M = 0$ であるための必要十分条件は，すべての極大イデアル
$\mathfrak m$ に対して $M_{\mathfrak m}$ が零であることを示せ．
\end{enumerate}
\end{exercise}

\begin{exercise}
\label{exercises-exercise-localization-is-field}
$A$ が相異なる素イデアルを三つ以上もつ整域であり，$f \in A$ が
主局所化 $A_f = \{1, f, f^2, \ldots \}^{-1}A$ を体にする元であるような
組 $(A, f)$ を見つけよ．
\end{exercise}

\begin{exercise}
\label{exercises-exercise-localize-finite-module-zero}
$A$ を環，$M$ を有限 $A$-加群，$S \subset A$ を乗法的集合とする．
$S^{-1}M = 0$ と仮定する．主局所化
$M_f = \{1, f, f^2, \ldots \}^{-1}M$ が零となる $f \in S$ が
存在することを示せ．
\end{exercise}

\begin{exercise}
\label{exercises-exercise-localization-is-quotient}
$A$ が環，$0 \not = I \not = A$ が零でない真のイデアル，
$S \subset A$ が乗法的部分集合であって，
$A$-代数として $A/I \cong S^{-1}A$ となる三つ組 $(A, I, S)$ の例を挙げよ．
\end{exercise}





\section{Nakayama の補題}
\label{exercises-section-nakayama}

\begin{exercise}
\label{exercises-exercise-nakayama}
$A$ を環とする．
$I$ を $A$ のイデアルとする．
$M$ を $A$-加群とする．
$x_1, \ldots, x_n \in M$ とする．
次を仮定する：
\begin{enumerate}
\item $M/IM$ は $x_1, \ldots, x_n$ で生成される，
\item $M$ は有限 $A$-加群である，
\item $I$ は $A$ のすべての極大イデアルに含まれる．
\end{enumerate}
$x_1, \ldots, x_n$ が $M$ を生成することを示せ．
（解法の提案：演習 \ref{exercises-exercise-localize-zero} と局所化の完全性を用いて，
$A$ の極大イデアルにおける局所化へ帰着せよ．
次に，$M$ を $x_1, \ldots, x_n$ で生成される部分加群で割った商を考え，
講義における Nakayama の補題の主張へ帰着せよ．）
\end{exercise}






\section{長さ}
\label{exercises-section-length}

\begin{definition}
\label{exercises-definition-length}
% P12-ERR-0425: frozen undefined-R module designation; see ../control/ERRATA_CANDIDATES.jsonl.
$A$ を環，$M$ を $A$-加群とする．$M$ の $R$-加群としての
{\it 長さ} とは，
$$
\text{length}_A(M)
=
\sup
\{
n
\mid
\exists\ 0 = M_0 \subset M_1 \subset \ldots \subset M_n = M,
\text{ }M_i \not = M_{i + 1}
\}.
$$
のことである．言い換えれば，部分加群の鎖の長さの上限である．
\end{definition}

\begin{exercise}
\label{exercises-exercise-length-is-one}
環 $A$ 上の加群 $M$ の長さが $1$ であるための必要十分条件は，
$A$ のある極大イデアル $\mathfrak m$ に対して
$A/\mathfrak m$ と同型であることを示せ．
\end{exercise}

\begin{exercise}
\label{exercises-exercise-length-easy}
次の各環上の各加群の長さを計算せよ．答えを簡潔に（！）説明せよ．
（短完全列における長さ函数の加法性を自由に用いてよい．
代数，補題 \ref{algebra-lemma-length-additive} を参照せよ．）
\begin{enumerate}
\item $\mathbf{Z}$ 上の $\mathbf{Z}/120\mathbf{Z}$ の長さ．
\item $\mathbf{C}[x]$ 上の $\mathbf{C}[x]/(x^{100} + x + 1)$ の長さ．
\item $\mathbf{R}[x]$ 上の $\mathbf{R}[x]/(x^4 + 2x^2 + 1)$ の長さ．
\end{enumerate}
\end{exercise}

\begin{exercise}
\label{exercises-exercise-compute-length}
$A = k[x, y]_{(x, y)}$ を原点におけるアフィン平面の局所環とする．
体 $k$ については望む仮定を置いてよい．
$f = x^3 + x^2y^2 + y^{100}$ および $g = y^3 - x^{999}$ とする．
$A$-加群としての $A/(f, g)$ の長さはいくつか．
（進め方の一例：さまざまな $n$ に対し，$A/{\mathfrak m}_A^n=
k[x, y]/(x, y)^n$ の形の商において $f$ と $g$ が生成するイデアルを考えよ．
% P12-ERR-0426: frozen ungrouped quotient-plus-power formula; see ../control/ERRATA_CANDIDATES.jsonl.
$A/(f, g)+{\mathfrak m}_A^n \cong A/(f, g)+{\mathfrak m}_A^{n + 1}$
となる $n$ を見つけ，NAK を用いよ．）
\end{exercise}






\section{随伴素イデアル}
\label{exercises-section-ass}

\noindent
随伴素イデアルは代数，節 \ref{algebra-section-ass} で論じる．

\begin{exercise}
\label{exercises-exercise-compute-ass}
次の各加群について，随伴素イデアルの集合を計算せよ．
\begin{enumerate}
\item $R = k[x, y]$ かつ $M = R/(xy(x + y))$，
\item $R = \mathbf{Z}[x]$ かつ $M = R/(300x + 75)$，
\item $R = k[x, y, z]$ かつ $M = R/(x^3, x^2y, xz)$．
\end{enumerate}
ここでも通常どおり $k$ は体とする．
\end{exercise}

\begin{exercise}
\label{exercises-exercise-prime-power-not-primary}
Noether 環 $R$ と素イデアル $\mathfrak p$ であって，
$\mathfrak p$ が $R/\mathfrak p^2$ の唯一の随伴素イデアルではないものの
例を挙げよ．
\end{exercise}

\begin{exercise}
\label{exercises-exercise-product-primes-not-only-associated}
$R$ を，包含関係にない素イデアル $\mathfrak p$，$\mathfrak q$ をもつ
Noether 環とする．すなわち，$\mathfrak p \not \subset \mathfrak q$ かつ
$\mathfrak q \not \subset \mathfrak p$ とする．
\begin{enumerate}
\item $N = R/(\mathfrak p \cap \mathfrak q)$ に対して
$\text{Ass}(N) = \{\mathfrak p, \mathfrak q\}$ となることを示せ．
\item 加群 $M = R/\mathfrak p \mathfrak q$ が
$\mathfrak p$ とも $\mathfrak q$ とも異なる随伴素イデアルをもちうることを，
例によって示せ．
\end{enumerate}
\end{exercise}





\section{Ext 群}
\label{exercises-section-ext}

\noindent
Ext 群は代数，節 \ref{algebra-section-ext} で定義する．

\begin{exercise}
\label{exercises-exercise-compute-ext-abelian-groups}
与えられた加群のすべての Ext 群 $\Ext^i(M, N)$ を，
$\mathbf{Z}$-加群の圏（アーベル群の圏とも呼ばれる）で計算せよ．
\begin{enumerate}
\item $M = \mathbf{Z}$ かつ $N = \mathbf{Z}$，
\item $M = \mathbf{Z}/4\mathbf{Z}$ かつ $N = \mathbf{Z}/8\mathbf{Z}$，
\item $M = \mathbf{Q}$ かつ $N = \mathbf{Z}/2\mathbf{Z}$，
\item $M = \mathbf{Z}/2\mathbf{Z}$ かつ $N = \mathbf{Q}/\mathbf{Z}$．
\end{enumerate}
\end{exercise}

\begin{exercise}
\label{exercises-exercise-compute-in-regular}
$R = k[x, y]$ とし，$k$ を体とする．
\begin{enumerate}
\item Koszul 複体
$$
0 \to R
\xrightarrow{
\left(
\begin{matrix}
y \\
-x
\end{matrix}
\right)
}
R^{\oplus 2} \xrightarrow{(x, y)} R \xrightarrow{f \mapsto f(0, 0)} k \to 0
$$
が完全であることを手計算で示せ．
\item $k = R/(x, y)$ を $R$-加群とみなし，$\Ext^i_R(k, k)$ を計算せよ．
\end{enumerate}
\end{exercise}

\begin{exercise}
\label{exercises-exercise-infinitely-many-nonzero-ext}
すべての $i \geq 0$ に対して $\Ext^i_R(M, N)$ が零でないような，
Noether 環 $R$ と有限加群 $M$，$N$ の例を挙げよ．
\end{exercise}

\begin{exercise}
\label{exercises-exercise-infinite-ext}
$\Ext^1_R(R/I, R/I)$ が有限 $R$-加群でないような環 $R$ と
イデアル $I$ の例を挙げよ．
（$R$ が Noether ならば，代数，補題
\ref{algebra-lemma-ext-noetherian} により，これは起こりえないことがわかっている．）
\end{exercise}







\section{深さ}
\label{exercises-section-depth}

\noindent
深さは代数，節 \ref{algebra-section-depth} で定義し，
双対化複体，節 \ref{dualizing-section-depth} でさらに研究する．

\begin{exercise}
\label{exercises-exercise-compute-depth}
$R$ を環，$I \subset R$ をイデアル，$M$ を $R$-加群とする．
次の各場合に $\text{depth}_I(M)$ を計算せよ．
\begin{enumerate}
\item $R = \mathbf{Z}$，$I = (30)$，$M = \mathbf{Z}$，
\item $R = \mathbf{Z}$，$I = (30)$，$M = \mathbf{Z}/(300)$，
\item $R = \mathbf{Z}$，$I = (30)$，$M = \mathbf{Z}/(7)$，
\item $R = k[x, y, z]/(x^2 + y^2 + z^2)$，$I = (x, y, z)$，$M = R$，
\item $R = k[x, y, z, w]/(xz, xw, yz, yw)$，$I = (x, y, z, w)$，
$M = R$．
\end{enumerate}
ここで $k$ は体とする．最後の二つの場合には，極大イデアル $I$ において
自由に局所化してよい．
\end{exercise}

\begin{exercise}
\label{exercises-exercise-depth-not-inherited-localization}
深さ $\geq 1$ の Noether 局所環 $(R, \mathfrak m, \kappa)$ と，
次を満たす素イデアル $\mathfrak p$ の例を挙げよ：
\begin{enumerate}
\item $\text{depth}_\mathfrak m(R) \geq 1$，
\item $\text{depth}_\mathfrak p(R_\mathfrak p) = 0$，
\item $\dim(R_\mathfrak p) \geq 1$．
\end{enumerate}
(3) を要求しなければ，この演習は易しすぎる．なぜか．
\end{exercise}

\begin{exercise}
\label{exercises-exercise-depth-torsion-free}
$(R, \mathfrak m)$ を Noether 局所整域とし，$M$ を有限 $R$-加群とする．
\begin{enumerate}
\item $M$ が捩れなしならば，$M$ の $R$ 上の深さが少なくとも $1$ であることを示せ．
\item 深さが $1$ に等しい例を挙げよ．
\end{enumerate}
\end{exercise}

\begin{exercise}
\label{exercises-exercise-depth-examples}
任意の $m \geq n \geq 0$ に対して，$\dim(R) = m$ かつ
$\text{depth}(R) = n$ となる Noether 局所環 $R$ の例を挙げよ．
\end{exercise}

\begin{exercise}
\label{exercises-exercise-make-depth-1}
$(R, \mathfrak m)$ を Noether 局所環とし，$M$ を有限 $R$-加群とする．
次の性質をもつ標準的な短完全列
$$
0 \to K \to M \to Q \to 0
$$
が存在することを示せ：
\begin{enumerate}
\item $\text{depth}(Q) \geq 1$，
\item $K$ は零であるか，または $\text{Supp}(K) = \{\mathfrak m\}$，
\item $\text{length}_R(K) < \infty$．
\end{enumerate}
ヒント：Noether 性を用いて (2) のような極大部分加群 $K$ が存在することを示し，
次に $Q = M/K$ が (1) を，$K$ が (3) を満たすことを示せ．
\end{exercise}

\begin{exercise}
\label{exercises-exercise-make-depth-2}
$(R, \mathfrak m)$ を Noether 局所環，$M$ を深さ $\geq 2$ の
有限 $R$-加群，$N \subset M$ を零でない部分加群とする．
\begin{enumerate}
\item $\text{depth}(N) \geq 1$ であることを示せ．
\item $\text{depth}(N) = 1$ であるための必要十分条件は，
商加群 $M/N$ が $\text{depth}(M/N) = 0$ を満たすことであることを示せ．
\item $N \subset N'$ が有限余長，すなわち
$\text{length}_R(N'/N) < \infty$ であり，かつ $N'$ の深さが
$\geq 2$ となる部分加群 $N' \subset M$ が存在することを示せ．
ヒント：演習 \ref{exercises-exercise-make-depth-1} を $M/N$ に適用し，
$N'$ を $K$ の逆像に選べ．
\end{enumerate}
\end{exercise}

\begin{exercise}
\label{exercises-exercise-Hartshorne-reduced}
$(R, \mathfrak m)$ を Noether 局所環とする．
$R$ は被約，すなわち $R$ は零でない冪零元をもたないと仮定する．
さらに，$R$ は相異なる二つの極小素イデアル $\mathfrak p$ と
$\mathfrak q$ をもつと仮定する．
\begin{enumerate}
\item
% P12-ERR-0427: frozen ungrouped quotient-ideal expressions; see ../control/ERRATA_CANDIDATES.jsonl.
$R$-加群の列
$$
0 \to R \to R/\mathfrak p \oplus R/\mathfrak q \to
R/\mathfrak p + \mathfrak q \to 0
$$
が完全であることを示せ（すべての箇所で確かめよ）．写像は
$x \mapsto (x \bmod \mathfrak p, x \bmod \mathfrak q)$ および
$(y \bmod \mathfrak p, z \bmod \mathfrak q) \mapsto
(y - z \bmod \mathfrak p + \mathfrak q)$
である．
\item $\text{depth}(R) \geq 2$ ならば，
$\dim(R/\mathfrak p + \mathfrak q) \geq 1$ であることを示せ．
\item $\text{depth}(R) \geq 2$ ならば，
$U = \Spec(R) \setminus \{\mathfrak m\}$ が連結位相空間であることを示せ．
\end{enumerate}
これは Hartshorne の連結性定理の非常に特殊な場合を証明している．
同定理は，$\text{depth} \geq 2$ の Noether 局所環の穴あきスペクトル
$U$ が連結であると述べる．
\end{exercise}

\begin{exercise}
\label{exercises-exercise-depth-2}
$(R, \mathfrak m)$ を Noether 局所環とし，
$x, y \in \mathfrak m$ を長さ $2$ の正則列とする．
任意の $n \geq 2$ に対して，
$$
x^{n - 1}y^{n - 1} = a x^n + b y^n
$$
を満たす $a, b \in R$ は存在しないことを示せ．
提案：まず $n = 2$ の場合を試し，議論の仕方を見いだせ．
注意：この結果には，単項式予想と呼ばれる大幅な一般化がある．
\end{exercise}







\section{Cohen--Macaulay 加群と環}
\label{exercises-section-CM}

\noindent
Cohen--Macaulay 加群は代数，節 \ref{algebra-section-CM} で研究し，
Cohen--Macaulay 環は代数，節 \ref{algebra-section-CM-ring} で研究する．

\begin{exercise}
\label{exercises-exercise-examples-CM}
次の各場合に，はい，または，いいえで答えよ．説明や証明は不要である．
\begin{enumerate}
\item $p$ を素数とする．
局所環 $\mathbf{Z}_{(p)}$ は Cohen--Macaulay 局所環か．
\item $p$ を素数とする．
局所環 $\mathbf{Z}_{(p)}$ は正則局所環か．
\item $k$ を体とする．
局所環 $k[x]_{(x)}$ は Cohen--Macaulay 局所環か．
\item $k$ を体とする．
局所環 $k[x]_{(x)}$ は正則局所環か．
\item $k$ を体とする．
局所環 $(k[x, y]/(y^2 - x^3))_{(x, y)} =
k[x, y]_{(x, y)}/(y^2 - x^3)$ は Cohen--Macaulay 局所環か．
\item $k$ を体とする．
局所環 $(k[x, y]/(y^2, xy))_{(x, y)} =
k[x, y]_{(x, y)}/(y^2, xy)$ は Cohen--Macaulay 局所環か．
\end{enumerate}
\end{exercise}

\section{特異点}
\label{exercises-section-singularities}

\begin{exercise}
\label{exercises-exercise-singularities}
$k$ を任意の体とする．$A = k[[x, y]]/(f)$ および
$B = k[[u, v]]/(g)$ とし，ここで $f = xy$，
$g = uv + \delta$，$\delta \in (u, v)^3$ とする．
$A$ と $B$ が同型な環であることを示せ．
\end{exercise}

\begin{remark}
\label{exercises-remark-singularities}
体 $k$ 上の曲線の特異点が通常二重点であるとは，
その点における曲線の完備局所環が $k'[[x, y]]/(f)$ の形であり，
次を満たすことをいう：
(a) $k'$ は $k$ の有限次分離拡大である；
(b) $f$ の初期項は次数二，すなわち，すべてが零ではない
$a, b, c\in k'$ によって $q = ax^2 + bxy + cy^2$ の形をしている；
(c) $q$ は $k'$ 上の非退化二次形式である
（標数 2 では，これは $b$ が零でないことを意味する）．
一般に，組 $(k', q)$ の各同型類に対して，そのような環の同型類が一つ存在する．
\end{remark}

\begin{exercise}
\label{exercises-exercise-periodic-resolution}
$R$ を環，$n \geq 1$ とする．$A$，$B$ を，係数が $R$ に属する
$n \times n$ 行列とし，$R$ のある非零因子 $f$ に対して
$AB = f 1_{n \times n}$ を満たすとする．$S = R/(f)$ とおく．
$$
\ldots \to
S^{\oplus n} \xrightarrow{B}
S^{\oplus n} \xrightarrow{A}
S^{\oplus n} \xrightarrow{B}
S^{\oplus n} \to \ldots
$$
が完全であることを示せ．
\end{exercise}

\section{構成可能集合}
\label{exercises-section-constructible}

\noindent
$k$ を代数閉体，たとえば複素数体 $\mathbf{C}$ とする．
$n \geq 0$ とする．多項式 $f \in k[x_1, \ldots, x_n]$ は，
値を代入することによって函数 $f : k^n \to k$ を定める．
部分集合 $Z \subset k^n$ が {\it 代数集合} であるとは，
多項式族の共通零点集合であることをいう．

\begin{exercise}
\label{exercises-exercise-finite-nr-equations}
代数集合が常に有限個の多項式の零点集合として書けることを証明せよ．
\end{exercise}

\noindent
上の記法のもとで，部分集合 $E \subset k^n$ が
{\it 構成可能} であるとは，多項式 $f$ による
$Z \cap \{f \not = 0\}$ の形の集合の有限合併であることをいう．

\begin{exercise}
\label{exercises-exercise-constructible-classical}
次を示せ：
\begin{enumerate}
\item 構成可能集合の補集合は構成可能集合である，
\item 構成可能集合の有限合併は構成可能集合である，
\item 構成可能集合の有限交叉は構成可能集合である，
\item 任意の構成可能集合 $E$ は有限な互いに素な合併
$E = \coprod E_i$ として書ける．ここで各 $E_i$ は
$Z \cap \{f \not = 0\}$ の形であり，$Z$ は代数集合，$f$ は多項式である．
\end{enumerate}
\end{exercise}

\begin{exercise}
\label{exercises-exercise-division-with-remainder}
$R$ を環とし，
$f = a_d x^d + a_{d - 1} x^{d - 1} + \ldots + a_0 \in R[x]$ とする．
（通常どおり，この記法は $a_0, \ldots, a_d \in R$ を意味する．）
$g \in R[x]$ とする．
$$
a_d^N g = q f + r
$$
かつ $\deg(r) < d$，すなわちある $c_i \in R$ に対して
$r = c_0 + c_1 x + \ldots + c_{d - 1}x^{d - 1}$
となる $N \geq 0$ および $r, q \in R[x]$ を見つけられることを証明せよ．
\end{exercise}

\section{Hilbert の零点定理}
\label{exercises-section-Hilbert-Nullstellensatz}

\begin{exercise}
\label{exercises-exercise-uncountable}
{\it 複素数を使った些末な議論！}
${\mathbf C}$ を複素数体とし，$V$ を ${\mathbf C}$ 上のベクトル空間とする．
線形作用素 $T : V \to V$ のスペクトルとは，作用素
$T - \lambda \text{id}_V$ が可逆でないような複素数
$\lambda \in {\mathbf C}$ の集合である．
\begin{enumerate}
\item $\mathbf{C}(X)$ の ${\mathbf C}$ 上の次元が非可算であることを示せ．
\item $V$ の次元が有限または可算ならば，$V$ 上の任意の線形作用素が
空でないスペクトルをもつことを示せ．
\item 有限生成 ${\mathbf C}$-代数 $R$ が体ならば，写像
${\mathbf C}\to R$ が同型であることを示せ．
\item ${\mathbf C}[x_1, \ldots, x_n]$ の任意の極大イデアル
${\mathfrak m}$ は，ある $\alpha_i \in {\mathbf C}$ によって
$(x_1-\alpha_1, \ldots, x_n-\alpha_n)$ の形に書けることを示せ．
\end{enumerate}
\end{exercise}

\begin{remark}
\label{exercises-remark-HNSS}
$k$ を体とする．このとき，任意の整数 $n\in {\mathbf N}$ と任意の極大イデアル
${\mathfrak m} \subset k[x_1, \ldots, x_n]$ に対して，商
$k[x_1, \ldots, x_n]/{\mathfrak m}$ は $k$ の有限次体拡大である．
これは講義の後半で示す．もちろん（確かめよ），このことから有限生成
$k$-代数の極大イデアルについて同様の主張が従う．
上の演習は $k = {\mathbf C}$ の場合にこれを証明している．
\end{remark}

\begin{exercise}
\label{exercises-exercise-Hilbert-Nullstellensatz}
$k$ を体とする．注 \ref{exercises-remark-HNSS} を用いよ．
\begin{enumerate}
\item $R$ を $k$-代数とする．$\dim_k R < \infty$ かつ
$R$ が整域であると仮定する．$R$ が体であることを示せ．
\item $R$ が有限生成 $k$-代数で，$f\in R$ が冪零でないと仮定する．
$f\not\in {\mathfrak m}$ を満たす極大イデアル
${\mathfrak m} \subset R$ が存在することを示せ．
\item $R$ が体上有限型でないとき，この主張が偽となることを例によって示せ．
\item 任意の根基イデアル $I \subset {\mathbf C}[x_1, \ldots, x_n]$ が，
それを含む極大イデアルの共通部分であることを示せ．
\end{enumerate}
\end{exercise}

\begin{remark}
\label{exercises-remark-Hilbert-Nullstellensatz}
これが Hilbert の零点定理である．すなわち，
$\Spec(k[x_1, \ldots, x_n])$ の閉部分集合
（先の演習により根基イデアルに対応する）は，
そこに含まれる閉点によって決定されると述べている．
\end{remark}

\begin{exercise}
\label{exercises-exercise-product-matrices-ring}
$A =
{\mathbf C}[x_{11}, x_{12}, x_{21}, x_{22}, y_{11}, y_{12}, y_{21}, y_{22}]$
とする．$I$ を $A$ のイデアルで，行列 $XY$ の成分によって生成されるものとする．
ここで
$$
X = \left(
\begin{matrix}
x_{11} & x_{12}\\
x_{21} & x_{22}
\end{matrix}
\right)
\quad\text{かつ}\quad
Y = \left(
\begin{matrix}
y_{11} & y_{12}\\
y_{21} & y_{22}
\end{matrix}
\right).
$$
$\Spec(A)$ の閉部分集合 $V(I)$ の既約成分を見つけよ．
（それらを記述し，各々を定める方程式を与えよ，という意味である．
書いた方程式が素イデアルを定めることは証明しなくてよい．）
ヒント：
\begin{enumerate}
\item Hilbert の零点定理を用いてよい．$V(I)$ の閉点全体を被覆する
既約局所閉部分集合を見つければ十分である．
\item 容易な成分が二つある．
\item 連続写像による既約集合の像は既約である．
\end{enumerate}
\end{exercise}

\section{次元}
\label{exercises-section-dimension}

\begin{exercise}
\label{exercises-exercise-dimension-bigger-one-finite-nr-primes}
素イデアルが有限個で，次元が $> 1$ である環 $A$ を構成せよ．
\end{exercise}

\begin{exercise}
\label{exercises-exercise-hypersurface-in-A2-dimension-one}
$f \in \mathbf{C}[x, y]$ を定数でない多項式とする．
$\mathbf{C}[x, y]/(f)$ の次元が $1$ であることを示せ．
\end{exercise}

\begin{exercise}
\label{exercises-exercise-dimension-polynomial-ring}
$(R, \mathfrak m)$ を Noether 局所環とし，$n \geq 1$ とする．
多項式環 $R[x_1, \ldots, x_n]$ において
$\mathfrak m' = (\mathfrak m, x_1, \ldots, x_n)$ とおく．
次を示せ：
$$
\dim(R[x_1, \ldots, x_n]_{\mathfrak m'}) = \dim(R) + n.
$$
\end{exercise}

\section{鎖状環}
\label{exercises-section-catenary}

\begin{definition}
\label{exercises-definition-catenary}
Noether 環 $A$ が {\it 鎖状} であるとは，
任意の三つの素イデアル
${\mathfrak p}_1 \subset {\mathfrak p}_2 \subset {\mathfrak p}_3$
に対して
$$
ht({\mathfrak p}_3 / {\mathfrak p}_1) = ht({\mathfrak p}_3/{\mathfrak p}_2) +
ht({\mathfrak p}_2/{\mathfrak p}_1).
$$
が成り立つことをいう．ここで $ht(\mathfrak p/\mathfrak q)$ は，
環 $A/\mathfrak q$ における $\mathfrak p/\mathfrak q$ の高さを意味する．
式で書けば，
$$
ht(\mathfrak p/\mathfrak q) =
\dim(A_\mathfrak p/\mathfrak qA_\mathfrak p) =
\dim((A/\mathfrak q)_\mathfrak p) =
\dim((A/\mathfrak q)_{\mathfrak p/\mathfrak q})
$$
位相空間 $X$ が {\it 鎖状} であるとは，
$T$ と $T'$ が閉かつ既約である $T \subset T' \subset X$ が与えられたとき，
既約閉部分集合の極大鎖
$$
T = T_0 \subset T_1 \subset \ldots \subset T_n = T'
$$
が存在し，そのような鎖がすべて同じ（有限の）長さをもつことをいう．
\end{definition}

\begin{exercise}
\label{exercises-exercise-catenary-the-same}
代数，定義 \ref{algebra-definition-catenary} で定義した鎖状性の概念が，
Noether 環については定義 \ref{exercises-definition-catenary} の概念と
一致することを示せ．
\end{exercise}

\begin{exercise}
\label{exercises-exercise-Noetherian-local-domain-dim-2-catenary}
次元 $2$ の Noether 局所整域が鎖状であることを示せ．
\end{exercise}

\begin{exercise}
\label{exercises-exercise-finite-type-over-field-catenary}
$k$ を体とする．有限型 $k$-代数が鎖状であることを示せ．
\end{exercise}

\begin{exercise}
\label{exercises-exercise-example-no-dim-function}
次元函数 $\delta : X \to \mathbf{Z}$ をもたない，有限，sober，かつ鎖状な
位相空間 $X$ の例を挙げよ．
ここで $\delta : X \to \mathbf{Z}$ が次元函数であるとは，
$x, y \in X$ に対して次が成り立つことをいう：
\begin{enumerate}
\item $x \leadsto y$ かつ $x \not = y$ ならば $\delta(x) > \delta(y)$，
\item $x \leadsto y$ かつ $\delta(x) \geq \delta(y) + 2$ ならば，
$x \leadsto z \leadsto y$ および $\delta(x) > \delta(z) > \delta(y)$ を満たす
$z \in X$ が存在する．
\end{enumerate}
空間を明確に記述し，次元函数が存在しえない理由を簡潔に説明せよ．
\end{exercise}

\section{分数体}
\label{exercises-section-fraction-fields}

\begin{exercise}
\label{exercises-exercise-find-fraction-field}
整域
$$
{\mathbf Q}[r, s, t]/(s^2-(r-1)(r-2)(r-3), t^2-(r + 1)(r + 2)(r + 3)).
$$
を考える．これと同型な分数体をもつ
${\mathbf Q}[x, y]/(f)$ の形の整域を見つけよ．
\end{exercise}

\section{超越次数}
\label{exercises-section-transcendence}

\begin{exercise}
\label{exercises-exercise-algebraic-extension}
$K'/K/k$ を体拡大とし，$K'$ は $K$ 上代数的とする．
$\text{trdeg}_k(K) = \text{trdeg}_k(K')$ を証明せよ．
（ヒント：$x_1, \ldots, x_d \in K$ が $k$ 上代数的独立であり，
$d < \text{trdeg}_k(K')$ ならば，
$k(x_1, \ldots, x_d) \subset K$ は代数拡大ではありえないことを示せ．）
\end{exercise}

\begin{exercise}
\label{exercises-exercise-growth-powers-subvector-space}
$k$ を体とし，$K/k$ を超越次数 $d$ の有限生成拡大とする．
$V, W \subset K$ が有限次元の $k$-部分ベクトル空間であるとき，
$$
VW = \{f \in K \mid f = \sum\nolimits_{i = 1, \ldots, n} v_i w_i
\text{ となるある }n\text{ と }v_i \in V, w_i \in W\}
$$
と書く．これは有限次元の $k$-部分ベクトル空間である．
$V^2 = VV$，$V^3 = V V^2$，等々とおく．
\begin{enumerate}
\item すべての $n \geq 1$ に対して $\dim V^n \geq \epsilon n^d$ となる
$V \subset K$ および $\epsilon > 0$ を見つけられることを示せ．
\item 逆に，任意の有限次元 $V \subset K$ に対して，すべての
$n \geq 1$ に対し $\dim V^n \leq C n^d$ となる $C > 0$ が存在することを示せ．
（進め方の一例：
まず $k[x_1, \ldots, x_d]$ の部分ベクトル空間について示せ．
次に $k(x_1, \ldots, x_d)$ の部分ベクトル空間について示せ．
最後に，$K/k(x_1, \ldots, x_d)$ が有限拡大である場合，
$K$ の $k(x_1, \ldots, x_d)$ 上の基底を選び，
この基底による展開を用いて議論せよ．）
\item $K$ の有限次元部分ベクトル空間の冪の増大度によって，
超越次数を定義し直せることを導け．
\end{enumerate}
これは $k$ 上の（非可換）代数の Gelfand--Kirillov 次元に関係している．
\end{exercise}

\section{ファイバーの次元}
\label{exercises-section-dimension-fibres}

\noindent
次元公式に関する問題をいくつか挙げる．
代数，節 \ref{algebra-section-dimension-formula} を参照せよ．

\begin{exercise}
\label{exercises-exercise-nr-components-fibre}
$k$ を好みの代数閉体とする．以下，$k[x]$ と $k[x, y]$ は多項式環を表す．
\begin{enumerate}
\item 任意の整数 $n \geq 0$ に対して，$A/xA$ のスペクトルが
ちょうど $n$ 個の既約成分をもつような，整域の有限型拡大
$k[x] \subset A$ を見つけよ．
\item すべての $\alpha \in k$ に対して $A/(x - \alpha)A$ のスペクトルが
空でなく可約となるような，整域の有限型拡大 $k[x] \subset A$ の例を作れ．
\item すべての $(\alpha, \beta) \in k^2 \setminus \{(0, 0)\}$ に対して
$A/(x - \alpha, y - \beta)A$ のスペクトルが
既約\footnote{既約ならば空でないことを思い出せ．}であり，
$A/(x, y)A$ のスペクトルが空でなく可約となるような，
整域の有限型拡大 $k[x, y] \subset A$ の例を作れ．
\end{enumerate}
\end{exercise}

\begin{exercise}
\label{exercises-exercise-codim-1}
$k$ を好みの代数閉体とする．$n \geq 1$ とする．
$k[x_1, \ldots, x_n]$ を多項式環とし，
$\mathfrak m = (x_1, \ldots, x_n)$ とおく．
$k[x_1, \ldots, x_n] \subset A$ を整域の有限型拡大とする．
$d = \dim(A)$ とおく．
\begin{enumerate}
\item $A/\mathfrak mA \not = 0$ ならば
$d - 1 \geq \dim(A/\mathfrak m A) \geq d - n$ であることを示せ．
\item 各値が実際に現れうることを例によって示せ．
\item $\Spec(A/\mathfrak m A)$ が異なる次元の既約成分をもちうることを
例によって示せ．
\end{enumerate}
\end{exercise}

\section{有限局所自由加群}
\label{exercises-section-finite-locally-free}

\begin{definition}
\label{exercises-definition-finite-locally-free}
$A$ を環とする．{\it 有限局所自由} $A$-加群 $M$ とは，
任意の ${\mathfrak p} \in \Spec(A)$ に対して，
$M_f$ が有限自由 $A_f$-加群となる
$f\in A$，$f \not \in {\mathfrak p}$ が存在するような加群であることを思い出そう．
$M$ が階数 $1$ の有限局所自由加群であるとき，$M$ を
{\it 可逆加群} という．すなわち，任意の ${\mathfrak p} \in \Spec(A)$ に対して，
$A_f$-加群として $M_f \cong A_f$ となる
$f\in A$，$f \not \in \mathfrak p$ が存在する．
\end{definition}

\begin{exercise}
\label{exercises-exercise-tensor-finite-locally-free}
有限局所自由加群のテンソル積が有限局所自由であることを証明せよ．
二つの可逆加群のテンソル積が可逆であることを証明せよ．
\end{exercise}

\begin{definition}
\label{exercises-definition-class-group}
$A$ を環とする．{\it $A$ の類群}，ときには {\it $A$ の Picard 群} とも呼ばれるものは，
可逆 $A$-加群の同型類全体の集合 $\Pic(A)$ に，
テンソル積で定まる群演算を入れたものである
（演習 \ref{exercises-exercise-tensor-finite-locally-free} を参照）．
\end{definition}

\noindent
$A$ の類群が自明であるための必要十分条件は，
すべての可逆加群が階数 1 の自由加群と同型であることであることに注意せよ．

\begin{exercise}
\label{exercises-exercise-class-group-trivial}
次の各環の類群が自明であることを示せ：
\begin{enumerate}
\item 体 $k$ 上の多項式環 $A = k[x]$，
\item 整数環 $A = \mathbf{Z}$，
\item 体 $k$ 上の多項式環 $A = k[x, y]$，
\item 体 $k$ 上の商 $k[x, y]/(xy)$．
\end{enumerate}
\end{exercise}

\begin{exercise}
\label{exercises-exercise-class-group-not-trivial}
$k$ を標数が $2$ でない体とし，$t_1, \ldots, t_n \in k$ は相異なり，
$n \geq 3$ は奇数，$f(x) = (x - t_1) \ldots (x - t_n)$ とする．
環 $A = k[x, y]/(y^2 - f(x))$ の類群が自明でないことを示せ．
（ヒント：イデアル $(y, x - t_1)$ が $\Pic(A)$ の非自明な元を
定めることを示せ．）
\end{exercise}

\begin{exercise}
\label{exercises-exercise-trace-det}
$A$ を環とする．
\begin{enumerate}
\item $M$ を有限局所自由 $A$-加群とし，
$\varphi : M \to M$ を自己準同型とする．
$\varphi$ の {\it トレース} と {\it 行列式} を定義・構成し，
この構成が「三つ組 $(A, M, \varphi)$ に関して関手的」であることを証明せよ．
\item
% P12-ERR-0428: frozen determinant assertion without equal-rank hypothesis; see ../control/ERRATA_CANDIDATES.jsonl.
$M, N$ が有限局所自由 $A$-加群であり，
$\varphi : M \to N$ および $\psi : N \to M$ ならば，
$\text{Trace}(\varphi \circ \psi) = \text{Trace}(\psi \circ \varphi)$ および
$\det(\varphi \circ \psi) = \det(\psi \circ \varphi)$ であることを示せ．
\item $M$ が有限局所自由の場合，$\text{Trace}$ が
$A$-線形写像 $\text{End}_A(M) \to A$ を定め，
$\det$ が乗法的写像 $\text{End}_A(M) \to A$ を定めることを示せ．
\end{enumerate}
\end{exercise}

\begin{exercise}
\label{exercises-exercise-trace-det-rings}
今度は $B$ を，$A$-加群として有限局所自由な $A$-代数とする．
言い換えると，$B$ は有限局所自由 $A$-代数である．
\begin{enumerate}
\item 演習 \ref{exercises-exercise-trace-det} の $\text{Trace}$ と $\det$ を用いて，
$\text{Trace}_{B/A}$ と $\text{Norm}_{B/A}$ を定義せよ．
\item $b\in B$ とし，$\pi : \Spec(B) \to \Spec(A)$ を誘導される射とする．
$\pi(V(b)) = V(\text{Norm}_{B/A}(b))$ を示せ．
（$V(f) = \{ {\mathfrak p} \mid f \in {\mathfrak p}\}$ を思い出せ．）
\item （基底変換．）$i : A \to A'$ を環準同型とし，
$B' = B \otimes_A A'$ とおく．
$i(\text{Norm}_{B/A}(b))$ が $\text{Norm}_{B'/A'}(b \otimes 1)$ に等しい理由を述べよ．
\item $B = A \times A \times A \times \ldots \times A$ および
$b = (a_1, \ldots, a_n)$ のとき，$\text{Norm}_{B/A}(b)$ を計算せよ．
\item 有限平坦写像 ${\mathbf Q}[x] \to {\mathbf Q}[y]$，
$x \to y^n$ に関する $y-y^3$ のノルムを計算せよ．
（ヒント：「基底変換」
$A = {\mathbf Q}[x] \subset A' = {\mathbf Q}(\zeta_n)(x^{1/n})$ を用いよ．）
\end{enumerate}
\end{exercise}

\section{貼り合わせ}
\label{exercises-section-glueing}

\begin{exercise}
\label{exercises-exercise-cover}
$A$ を環とし，$M$ を $A$-加群とする．
$f_i$，$i \in I$ を，次を満たす $A$ の元の族とする：
$$
\Spec(A) = \bigcup D(f_i).
$$
\begin{enumerate}
\item $M_{f_i}$ が有限 $A_{f_i}$-加群ならば，
$M$ が有限 $A$-加群であることを示せ．
\item $M_{f_i}$ が平坦 $A_{f_i}$-加群ならば，
$M$ が平坦 $A$-加群であることを示せ．
（正しい仕方で考えれば，これはいささか些末なことである．）
\end{enumerate}
\end{exercise}

\begin{remark}
\label{exercises-remark-cover}
代数幾何学の言葉では，これは「有限生成である」あるいは「平坦である」という性質が，
（適切な意味で）Zariski 位相に関して局所的であることを意味する．
「有限表示である」という性質についても同じことを示せる．
\end{remark}

\begin{exercise}
\label{exercises-exercise-cover-ring-map}
$A \to B$ を環準同型とする．
$f_i \in A$，$i \in I$ および $g_j \in B$，$j \in J$ を，次を満たす元の族とする：
$$
\Spec(A) = \bigcup D(f_i)
\quad\text{かつ}\quad
\Spec(B) = \bigcup D(g_j).
$$
すべての $i, j$ に対して $A_{f_i} \to B_{f_ig_j}$ が有限型ならば，
$A \to B$ が有限型であることを示せ．
\end{exercise}

\section{上昇定理と下降定理}
\label{exercises-section-going-up}

\begin{definition}
\label{exercises-definition-GU-GD}
$\phi : A \to B$ を環準同型とする．
次の条件を満たすとき，$\phi$ に対して {\it 上昇定理} が成り立つという：
\begin{itemize}
\item[(GU)] ${\mathfrak p} \subset {\mathfrak p}'$ を満たす任意の
${\mathfrak p}, {\mathfrak p}' \in \Spec(A)$ と，
${\mathfrak p}$ の上にある任意の $P \in \Spec(B)$ に対して，
$P \subset P'$ を満たし，${\mathfrak p}'$ の上にある
$P'\in \Spec(B)$ が存在する．
\end{itemize}
同様に，次の条件を満たすとき，$\phi$ に対して {\it 下降定理} が成り立つという：
\begin{itemize}
\item[(GD)] ${\mathfrak p} \subset {\mathfrak p}'$ を満たす任意の
${\mathfrak p}, {\mathfrak p}' \in \Spec(A)$ と，
${\mathfrak p}'$ の上にある任意の $P' \in \Spec(B)$ に対して，
$P \subset P'$ を満たし，${\mathfrak p}$ の上にある
$P\in \Spec(B)$ が存在する．
\end{itemize}
\end{definition}

\begin{exercise}
\label{exercises-exercise-GU-GD}
次の各場合に (GU)，(GD) が成り立つかどうかを判定し，その理由を説明せよ．
（見つけられるどの命題・定理・補題を用いてもよいが，各場合に仮定を確かめよ．）
\begin{enumerate}
\item $k$ は体，$A = k$，$B = k[x]$．
\item $k$ は体，$A = k[x]$，$B = k[x, y]$．
\item $A = {\mathbf Z}$，$B = {\mathbf Z}[1/11]$．
\item $k$ は代数閉体，$A = k[x, y]$，
$B = k[x, y, z]/(x^2-y, z^2-x)$．
\item $A = {\mathbf Z}$，$B = {\mathbf Z}[i, 1/(2 + i)]$．
\item $A = {\mathbf Z}$，$B = {\mathbf Z}[i, 1/(14 + 7i)]$．
\item $k$ は代数閉体，$A = k[x]$，
$B = k[x, y, 1/(xy-1)]/(y^2-y)$．
\end{enumerate}
\end{exercise}

\begin{exercise}
\label{exercises-exercise-image}
$A$ を環とし，$B = A[x]$ を $A$ 上の一変数多項式代数とする．
$f = a_0 + a_1 x + \ldots + a_r x^r \in B = A[x]$ とする．
$\Spec(A)$ における $D(f)$ の像が
$D(a_0) \cup \ldots \cup D(a_r)$ に等しいことを丁寧に証明せよ．
\end{exercise}

\begin{exercise}
\label{exercises-exercise-images}
$k$ を代数閉体とする．次の各写像の $\Spec(k[x, y])$ における像を計算せよ：
\begin{enumerate}
\item $\Spec(k[x, yx^{-1}]) \to \Spec(k[x, y])$，ここで
$k[x, y] \subset k[x, yx^{-1}] \subset k[x, y, x^{-1}]$．
（ヒント：混乱を避けるため，元 $yx^{-1}$ に別の名前を与えよ．）
\item $\Spec(k[x, y, a, b]/(ax-by-1))\to \Spec(k[x, y])$．
\item $x \mapsto t^2$ および $y \mapsto t^3$ により誘導される
$\Spec(k[t, 1/(t-1)]) \to \Spec(k[x, y])$．
\item $k = {\mathbf C}$（複素数），
$\Spec(k[s, t]/(s^3 + t^3-1)) \to \Spec(k[x, y])$，ここで
$x\mapsto s^2$，$y \mapsto t^2$．
\end{enumerate}
\end{exercise}

\begin{remark}
\label{exercises-remark-elimination-theory}
上のように像を求めるには，通常，消去法を用いる．
\end{remark}

\section{Fitting イデアル}
\label{exercises-section-fitting-ideals}

\begin{exercise}
\label{exercises-exercise-fitting}
$R$ を環，$M$ を有限 $R$-加群とする．$M$ の表示
$$
\bigoplus\nolimits_{j \in J} R \longrightarrow R^{\oplus n}
\longrightarrow M \longrightarrow 0.
$$
を選ぶ．写像 $R^{\oplus n} \to M$ は $M$ の元の列
$x_1, \ldots, x_n$ によって与えられることに注意せよ．
元 $x_i$ は $M$ の {\it 生成元} である．
写像 $\bigoplus_{j \in J} R \to R^{\oplus n}$ は，係数が $R$ に属する
$n \times J$ 行列 $A$ によって与えられる．
言い換えると，$A = (a_{ij})_{i = 1, \ldots, n, j \in J}$ である．
$A$ の列 $(a_{1j}, \ldots, a_{nj})$，$j \in J$ を {\it 関係式} という．
$\sum r_i x_i = 0$ を満たす任意のベクトル $(r_i) \in R^{\oplus n}$ は，
$A$ の列の線形結合である．もちろん，任意の有限 $R$-加群は多くの異なる表示をもつ．
\begin{enumerate}
\item $A$ の $(n - k) \times (n - k)$ 小行列式で生成されるイデアルが，
表示の選択に依存しないことを示せ．このイデアルを
{\it $M$ の第 $k$ Fitting イデアル} といい，$Fit_k(M)$ と書く．
\item
$Fit_0(M) \subset Fit_1(M) \subset Fit_2(M) \subset \ldots$
を示せ．（ヒント：行列式は一つの列に沿って展開して計算できることを用いよ．）
\item 次が同値であることを示せ：
\begin{enumerate}
\item $Fit_{r - 1}(M) = (0)$ かつ $Fit_r(M) = R$，
\item $M$ は階数 $r$ の局所自由加群である．
\end{enumerate}
\end{enumerate}
\end{exercise}

\section{Hilbert 函数}
\label{exercises-section-hilbert}

\begin{definition}
\label{exercises-definition-numerical-polynomial}
{\it 整数値多項式} とは，任意の整数 $n$ に対して
$f(n) \in {\mathbf Z}$ となる多項式 $f(x) \in {\mathbf Q}[x]$ のことである．
\end{definition}

\begin{definition}
\label{exercises-definition-graded-module}
環 $A$ 上の {\it 次数付き加群} $M$ とは，$A$-部分加群への直和分解
$
\bigoplus\nolimits_{n \in {\mathbf Z}} M_n
$
を備えた $A$-加群 $M$ のことである．すべての $M_n$ が有限
$A$-加群であるとき，$M$ は {\it 局所有限} であるという．
$A$ を Noether 環とし，$\varphi$ を有限 $A$-加群上の
{\it Euler--Poincar\'e 函数} とする．これは，任意の有限生成 $A$-加群 $M$ に
整数 $\varphi(M) \in {\mathbf Z}$ が与えられ，任意の短完全列
$$
0
\longrightarrow
M'
\longrightarrow
M
\longrightarrow
M''
\longrightarrow
0
$$
に対して $\varphi(M) = \varphi(M') + \varphi(M'')$ となることを意味する．
局所有限次数付き加群 $M$ の（$\varphi$ に関する）
{\it Hilbert 函数} とは，函数 $\chi_\varphi(M, n) = \varphi(M_n)$ のことである．
十分大きいすべての整数 $n$ に対して
$\chi_\varphi(M, n) = P_\varphi(n)$ となる整数値多項式 $P_\varphi$ が存在するとき，
$M$ は {\it Hilbert 多項式} をもつという．
\end{definition}

\begin{definition}
\label{exercises-definition-graded-algebra}
{\it 次数付き $A$-代数} とは，次数付き $A$-加群
$B = \bigoplus_{n \geq 0} B_n$ と，$A$-双線形写像
$$
B \times B \longrightarrow B, \ (b, b') \longmapsto bb'
$$
であって，$B$ を $A$-代数にし，
$B_n \cdot B_m \subset B_{n + m}$ を満たすもののことである．
最後に，{\it 次数付き $A$-代数 $B$ 上の次数付き加群 $M$} とは，
次数付き $A$-加群 $M$ に，（両立する）$B$-加群構造で
$B_n \cdot M_d \subset M_{n + d}$ を満たすものを与えたものである．
これで，{\it 次数付き加群・環の準同型}，{\it 次数付き部分加群}，
{\it 次数付きイデアル}，{\it 次数付き加群の完全列}，等々を定義できる．
\end{definition}

\begin{exercise}
\label{exercises-exercise-Euler-Poincare-field}
$A = k$ を体とする．この場合，有限 $A$-加群上の
Euler--Poincar\'e 函数としてどのようなものがあるか，すべて求めよ．
\end{exercise}

\begin{exercise}
\label{exercises-exercise-Euler-Poincare-Z}
$A ={\mathbf Z}$ とする．この場合，有限 $A$-加群上の
Euler--Poincar\'e 函数としてどのようなものがあるか，すべて求めよ．
\end{exercise}

\begin{exercise}
\label{exercises-exercise-Euler-Poincare-node}
$A = k[x, y]/(xy)$ とし，$k$ を代数閉体とする．この場合，有限
$A$-加群上の Euler--Poincar\'e 函数としてどのようなものがあるか，すべて求めよ．
\end{exercise}

\begin{exercise}
\label{exercises-exercise-kernel-locally-finite}
$A$ を Noether 環とする．局所有限次数付き $A$-加群の写像の核が
局所有限であることを示せ．
\end{exercise}

\begin{exercise}
\label{exercises-exercise-no-hilbert}
$k$ を体とし，$A = k$ および $B = k[x, y]$ とおく．
次数は $\deg(x) = 2$ と $\deg(y) = 3$ によって定める．
$\varphi(M) = \dim_k(M)$ とする．次数付き $k$-加群としての
$B$ の Hilbert 函数を計算せよ．この場合，Hilbert 多項式は存在するか．
\end{exercise}

\begin{exercise}
\label{exercises-exercise-no-hilbert-or-is-there}
$k$ を体とし，$A = k$ および $B = k[x, y]/(x^2, xy)$ とおく．
次数は $\deg(x) = 2$ と $\deg(y) = 3$ によって定める．
$\varphi(M) = \dim_k(M)$ とする．次数付き $k$-加群としての
$B$ の Hilbert 函数を計算せよ．この場合，Hilbert 多項式は存在するか．
\end{exercise}

\begin{exercise}
\label{exercises-exercise-hilbert-to-compute}
$k$ を体とし，$A = k$ とおく．$\varphi(M) = \dim_k(M)$ とする．
$d\in {\mathbf N}$ を固定する．次数付き $A$-代数
$B = k[x, y, z]/(x^d + y^d + z^d)$ を考える．ここで $x, y, z$ は
それぞれ次数 $1$ をもつ．$B$ の Hilbert 函数を計算せよ．
この場合，Hilbert 多項式は存在するか．
\end{exercise}

\section{環の Proj}
\label{exercises-section-proj-ring}

\begin{definition}
\label{exercises-definition-homogeneous-ideal}
$R$ を次数付き環とする．{\it 斉次}イデアルとは，単に $R$ の次数付き部分加群でもある
イデアル $I \subset R$ のことである．同値なことだが，斉次元で生成される
イデアルのことである．さらに同値なことだが，$f \in I$ であり，
$$
f = f_0 + f_1 + \ldots + f_n
$$
が $R$ における $f$ の斉次成分への分解ならば，各 $i$ に対して
$f_i \in I$ となることである．
\end{definition}

\begin{definition}
\label{exercises-definition-Proj-R}
次数付き環 $R$ の {\it 斉次スペクトル $\text{Proj}(R)$} を，
$R_{+} \not \subset {\mathfrak p}$ を満たす $R$ の斉次素イデアル
${\mathfrak p}$ 全体の集合として定義する．
$\text{Proj}(R)$ は $\Spec(R)$ の部分集合であり，
したがって自然な誘導位相をもつことに注意せよ．
\end{definition}

\begin{definition}
\label{exercises-definition-Dplus-Vplus}
$R = \oplus_{d \geq 0} R_d$ を次数付き環，$f\in R_d$ とし，$d \geq 1$ と仮定する．
{\it $R_{(f)}$} を，$r$ が斉次で $\deg(r) = nd$ を満たす
$r/f^n$ の形の元からなる $R_f$ の部分環として定義する．さらに，
$$
D_{+}(f) = \{ {\mathfrak p} \in \text{Proj}(R) | f \not\in {\mathfrak p} \}.
$$
と定める．最後に，斉次イデアル $I \subset R$ に対して
$V_{+}(I) = V(I) \cap \text{Proj}(R)$ と定める．
\end{definition}

\begin{exercise}
\label{exercises-exercise-topology-proj}
$\text{Proj}(R)$ 上の位相について．上の定義と記法のもとで，次の主張を証明せよ．
\begin{enumerate}
\item $D_{+}(f)$ が $\text{Proj}(R)$ において開であることを示せ．
\item $D_{+}(ff') = D_{+}(f) \cap D_{+}(f')$ を示せ．
\item $g = g_0 + \ldots + g_m$ を $R$ の元とし，$g_i \in R_i$ とする．
$D(g) \cap \text{Proj}(R)$ を，$D_{+}(g_i)$，$i \geq 1$ および
$D(g_0) \cap \text{Proj}(R)$ を用いて表せ．証明は不要である．
\item $g\in R_0$ を次数 $0$ の斉次元とする．
$D(g) \cap \text{Proj}(R)$ を，正次数斉次元の適切な族
$f_\alpha \in R$ に対する $D_{+}(f_\alpha)$ を用いて表せ．
\item 開集合族 $\{D_{+}(f)\}$ が $\text{Proj}(R)$ の位相の基底を
なすことを示せ．
\item
\label{exercises-item-bijection}
標準的な全単射 $D_{+}(f) \to \Spec(R_{(f)})$ が存在することを示せ．
% P12-ERR-0429: frozen imperative typo; see ../control/ERRATA_CANDIDATES.jsonl.
（ヒント：$\Spec$ の場合の証明をまね，ただし途中のどこかで
イデアルの根基を取り入れよ．）
\item (\ref{exercises-item-bijection}) の写像が同相写像であることを示せ．
\item $\text{Proj}(R)$ が準コンパクトでないような $R$ の例を挙げよ．
証明は不要である．
\item 任意の閉部分集合 $T \subset \text{Proj}(R)$ が，ある斉次イデアル
$I \subset R$ に対して $V_{+}(I)$ の形になることを示せ．
\end{enumerate}
\end{exercise}

\begin{remark}
\label{exercises-remark-continuous-proj-spec}
連続写像 $ \text{Proj}(R) \longrightarrow \Spec(R_0) $ が存在する．
\end{remark}

\begin{exercise}
\label{exercises-exercise-iso-polynomial-ring-one-variable}
$R = A[X]$ かつ $\deg(X) = 1$ ならば，自然な写像
$\text{Proj}(R) \to \Spec(A)$ が全単射であり，実際には同相写像であることを示せ．
\end{exercise}

\begin{exercise}
\label{exercises-exercise-blowing-up-I}
吹き上げ：第 I 部．
この演習では $R = Bl_I(A) = A \oplus I \oplus I^2 \oplus \ldots$ とする．
自然な写像 $b : \text{Proj}(R) \to \Spec(A)$ を考える．
$U = \Spec(A) - V(I)$ とおく．
$$
b : b^{-1}(U) \longrightarrow U
$$
が同相写像であることを示せ．
したがって $U$ を $\text{Proj}(R)$ の開部分集合とみなすことができる．
$Z \subset \Spec(A)$ を生成点 $\xi \in Z$ をもつ既約閉部分スキームとする．
$\xi \not\in V(I)$，言い換えると $Z \not\subset V(I)$，
さらに言い換えると $\xi \in U$，さらに言い換えると
$Z\cap U \not = \emptyset$ と仮定する．
$Z$ の {\it 狭義変換} $Z'$ を，$\xi$ の上にある一意な点 $\xi'$ の閉包として定める．
別の言い方をすれば，$Z'$ は局所閉部分集合
$Z\cap U \subset U \subset \text{Proj}(R)$ の $\text{Proj}(R)$ における閉包である．
\end{exercise}

\begin{exercise}
\label{exercises-exercise-blowing-up-II}
吹き上げ：第 II 部．
$A = k[x, y]$ とし，$k$ を体，$I = (x, y)$ とする．
$R$ を $A$ と $I$ に関する吹き上げ代数とする．
\begin{enumerate}
\item $Z_1 = V(\{x\})$ と $Z_2 = V(\{y\})$ の狭義変換が互いに素であることを示せ．
\item $Z_1 = V(\{x\})$ と $Z_2 = V(\{x-y^2\})$ の狭義変換が
互いに素でないことを示せ．
\item $V(J) = V(I)$ を満たし，$J$ に沿う吹き上げにおける
$Z_1 = V(\{x\})$ と $Z_2 = V(\{x-y^2\})$ の狭義変換が
互いに素となるようなイデアル $J \subset A$ を見つけよ．
\end{enumerate}
\end{exercise}

\begin{exercise}
\label{exercises-exercise-proj-when-empty}
$R$ を次数付き環とする．
\begin{enumerate}
\item
% P12-ERR-0430: frozen double-greater-than tokens; see ../control/ERRATA_CANDIDATES.jsonl.
すべての $n >> 0$ に対して $R_n = (0)$ ならば，
$\text{Proj}(R)$ が空であることを示せ．
\item $R$ が整域で $R_{+}$ が零でなければ，
$\text{Proj}(R)$ が既約位相空間であることを示せ．
（空の位相空間は既約でないことを思い出せ．）
\end{enumerate}
\end{exercise}

\begin{exercise}
\label{exercises-exercise-blowing-up-III}
吹き上げ：第 III 部．
狭義変換の定義における $A$，$I$ および $U$，$Z$ を考える．
ある素イデアル ${\mathfrak p}$ に対して $Z = V({\mathfrak p})$ とする．
$\bar A = A/{\mathfrak p}$ とし，$\bar I$ を $\bar A$ における $I$ の像とする．
\begin{enumerate}
\item
% P12-ERR-0431: frozen spaced definitional-equality tokens; see ../control/ERRATA_CANDIDATES.jsonl.
全射な環準同型
$R: = Bl_I(A) \to \bar R: = Bl_{\bar I}(\bar A)$ が存在することを示せ．
\item
% P12-ERR-0432: frozen manual locator 5(b); see ../control/ERRATA_CANDIDATES.jsonl.
上の環準同型が $\text{Proj}(\bar R)$ から $Z$ の狭義変換 $Z'$ への
全単射を誘導することを示せ．（これはそれほど容易ではない．ヒント：上の 5(b) を用いよ．）
\item 狭義変換は $Z' = V_{+}(P)$ となることを導け．ここで
$P \subset R$ は $P_d = I^d \cap {\mathfrak p}$ で定まる斉次イデアルである．
\item $Z_1 = V({\mathfrak p})$ と $Z_2 = V({\mathfrak q})$ を，
素イデアルで定まる既約閉部分集合で，$Z_1 \not \subset Z_2$ かつ
$Z_2 \not \subset Z_1$ を満たすものとする．
イデアル $I = {\mathfrak p} + {\mathfrak q}$ の吹き上げが，
$Z_1$ と $Z_2$ の狭義変換を分離すること，すなわち
$Z_1' \cap Z_2' = \emptyset$ を示せ．
（ヒント：(c) の斉次イデアル $P$ と $Q$ を考え，$V(P + Q)$ を考えよ．）
\end{enumerate}
\end{exercise}

% Frozen source: exercises.tex lines 2053--2236.
\section{次元 1 の Cohen--Macaulay 環}
\label{exercises-section-CM-dim-1}

\begin{definition}
\label{exercises-definition-CM}
Noether 局所環 $A$ が次元 $d$ の {\it Cohen--Macaulay 環}であるとは，
その次元が $d$ であり，$A$ のパラメータ系 $x_1, \ldots, x_d$ であって，
$i = 1, \ldots, d$ に対し $x_i$ が $A/(x_1, \ldots, x_{i-1})$
の非零因子となるものが存在することをいう．
\end{definition}

\begin{exercise}
\label{exercises-exercise-CM-dim-1-I}
次元 1 の Cohen--Macaulay 環．第 I 部：理論．
\begin{enumerate}
% P12-ERR-0433: source omits the head noun "ring".
\item $(A, {\mathfrak m})$ を $\dim A = 1$ である Noether 局所環とする．
$x\in {\mathfrak m}$ が零因子でなければ，次が成り立つことを示せ．
\begin{enumerate}
\item $\dim A/xA = 0$，言い換えると $A/xA$ は Artin 環であり，
さらに言い換えると $\{x\}$ は $A$ のパラメータ系である．
% P12-ERR-0434: duplicated finite verb in the source.
\item $A$ は埋め込まれた素イデアルを持たない．
\end{enumerate}
\item 逆に，$(A, {\mathfrak m})$ を次元 $1$ の Noether 局所環とする．
$A$ が埋め込まれた素イデアルを持たなければ，
${\mathfrak m}$ に非零因子が存在することを示せ．
\end{enumerate}
\end{exercise}

\begin{exercise}
\label{exercises-exercise-CM-dim-1-II}
次元 1 の Cohen--Macaulay 環．第 II 部：例．
\begin{enumerate}
\item $A$ を $k[x, y]/(x^2, xy)$ の $(x, y)$ における局所環とする．
\begin{enumerate}
\item $A$ の次元が 1 であることを示せ．
\item ${\mathfrak m}\subset A$ のすべての元が零因子であることを証明せよ．
\item $\dim A/zA = 0$ となる $z\in {\mathfrak m}$ を求めよ
（証明は不要）．
\end{enumerate}
\item $A$ を $k[x, y]/(x^2)$ の $(x, y)$ における局所環とする．
${\mathfrak m}$ の非零因子を求めよ（証明は不要）．
\end{enumerate}
\end{exercise}

\begin{exercise}
\label{exercises-exercise-embedding-dim-1}
埋め込み次元 $1$ の局所環．
$(A, {\mathfrak m}, k)$ を埋め込み次元 $1$ の Noether 局所環，
すなわち
$$
\dim_k {\mathfrak m}/{\mathfrak m}^2 = 1.
$$
を満たすものとする．
函数 $f(n) = \dim_k {\mathfrak m}^n/{\mathfrak m}^{n + 1}$ は
値 $1$ の定数函数であるか，またはその値が
$$
1, 1, \ldots, 1, 0, 0, 0, 0, 0, \ldots
$$
となることを示せ．
\end{exercise}

\begin{exercise}
\label{exercises-exercise-regular-local-dim-1}
次元 $1$ の正則局所環．
$(A, {\mathfrak m}, k)$ を次元 $1$ の正則 Noether 局所環とする．
これは $A$ の次元が $1$ であり，埋め込み次元も $1$，
すなわち
$$
\dim_k {\mathfrak m}/{\mathfrak m}^2 = 1.
$$
であることを意味することを思い出そう．
$x\in{\mathfrak m}$ を，${\mathfrak m}/{\mathfrak
m}^2$ における類が零でない任意の元とする．
\begin{enumerate}
% P12-ERR-0543: preserve the universal quantifier, including zero.
\item ${\mathfrak m}$ の各元 $y$ に対して整数 $n$ が存在し，
単元 $u\in A^\ast$ を用いて $y$ を $y = ux^n$ と書けることを示せ．
\item $x$ が $A$ の非零因子であることを示せ．
\item $A$ が整域であると結論せよ．
\end{enumerate}
\end{exercise}

\begin{exercise}
\label{exercises-exercise-nonzerodivisor-graded}
$(A, {\mathfrak m}, k)$ を随伴次数環 $Gr_{\mathfrak m}(A)$
を持つ Noether 局所環とする．
\begin{enumerate}
\item $x\in {\mathfrak m}^d$ が $Gr_{\mathfrak m}(A)$ の次数 $d$ にある
非零因子 $\bar x \in {\mathfrak m}^d/{\mathfrak m}^{d + 1}$
に写ると仮定する．$x$ が非零因子であることを示せ．
\item $A$ の深さは少なくとも $1$ であると仮定する．
すなわち，非零因子 $y \in {\mathfrak m}$ が存在すると仮定する．
この場合にはさらに強い結果が得られる．$x\in {\mathfrak m}^d$ が
$Gr_{\mathfrak m}(A)$ の次数 $d$ にある元
$\bar x \in {\mathfrak m}^d/{\mathfrak m}^{d + 1}$ に写り，
これが十分高い次数で非零因子であることだけを仮定する．
つまり，$\exists N$ であって，すべての $n \geq N$ に対し
$\bar x$ を掛ける写像
$$
{\mathfrak m}^n/{\mathfrak m}^{n + 1} \longrightarrow
{\mathfrak m}^{n + d}/{\mathfrak m}^{n + d + 1}
$$
が単射である．このとき $x$ が非零因子であることを示せ．
\end{enumerate}
\end{exercise}

\begin{exercise}
\label{exercises-exercise-embedding-2-dim-1}
$(A, {\mathfrak m}, k)$ を次元 $1$ の Noether 局所環とする．
また $A$ の埋め込み次元が $2$，すなわち
$$
\dim_k {\mathfrak m}/{\mathfrak m}^2 = 2.
$$
であると仮定する．
記法：$f(n) = \dim_k {\mathfrak m}^n/{\mathfrak m}^{n + 1}$．
生成元 $x, y \in {\mathfrak m}$ を選び，ある斉次イデアル $I$ を用いて
$Gr_{\mathfrak m}(A) = k[\bar x, \bar y]/I$ と書く．
\begin{enumerate}
\item 斉次元 $F\in k[\bar x, \bar y]$ が存在し，$I \subset (F)$
であって，十分高いすべての次数では等号が成り立つことを示せ．
\item $f(n) \leq n + 1$ を示せ．
\item $f(n) < n + 1$ ならば $n \geq \deg(F)$ であることを示せ．
\item $f(n) < n + 1$ ならば $f(n + 1) \leq f(n)$ であることを示せ．
% P12-ERR-0443: preserve literal >> in the frozen formula.
\item すべての $n >> 0$ に対して $f(n) = \deg(F)$ であることを示せ．
\end{enumerate}
\end{exercise}

\begin{exercise}
\label{exercises-exercise-CM-dim-1-embedding-dim-2}
次元 1，埋め込み次元 2 の Cohen--Macaulay 環．
$(A, {\mathfrak m}, k)$ を次元 $1$ の Cohen--Macaulay
Noether 局所環とする．また $A$ の埋め込み次元が $2$，
すなわち
$$
\dim_k {\mathfrak m}/{\mathfrak m}^2 = 2.
$$
であると仮定する．
記法：$f$，$F$，$x, y\in {\mathfrak m}$，$I$ は
演習 \ref{exercises-exercise-embedding-2-dim-1} のとおりとする．
上の問題の結果はどれを使ってもよい．
\begin{enumerate}
\item $z\in {\mathfrak m}$ は，${\mathfrak m}/{\mathfrak m}^2$
における類が $F$ と互いに素な一次形式
$\alpha \bar x + \beta \bar y \in k[\bar x, \bar y]$
であるような元とする．
\begin{enumerate}
\item $z$ が $A$ 上の非零因子であることを示せ．
\item $d = \deg(F)$ とする．十分大きいすべての $n$ に対して
${\mathfrak m}^n = z^{n + 1-d}{\mathfrak m}^{d-1}$
となることを示せ．（ヒント：まず $Gr_{\mathfrak m}(A)$ について
知っていることから
$z^{n + 1-d}{\mathfrak m}^{d-1} \to {\mathfrak m}^n/{\mathfrak m}^{n + 1}$
が全射であることを示せ．次に NAK を使え．）
\end{enumerate}
\item $k$ にどのような条件を課せば，このような $z$ の存在が保証されるか．
（証明は不要．あまりに簡単だからである．）

\noindent
ここからは上のような $z$ が存在すると仮定する．
これは無害な仮定であることが分かる
（以下の (d)，(e) の結果を得るためには，
この仮定の成り立つ状況に帰着できるという意味で）．
\item すべての $\ell \geq d$ に対して
${\mathfrak m}^\ell = z^{\ell - d + 1} {\mathfrak m}^{d-1}$
であることを示せ．
\item $I = (F)$ と結論せよ．
% P12-ERR-0435: the frozen sequence starts at 2 without an explicit starting index.
\item 函数 $f$ の値が
$$
2, 3, 4, \ldots, d-1, d, d, d, d, d, d, d, \ldots
$$
となると結論せよ．
\end{enumerate}
\end{exercise}

\begin{remark}
\label{exercises-remark-CM-dim-1-embedding-dim-2}
これは，次元 1，埋め込み次元 2 の Noether 局所 Cohen--Macaulay 環が，
2 次元正則局所環 $B$ を用いて $B/FB$ の形に書けることを示唆する．
これはほぼ正しい（環について適切な「良さ」の条件の下で）．
\end{remark}

% Frozen source: exercises.tex lines 2237--2359.
\section{無限個の素数}
\label{exercises-section-many-primes}

\noindent
無限個の素数が可逆でない環について，奇妙な問いを集めた節である．

\begin{exercise}
\label{exercises-exercise-not-in-Q}
次の二つの性質を持つ有限型 ${\mathbf Z}$-代数 $R$ の例を与えよ．
\begin{enumerate}
\item 環準同型 $R \to {\mathbf Q}$ は存在しない．
\item 各素数 $p$ に対して極大イデアル
${\mathfrak m} \subset R$ が存在し，$R/{\mathfrak m} \cong {\mathbf F}_p$
となる．
\end{enumerate}
\end{exercise}

\begin{exercise}
\label{exercises-exercise-strange-fp-1}
$f \in {\mathbf Z}[x, u]$ に対し，$f_p(x)
= f(x, x^p) \bmod p \in {\mathbf F}_p[x]$ と定める．
次の二つの性質を持つ $f \in {\mathbf Z}[x, u]$ の例を与えよ．
\begin{enumerate}
\item $f_p$ が ${\mathbf F}_p$ に零点を持たないような
$p$ が無限個存在する．
% P12-ERR-0443: preserve literal >> in the frozen formula.
\item すべての $p >> 0$ に対して，多項式 $f_p$ は
一次または二次の因子を持つ．
\end{enumerate}
\end{exercise}

\begin{exercise}
\label{exercises-exercise-strange-fp-2}
$f \in {\mathbf Z}[x, y, u, v]$ に対し，$f_p(x, y)
= f(x, y, x^p, y^p) \bmod p \in {\mathbf F}_p[x, y]$ と定める．
% P12-ERR-0443: preserve literal >> in the frozen formula.
すべての $p >> 0$ に対して $f_p$ が可約となるような $f$ の
「興味深い」例を与えよ．
例えば，$f = xv-yu$ として $f_p = xy^p-x^py = xy(x^{p-1}-y^{p-1})$
となる例は「面白くない」．$x, u$ だけに依存する $f$ も
「面白くない」，等々．
\end{exercise}

\begin{remark}
\label{exercises-remark-strange-fp}
$h \in {\mathbf Z}[y]$ を次数 $d$ のモニック多項式とする．このとき：
\begin{enumerate}
\item 写像 $A = {\mathbf Z}[x] \to B ={\mathbf Z}[y]$，
$x \mapsto h$ は階数 $d$ の有限局所自由な写像である．
% P12-ERR-0436: preserve the frozen domain-generator token y.
\item すべての素数 $p$ に対して，写像
$A_p = {\mathbf F}_p[x]\to B_p = {\mathbf F}_p[y]$，
$y \mapsto h(y) \bmod p$ は階数 $d$ の有限局所自由な写像である．
\end{enumerate}
\end{remark}

\begin{exercise}
\label{exercises-exercise-strange-fp-3}
$h, A, B, A_p, B_p$ は注意でのものとする．
$f \in {\mathbf Z}[x, u]$ に対し
$f_p(x) = f(x, x^p) \bmod p \in {\mathbf F}_p[x]$ と定める．
$g \in {\mathbf Z}[y, v]$ に対しては
$g_p(y) = g(y, y^p) \bmod p \in {\mathbf F}_p[y]$ と定める．
\begin{enumerate}
\item 次の性質を持つ $f$ が存在しないような $h$ と $g$ の例を与えよ．
% P12-ERR-0437: preserve bare Norm in the frozen formula.
$$
f_p  =  Norm_{B_p/A_p}(g_p).
$$
\item 上のような $h$ と $g$ の任意の選択に対して，
零でない $f$ が存在し，すべての $p$ に対して
% P12-ERR-0437: preserve bare Norm; only the text payload is translated.
$$
Norm_{B_p/A_p}(g_p)\quad\text{は次を割り切る}\quad f_p .
$$
となることを示せ．望むなら $h = y^n$ の場合，
さらには $n = 2$ の場合に制限してもよいが，これは一般に成り立つ．
% P12-ERR-0438: preserve the frozen manual exercise numbers.
\item これと前の問題群の演習 6，7 との関係を論ぜよ．
\end{enumerate}
\end{exercise}

\begin{exercise}
\label{exercises-exercise-strange-fp-unsolved}
未解決問題．本当に難しいかもしれないし，易しいかもしれない．
私には分からない．
\begin{enumerate}
\item 無限個の $p$ に対して $f_p$ が既約となる
$f \in {\mathbf Z}[x, u]$ は存在するか．
（ヒント：存在する．$f(x, u) = u - x - 1$ でも
$f(x, u) = u^2 - x^2 + 1$ でもそうなる．）
\item $f \in {\mathbf Z}[x, u]$ を零でないものとし，
十分大きいすべての $p$ に対して $\deg_x(f_p) = dp + d'$ と仮定する．
（言い換えると $\deg_u(f) = d$ であり，
$f$ における $u^d$ の係数 $c$ は $\deg_x(c) = d'$ を満たす．）
$d = d_1 + d_2$ および $d' = d'_1 + d'_2$ と書け，
$d_1, d_2 > 0$，$d'_1, d'_2 \geq 0$ であって，
十分大きいすべての $p$ に対して因数分解
$$
f_p = f_{1, p} f_{2, p}
$$
が存在し，$\deg_x(f_{1, p}) = d_1p + d'_1$ となると仮定する．
% P12-ERR-0439: preserve the frozen manual reference to Exercise 4.
$f$ は演習 4 のようなノルムの構成から得られるというのは正しいか．
% P12-ERR-0437: preserve bare Norm in the frozen formula.
% P12-ERR-0443: preserve literal >> in the frozen formula.
（より正確には，すべての $p >> 0$ に対して
$Norm_{B_p/A_p}(g_p)$ が $f_p$ を割り切るような $h$ と $g$ は存在するか．）
\item (b) と同様の問いを，今度は
$f \in {\mathbf Z}[x_1, x_2, u_1, u_2]$ が既約であり，
% P12-ERR-0443: preserve literal >> in the frozen formula.
すべての $p >> 0$ に対して
$f_p(x_1, x_2) = f(x_1, x_2, x_1^p, x_2^p) \bmod p$
が因数分解するということだけを仮定して考えよ．
\end{enumerate}
\end{exercise}

% Frozen source: exercises.tex lines 2360--2571.
\section{フィルター付き導来圏}
\label{exercises-section-filtered-derived}

\noindent
この節の演習を解くには，ホモロジー，節 \ref{homology-section-filtrations}
の内容を読んでほしい．$A$ が $\mathcal{A}$ のフィルター付き対象であるというとき，
$A$ には記法から省略したフィルトレーション $F$ が与えられているものとする．

\begin{exercise}
\label{exercises-exercise-split-injective}
$\mathcal{A}$ をアーベル圏とする．
$I$ を $\mathcal{A}$ のフィルター付き対象とする．
$I$ のフィルトレーションは有限であり，各 $\text{gr}^p(I)$ は
$\mathcal{A}$ の入射的対象であると仮定する．
同型 $I \cong \bigoplus \text{gr}^p(I)$ が存在し，
% P12-ERR-0440: preserve the frozen repeated exponent p in the summand.
フィルトレーション $F^p(I)$ が $\bigoplus_{p' \geq p} \text{gr}^p(I)$
に対応することを示せ．
\end{exercise}

\begin{exercise}
\label{exercises-exercise-filtered-injective}
$\mathcal{A}$ をアーベル圏とする．
$I$ を $\mathcal{A}$ のフィルター付き対象とする．
$I$ のフィルトレーションは有限であると仮定する．
次が同値であることを示せ．
\begin{enumerate}
\item フィルター付き対象の実線の図式
$$
\xymatrix{
A \ar[r]_\alpha \ar[d] & B \ar@{-->}[ld] \\
I &
}
$$
% P12-ERR-0441: supply the Japanese predicate without changing the conditions.
であって，(\romannumeral1) $A$ と $B$ のフィルトレーションが有限であり，
(\romannumeral2) $\text{gr}(\alpha)$ が単射であるような任意のものに対して，
図式を可換にする点線の射が存在する．
\item 各 $\text{gr}^p I$ は入射的である．
\end{enumerate}
\end{exercise}

\noindent
% P12-ERR-0441: the source predicate omits "is".
有限なフィルトレーションを持つ対象の射 $\alpha : A \to B$ が与えられたとき，
$\text{gr}(\alpha)$ が単射であるということは，$\alpha$ が
圏 $\text{Fil}(\mathcal{A})$ における{\it 厳密なモノ射}であるということと
同じであることに注意しよう．すなわち，モノ射であることは $\Ker(\alpha) = 0$
を意味し，厳密であることは，これから $\Ker(\text{gr}(\alpha)) = 0$
も従うことを意味する．
ホモロジー，補題 \ref{homology-lemma-characterize-strict} を参照．
（核を持つ任意の加法圏でも意味をなすが，
射について「単射」という語はアーベル圏においてのみ用いる．）
% P12-ERR-0442: source subject-verb agreement error.
上の演習は次の定義を正当化する．

\begin{definition}
\label{exercises-definition-injective-filtered}
$\mathcal{A}$ をアーベル圏とする．
$I$ を $\mathcal{A}$ のフィルター付き対象とする．
$I$ のフィルトレーションは有限であると仮定する．
各 $\text{gr}^p(I)$ が $\mathcal{A}$ の入射的対象であるとき，
$I$ は{\it フィルター付き入射的}であるという．
\end{definition}

\noindent
至るところで「有限なフィルトレーションを持つ」と繰り返さずに済むように，
次の定義を行う．

\begin{definition}
\label{exercises-definition-finite-filtration-category}
$\mathcal{A}$ をアーベル圏とする．
$\text{Fil}(\mathcal{A})$ の充満部分圏で，
フィルトレーションが有限であるような
$A \in \Ob(\text{Fil}(\mathcal{A}))$ を対象とするものを
{\it $\text{Fil}^f(\mathcal{A})$} と表す．
\end{definition}

\begin{exercise}
\label{exercises-exercise-inject-into-injective}
$\mathcal{A}$ をアーベル圏とする．
$\mathcal{A}$ は十分多くの入射的対象を持つと仮定する．
$A$ を $\text{Fil}^f(\mathcal{A})$ の対象とする．
$A$ から $\text{Fil}^f(\mathcal{A})$ のフィルター付き入射的対象 $I$ への
厳密なモノ射 $\alpha : A \to I$ が存在することを示せ．
\end{exercise}

\begin{definition}
\label{exercises-definition-filtered-quasi-isomorphism}
$\mathcal{A}$ をアーベル圏とする．
$\alpha : K^\bullet \to L^\bullet$ を
$\text{Fil}(\mathcal{A})$ の複体の射とする．
各 $p \in \mathbf{Z}$ に対して射
$\text{gr}^p(K^\bullet) \to \text{gr}^p(L^\bullet)$ が擬同型であるとき，
$\alpha$ は{\it フィルター付き擬同型}であるという．
\end{definition}

\begin{definition}
\label{exercises-definition-filtered-acyclic}
$\mathcal{A}$ をアーベル圏とする．
$K^\bullet$ を $\text{Fil}^f(\mathcal{A})$ の複体とする．
各 $p \in \mathbf{Z}$ に対して複体 $\text{gr}^p(K^\bullet)$ が非輪状であるとき，
$K^\bullet$ は{\it フィルター付き非輪状}であるという．
\end{definition}

\begin{exercise}
\label{exercises-exercise-filtered-quasi-isomorphism}
$\mathcal{A}$ をアーベル圏とする．
$\alpha : K^\bullet \to L^\bullet$ を
$\text{Fil}^f(\mathcal{A})$ の下に有界な複体の射とする．
（上付き添字 $f$ に注意．）
次が同値であることを示せ．
\begin{enumerate}
\item $\alpha$ はフィルター付き擬同型である．
\item 各 $p \in \mathbf{Z}$ に対して，
写像 $\alpha : F^pK^\bullet \to F^pL^\bullet$ は擬同型である．
\item 各 $p \in \mathbf{Z}$ に対して，写像
$\alpha : K^\bullet/F^pK^\bullet \to L^\bullet/F^pL^\bullet$
は擬同型である．
\item $\alpha$ の錐（導来圏，定義 \ref{derived-definition-cone} を参照）
はフィルター付き非輪状複体である．
\end{enumerate}
さらに，$\alpha$ がフィルター付き擬同型ならば，
$\alpha$ は通常の意味でも擬同型であることを示せ．
\end{exercise}

\begin{exercise}
\label{exercises-exercise-injective-resolution}
$\mathcal{A}$ をアーベル圏とする．
$\mathcal{A}$ は十分多くの入射的対象を持つと仮定する．
$A$ を $\text{Fil}^f(\mathcal{A})$ の対象とする．
$\text{Fil}^f(\mathcal{A})$ の複体 $I^\bullet$ と
射 $A[0] \to I^\bullet$ であって，次を満たすものが存在することを示せ．
\begin{enumerate}
\item 各 $I^p$ はフィルター付き入射的である．
\item $p < 0$ に対して $I^p = 0$ である．
\item $A[0] \to I^\bullet$ はフィルター付き擬同型である．
\end{enumerate}
\end{exercise}

\begin{exercise}
\label{exercises-exercise-injective-resolution-complex}
$\mathcal{A}$ をアーベル圏とする．
$\mathcal{A}$ は十分多くの入射的対象を持つと仮定する．
$K^\bullet$ を $\text{Fil}^f(\mathcal{A})$ の対象からなる下に有界な複体とする．
フィルター付き擬同型 $\alpha : K^\bullet \to I^\bullet$ が存在し，
$I^\bullet$ は $\text{Fil}^f(\mathcal{A})$ の下に有界な複体で，
その各項 $I^n$ がフィルター付き入射的であることを示せ．
実際，$\alpha$ は，各 $\alpha^n$ が厳密なモノ射となるように選べる．
\end{exercise}

\begin{exercise}
\label{exercises-exercise-morphisms-lift}
$\mathcal{A}$ をアーベル圏とする．
$\text{Fil}^f(\mathcal{A})$ の複体の実線の図式
$$
\xymatrix{
K^\bullet \ar[r]_\alpha \ar[d]_\gamma & L^\bullet \ar@{-->}[dl]^\beta \\
I^\bullet
}
$$
を考える．$K^\bullet$，$L^\bullet$，$I^\bullet$ は下に有界で，
各 $I^n$ はフィルター付き入射的対象であると仮定する．
また $\alpha$ はフィルター付き擬同型であると仮定する．
\begin{enumerate}
\item 図式をホモトピーを除いて可換にする複体の射 $\beta$ が存在する．
\item $\alpha$ が各次数で厳密なモノ射ならば，
図式を可換にする $\beta$ を見つけることができる．
\end{enumerate}
\end{exercise}

\begin{exercise}
\label{exercises-exercise-acyclic-is-zero}
$\mathcal{A}$ をアーベル圏とする．
% P12-ERR-0444: preserve the duplicated K in the frozen introduction.
$K^\bullet$，$K^\bullet$ を $\text{Fil}^f(\mathcal{A})$ の複体とする．
次を仮定する．
\begin{enumerate}
\item $K^\bullet$ は下に有界で，フィルター付き非輪状である．
\item $I^\bullet$ は下に有界で，フィルター付き入射的対象からなる．
\end{enumerate}
このとき，任意の射 $K^\bullet \to I^\bullet$ は零とホモトピックである．
\end{exercise}

\begin{exercise}
\label{exercises-exercise-morphisms-equal-up-to-homotopy}
$\mathcal{A}$ をアーベル圏とする．
$\text{Fil}^f(\mathcal{A})$ の複体の実線の図式
$$
\xymatrix{
K^\bullet \ar[r]_\alpha \ar[d]_\gamma & L^\bullet \ar@{-->}[dl]^{\beta_i} \\
I^\bullet
}
$$
を考える．$K^\bullet$，$L^\bullet$，$I^\bullet$ は下に有界であり，
各 $I^n$ はフィルター付き入射的対象であると仮定する．
また $\alpha$ はフィルター付き擬同型であると仮定する．
図式をホモトピーを除いて可換にする任意の二つの射
$\beta_1, \beta_2$ はホモトピックである．
\end{exercise}

% Frozen source: exercises.tex lines 2572--2649.
\section{正則函数}
\label{exercises-section-regular-functions}

\begin{exercise}
\label{exercises-exercise-extra-function}
座標 $s, t$ を持つ $\mathbf{C}^2$ において，
方程式 $t^2 = s^5 + 8$ で与えられるアファイン曲線 $X$ を考える．
$x \in X$ を座標 $(1, 3)$ を持つ点とする．
$U = X \setminus \{x\}$ とする．
$U$ 上の正則函数であって，$\mathbf{C}^2$ 上の正則函数の制限ではないもの，
すなわち $s$ と $t$ の多項式の $U$ への制限ではないものが存在することを証明せよ．
\end{exercise}

\begin{exercise}
\label{exercises-exercise-no-extra-function}
$n \geq 2$ とする．$E \subset \mathbf{C}^n$ を有限部分集合とする．
$\mathbf{C}^n \setminus E$ 上の任意の正則函数が多項式であることを示せ．
\end{exercise}

\begin{exercise}
\label{exercises-exercise-cone}
$X \subset \mathbf{C}^n$ をアファイン多様体とする．
$x = (a_1, \ldots, a_n) \in X$ および $\lambda \in \mathbf{C}$ から
$(\lambda a_1, \ldots, \lambda a_n) \in X$ が従うとき，
$X$ は{\it 錐}であるということにする．
もちろん，$\mathfrak p \subset \mathbf{C}[x_1, \ldots, x_n]$ が
$x_1, \ldots, x_n$ の斉次多項式で生成される素イデアルならば，
アファイン多様体 $X = V(\mathfrak p) \subset \mathbf{C}^n$ は錐である．
逆に，錐に対応する素イデアル
$\mathfrak p \subset \mathbf{C}[x_1, \ldots, x_n]$ は
$x_1, \ldots, x_n$ の斉次多項式で生成できることを示せ．
\end{exercise}

\begin{exercise}
\label{exercises-exercise-extra-function-cone}
錐であるようなアファイン多様体 $X \subset \mathbf{C}^n$
（演習 \ref{exercises-exercise-cone} を参照）と，
$U = X \setminus \{(0, \ldots, 0)\}$ 上の正則函数 $f$ であって，
$\mathbf{C}^n$ 上の多項式函数の制限ではないものの例を与えよ．
\end{exercise}

\begin{exercise}
\label{exercises-exercise-regular-functions}
この演習では，代数閉体でない体上の正則函数に何が起こるかを調べる．
$k$ を体とする．$Z \subset k^n$ を Zariski 局所閉部分集合，
すなわちイデアル $I \subset J \subset k[x_1, \ldots, x_n]$ が存在して
$$
Z = \{a \in k^n \mid
f(a) = 0\ \forall\ f \in I,\ \exists\ g \in J,\ g(a) \not = 0\}.
$$
となるものとする．
函数 $\varphi : Z \to k$ が{\it 正則}であるとは，
各 $z \in Z$ に対して Zariski 開近傍 $z \in U \subset Z$ と
多項式 $f, g \in k[x_1, \ldots, x_n]$ が存在し，
すべての $u \in U$ に対して $g(u) \not = 0$ であり，
すべての $u \in U$ に対して $\varphi(u) = f(u)/g(u)$ となることをいう．
\begin{enumerate}
\item $k = \bar k$ かつ $Z = k^n$ ならば，
正則函数は多項式で与えられることを示せ．
（この議論を以前に見たことがなければ取り組めばよい．）
\item $k$ が有限ならば，(a) すべての函数 $\varphi$ は正則であり，
(b) 正則函数の環は $k$ 上有限次元であることを示せ．
（望むなら $Z = k^n$，さらには $n = 1$ としてもよい．）
\item $k = \mathbf{R}$ の場合，$Z = \mathbf{R}$ 上の正則函数で，
多項式では与えられないものの例を与えよ．
\item $k = \mathbf{Q}_p$ の場合，$Z = \mathbf{Q}_p$ 上の正則函数で，
多項式では与えられないものの例を与えよ．
\end{enumerate}
\end{exercise}

% Frozen source: exercises.tex lines 2650--2810.
\section{層}
\label{exercises-section-sheaves}

\noindent
圏 $\mathcal{C}$ の射 $f : X \to Y$ が{\it モノ射}であるとは，
任意の射の対 $a, b : W \to X$ に対して
$f \circ a = f \circ b \Rightarrow a = b$ が成り立つことをいう．
集合の圏におけるモノ射は集合の単射写像である．

\begin{exercise}
\label{exercises-exercise-mono-sheaves-sets}
% P12-ERR-0445: source indefinite-article error.
集合の層の写像が（集合の層の圏における）モノ射であることと，
すべての茎に誘導される写像が単射であることが同値であることを，丁寧に証明せよ．
\end{exercise}

\noindent
圏 $\mathcal{C}$ の射 $f : X \to Y$ が{\it 同型}であるとは，
射 $g : Y \to X$ が存在して $f \circ g = \text{id}_Y$ かつ
$g \circ f = \text{id}_X$ となることをいう．
集合の圏における同型は集合の全単射写像である．

\begin{exercise}
\label{exercises-exercise-isomorphism-sheaves-sets}
集合の層の写像が（集合の層の圏における）同型であることと，
すべての茎に誘導される写像が全単射であることが同値であることを，丁寧に証明せよ．
\end{exercise}

\noindent
圏 $\mathcal{C}$ の射 $f : X \to Y$ が{\it エピ射}であるとは，
任意の射の対 $a, b : Y \to Z$ に対して
$a \circ f = b \circ f \Rightarrow a = b$ が成り立つことをいう．
集合の圏におけるエピ射は集合の全射写像である．

\begin{exercise}
\label{exercises-exercise-epi-sheaves-sets}
集合の層の写像が（集合の層の圏における）エピ射であることと，
すべての茎に誘導される写像が全射であることが同値であることを，丁寧に証明せよ．
\end{exercise}

\begin{exercise}
\label{exercises-exercise-adjoint-push-pull}
$f : X \to Y$ を位相空間の写像とする．
{\bf 集合}の層に対する順像 $f_\ast$ と逆像 $f^{-1}$ が
随伴関手の対をなすことを証明せよ．
\end{exercise}

\begin{exercise}
\label{exercises-exercise-j-shriek}
$j : U \to X$ を開埋め込みとする．次を示せ．
\begin{enumerate}
\item 逆像 $j^{-1} : \Sh(X) \to \Sh(U)$ は，
{\it 空集合による延長}と呼ばれる左随伴 $j_{!} : \Sh(U) \to \Sh(X)$ を持つ．
\item $\mathcal{G} \in \Sh(U)$ に対して
$j_{!}({\mathcal G})$ の茎を特徴付けよ．
\item 逆像 $j^{-1} : \textit{Ab}(X) \to \textit{Ab}(U)$ は，
{\it 零による延長}と呼ばれる左随伴
$j_{!} : \textit{Ab}(U) \to \textit{Ab}(X)$ を持つ．
\item $\mathcal{G} \in \textit{Ab}(U)$ に対して
$j_{!}({\mathcal G})$ の茎を特徴付けよ．
\end{enumerate}
零による延長は空集合による延長とは異なることに注意せよ！
\end{exercise}

\begin{exercise}
\label{exercises-exercise-not-locally-generated-by-sections}
$X = \mathbf{R}$ に通常の位相を入れる．
$\mathcal{O}_X = \underline{\mathbf{Z}/2\mathbf{Z}}_X$ とする．
$i : Z = \{0\} \to X$ を包含写像とし，
$\mathcal{O}_Z = \underline{\mathbf{Z}/2\mathbf{Z}}_Z$ とする．
次を証明せよ（最初の三つは定義から従うが，
定義が明確でなければ，それらを明らかにすべきである）．
\begin{enumerate}
\item $i_*\mathcal{O}_Z$ は摩天楼層である．
\item 標準的な全射写像
$\underline{\mathbf{Z}/2\mathbf{Z}}_X \to
i_*\underline{\mathbf{Z}/2\mathbf{Z}}_Z$ が存在する．
その核を $\mathcal{I} \subset \mathcal{O}_X$ と表す．
\item $\mathcal{I}$ は $\mathcal{O}_X$ のイデアル層である．
\item $X$ 上の層 $\mathcal{I}$ は，
（加群，定義 \ref{modules-definition-locally-generated} の意味で）
切断によって局所的に生成することができない．
\end{enumerate}
\end{exercise}

\begin{exercise}
\label{exercises-exercise-quotient-j-shriek-Z}
$X$ を位相空間とする．
${\mathcal F}$ を $X$ 上のアーベル群の層とする．
${\mathcal F}$ は，各項が
$$
j_{!}(\underline{{\mathbf Z}}_U)
$$
の形をした層の（非常に大きいかもしれない）直和の商であることを示せ．
ここで $U \subset X$ は開であり，
$\underline{{\mathbf Z}}_U$ は $U$ 上の値 ${\mathbf Z}$ の定数層を表す．
\end{exercise}

\begin{remark}
\label{exercises-remark-direct-sum-stalk-abelian}
$X$ を位相空間とする．
アーベル群の層の圏における層の族 $\{{\mathcal F}_i\}_{i\in I}$ の直和は，
% P12-ERR-0446: preserve the unindexed direct sum in the frozen presheaf.
前層 $U \mapsto \oplus {\mathcal F}_i(U)$ に付随する層である．
したがって，直和の点 $x$ における茎は，
${\mathcal F}_i$ の $x$ における茎の直和である．
\end{remark}

\begin{exercise}
\label{exercises-exercise-product-over-points}
$X$ を位相空間とする．
$x \in X$ で添字付けられたアーベル群の集まり $A_x$ が与えられたと仮定する．
対応 $U \mapsto \prod_{x \in U} A_x$ と明らかな制限写像によって，
アーベル群の層 $\mathcal{G}$ が定まることを示せ．
通常は $x \in X$ に対して $\mathcal{G}_x = A_x$ となるわけではないことを，
例によって示せ．
\end{exercise}

\begin{exercise}
\label{exercises-exercise-modified-product-over-points}
$X$，$A_x$，$\mathcal{G}$ は演習 \ref{exercises-exercise-product-over-points}
におけるものとする．
$\mathcal{B}$ を $X$ の位相の基とする．
位相，定義 \ref{topology-definition-base} を参照．
$U \in \mathcal{B}$ に対し，$A_U$ を部分群
$A_U \subset \mathcal{G}(U) = \prod_{x \in U} A_x$ とする．
$U, V \in \mathcal{B}$ が $U \subset V$ を満たすとき，
制限写像は $A_V$ を $A_U$ に写すと仮定する．
開集合 $U \subset X$ に対して
$$
\mathcal{F}(U) =
\left\{
(s_x)_{x \in U}
\middle|
\begin{matrix}
\text{各 }x\text{（}U\text{ の元）に対し，ある } V \in \mathcal{B} \\
x \in V \subset U\text{ が存在して } (s_y)_{y \in V} \in A_V
\end{matrix}
\right\}
$$
とおく．$\mathcal{F}$ が $X$ 上のアーベル群の層を定めることを示せ．
通常は $U \in \mathcal{B}$ に対して $\mathcal{F}(U) = A_U$
となるわけではないことを，例によって示せ．
\end{exercise}

\begin{exercise}
\label{exercises-exercise-exact-but-not-a-stalk-functor}
位相空間 $X$ と関手
$$
F : \Sh(X) \longrightarrow \textit{Sets}
$$
であって，完全である
（有限積および等化子と可換し，有限余積および余等化子と可換する．
圏，節 \ref{categories-section-exact-functor} を参照）が，
$F$ が茎関手 $\mathcal{F} \mapsto \mathcal{F}_x$ と同型となるような点
$x \in X$ は存在しないものの例を与えよ．
\end{exercise}

% Frozen source: exercises.tex lines 2811--2969.
\section{スキーム}
\label{exercises-section-schemes}

\noindent
$LRS$ を局所環付き空間の圏とする．
アファインスキームとは，$LRS$ の対象であって，
$(\Spec(A), {\mathcal O}_{\Spec(A)})$ の形の対と $LRS$ において同型なものである．
スキームとは，$LRS$ の対象 $(X, {\mathcal O}_X)$ であって，
各点 $x\in X$ が開近傍 $U \subset X$ を持ち，
対 $(U, {\mathcal O}_X|_U)$ がアファインスキームとなるものである．

\begin{exercise}
\label{exercises-exercise-one-point}
スキームでない $1$ 点の局所環付き空間を見つけよ．
\end{exercise}

\begin{exercise}
\label{exercises-exercise-two-points}
$X$ を，その台位相空間が 2 点を持つスキームとする．
$X$ がアファインスキームであることを示せ．
\end{exercise}

\begin{exercise}
\label{exercises-exercise-discrete-finite-set-points}
$X$ を，その台位相空間が有限離散集合であるスキームとする．
$X$ がアファインスキームであることを示せ．
\end{exercise}

\begin{exercise}
\label{exercises-exercise-three-points}
三つの点を持つ非アファインスキームが存在することを示せ．
\end{exercise}

\begin{exercise}
\label{exercises-exercise-quasi-compact-closed-point}
$X$ を空でない準コンパクトスキームとする．
$X$ が閉点を持つことを示せ．
\end{exercise}

\begin{remark}
\label{exercises-remark-open-immersion}
$(X, {\mathcal O}_X)$ が環付き空間で，$U \subset X$ が開部分集合ならば，
$(U, {\mathcal O}_X|_U)$ は環付き空間である．
記法：${\mathcal O}_U = {\mathcal O}_X|_U$．
% P12-ERR-0447: source singular/plural grammar.
環付き空間の標準的な射
$$
j : (U, {\mathcal O}_U) \longrightarrow (X, {\mathcal O}_X).
$$
が存在する．$(X, {\mathcal O}_X)$ が局所環付き空間ならば，
$(U, {\mathcal O}_U)$ もそうであり，$j$ は局所環付き空間の射である．
$(X, {\mathcal O}_X)$ がスキームならば，
$(U, {\mathcal O}_U)$ もそうであり，$j$ はスキームの射である．
$(U, {\mathcal O}_U)$ は $(X, {\mathcal O}_X)$ の{\it 開部分スキーム}であるといい，
$j$ は{\it 開埋め込み}であるという．より一般に，
上のような射 $j : (U, {\mathcal O}_U) \to (X, {\mathcal O}_X)$ と
{\it 同型}である任意の射
$j' : (V, {\mathcal O}_V) \to (X, {\mathcal O}_X)$ を開埋め込みと呼ぶ．
\end{remark}

\begin{exercise}
\label{exercises-exercise-open-affine-not-affine}
% P12-ERR-0448: preserve the frozen missing restriction subscript.
アファインスキーム $(X, {\mathcal O}_X)$ と開集合 $U \subset X$ であって，
$(U, {\mathcal O}_X|U)$ がアファインスキームでないものの例を与えよ．
\end{exercise}

\begin{exercise}
\label{exercises-exercise-morphism-does-not-extend}
% P12-ERR-0449: source imperative grammar.
アファインスキームの対 $(X, {\mathcal O}_X)$，$(Y, {\mathcal O}_Y)$，
$X$ の開部分スキーム $(U, {\mathcal O}_X|_U)$，
およびスキームの射 $(U, {\mathcal O}_X|_U) \to (Y, {\mathcal O}_Y)$ であって，
スキームの射 $(X, {\mathcal O}_X) \to (Y, {\mathcal O}_Y)$ に
延長できないものの例を与えよ．
\end{exercise}

\begin{exercise}
\label{exercises-exercise-closed-subscheme-does-not-extend}
（これはかなり難しい．）
% P12-ERR-0450: source imperative/article grammar.
スキーム $X$，開部分スキーム $U \subset X$，および閉部分スキーム
$Z \subset U$ であって，$Z$ が $X$ の閉部分スキームに延長できないものの例を与えよ．
\end{exercise}

\begin{exercise}
\label{exercises-exercise-not-morphism-schemes}
スキーム $X$，体 $K$，および環付き空間の射 $\Spec(K) \to X$ であって，
スキームの射ではないものの例を与えよ．
\end{exercise}

\begin{exercise}
\label{exercises-exercise-just-kidding}
\cite[Chapter II]{H} の第 1，2 節の演習をすべて解け……\ \ 冗談である！
\end{exercise}

\begin{definition}
\label{exercises-definition-integral}
スキーム $X$ が{\it 整}であるとは，$X$ が空でなく，
空でない任意のアファイン開集合 $U \subset X$ に対して，
環 $\Gamma(U, \mathcal{O}_X) = \mathcal{O}_X(U)$ が整域であることをいう．
\end{definition}

\begin{exercise}
\label{exercises-exercise-morphism-integral-schemes-surjective-stalks-not-closed}
{\it 整}スキームの射 $f : X \to Y$ であって，
誘導される写像 ${\mathcal O}_{Y, f(x)}
\to {\mathcal O}_{X, x}$ はすべての $x\in X$ に対して全射だが，
$f$ は閉埋め込みでないものの例を与えよ．
\end{exercise}

\begin{exercise}
\label{exercises-exercise-fibre-product-affines-not-affine}
ファイバー積 $X \times_S Y$ であって，
$X$ と $Y$ はアファインだが $X \times_S Y$ はそうでないものの例を与えよ．
\end{exercise}

\begin{remark}
\label{exercises-remark-separated-base-fibre-product-affines-affine}
$S$ が分離的ならば，これは起こりえないことが分かる．理由が分かるだろうか．
\end{remark}

\begin{exercise}
\label{exercises-exercise-not-geometrically-integral}
% P12-ERR-0451: source article/compound-modifier grammar.
${\mathbf Q}$ 上有限型の 1 次元整スキーム $V$ であって，
$\Spec({\mathbf C}) \times_{\Spec({\mathbf Q})} V$
が整でないものの例を与えよ．
\end{exercise}

\begin{exercise}
\label{exercises-exercise-not-geometrically-reduced}
% P12-ERR-0452: source article/compound-modifier grammar.
体 $k$ 上有限型の 1 次元整スキーム $V$ であって，
ある有限体拡大 $k'/k$ に対して $\Spec(k') \times_{\Spec(k)} V$
が被約でないものの例を与えよ．
\end{exercise}

\begin{remark}
\label{exercises-remark-affine-dimension}
スキームがアファインならば，次元は対応する環の Krull 次元と同じである．
したがって，前学期の結果を次元の計算に使うことができる．
% P12-ERR-0453: source possessive grammar.
\end{remark}

% Frozen source: exercises.tex lines 2970--3157.
\section{射}
\label{exercises-section-morphisms}

\noindent
射 $\pi : X \to S$ が与えられたとき，
それが切断または有理切断を持つかどうかは重要な問題である．
以下に例となる演習をいくつか挙げる．

\begin{exercise}
\label{exercises-exercise-no-section}
スキームの射
$$
\pi :
X = \Spec(\mathbf{C}[x, t, 1/xt])
\longrightarrow
S = \Spec(\mathbf{C}[t]).
$$
を考える．
\begin{enumerate}
\item $\pi \circ \sigma = \text{id}_S$ となる射
$\sigma : S \to X$ は存在しないことを示せ．
\item 空でない開集合 $U \subset S$ と射 $\sigma : U \to X$ であって，
$\pi \circ \sigma = \text{id}_U$ となるものは存在することを示せ．
\end{enumerate}
\end{exercise}

\begin{exercise}
\label{exercises-exercise-no-rational-section}
スキームの射
$$
\pi :
X = \Spec(\mathbf{C}[x, t]/(x^2 + t))
\longrightarrow
S = \Spec(\mathbf{C}[t]).
$$
を考える．
空でない開集合 $U \subset S$ と射 $\sigma : U \to X$ であって，
$\pi \circ \sigma = \text{id}_U$ となるものは存在しないことを示せ．
\end{exercise}

\begin{exercise}
\label{exercises-exercise-has-rational-section}
$A, B, C \in \mathbf{C}[t]$ を零でない多項式とする．
スキームの射
$$
\pi :
X = \Spec(\mathbf{C}[x, y, t]/(A + Bx^2 + Cy^2))
\longrightarrow
S = \Spec(\mathbf{C}[t]).
$$
を考える．
空でない開集合 $U \subset S$ と射 $\sigma : U \to X$ であって，
$\pi \circ \sigma = \text{id}_U$ となるものは存在することを示せ．
（ヒント：形式的に，ある $d$ に対して次数 $\leq d$ の
$X, Y, Z \in \mathbf{C}[t]$ を用いて $x = X/Z$，$y = Y/Z$ と書き，
これが方程式の解となる条件を求めよ．
% P12-ERR-0454: preserve literal >> in the frozen formula.
次に次元論を用いて，$d >> 0$ ならば方程式を満たす
零でない $X, Y, Z$ を見つけられることを示せ．）
\end{exercise}

\begin{remark}
\label{exercises-remark-tsen}
演習 \ref{exercises-exercise-has-rational-section} は「Tsen の定理」の特別な場合である．
演習 \ref{exercises-exercise-no-section-curve} は，
この方法が低い次数の方程式に限られることを示している
（底とファイバーの次元が 1 のときは二次曲線）．
\end{remark}

\begin{exercise}
\label{exercises-exercise-no-section-curve}
スキームの射
$$
\pi :
X = \Spec(\mathbf{C}[x, y, t]
/(1 + t x^3 + t^2 y^3))
\longrightarrow
S = \Spec(\mathbf{C}[t])
$$
を考える．
空でない開集合 $U \subset S$ と射 $\sigma : U \to X$ であって，
$\pi \circ \sigma = \text{id}_U$ となるものは存在しないことを示せ．
\end{exercise}

\begin{exercise}
\label{exercises-exercise-no-section-surface}
スキーム
$$
X = \Spec(\mathbf{C}[\{x_i\}_{i = 1}^{8}, s, t]
/(1 + s x_1^3 + s^2 x_2^3 +
t x_3^3 + st x_4^3 + s^2t x_5^3 +
t^2 x_6^3 + st^2 x_7^3 + s^2t^2 x_8^3))
$$
および
$$
S = \Spec(\mathbf{C}[s, t])
$$
とスキームの射
$$
\pi : X \longrightarrow S
$$
を考える．
空でない開集合 $U \subset S$ と射 $\sigma : U \to X$ であって，
$\pi \circ \sigma = \text{id}_U$ となるものは存在しないことを示せ．
\end{exercise}

\begin{exercise}
\label{exercises-exercise-for-number-theorists}
（数論を学ぶ人向け．）閉部分スキーム
$$
Z \subset \Spec\left({\mathbf Z}[x, \frac{1 }{ x(x-1)(2x-1)}]\right)
$$
であって，射 $Z \to \Spec({\mathbf Z})$ が有限かつ全射となるものの例を与えよ．
\end{exercise}

\begin{exercise}
\label{exercises-exercise-quasi-section}
数論が好みでなければ，次を見るという変種に取り組んでもよい．
$$
\Spec\left({\mathbf F}_p[t, x, \frac{1 }{ x(x-t)(tx-1)}]\right)
\longrightarrow
\Spec({\mathbf F}_p[t])
$$
% P12-ERR-0455: source adverbial grammar.
上のスキームの閉部分スキームで，下のスキームへ有限全射で写るものを見つけよ．
（ここで基礎体が有限であることには理論的な理由がある．
もっとも，この特別な場合には必要でないかもしれない．）
\end{exercise}

\begin{remark}
\label{exercises-remark-interpretation-skolem-noether}
演習 \ref{exercises-exercise-for-number-theorists} と \ref{exercises-exercise-quasi-section}
では，射 $f : X \to S$ が与えられ，
$S$ は Dedekind 環 $A$ のスペクトルであり，
$f$ は平坦かつ全射で，生成ファイバーは幾何学的に既約である．
いずれの場合にも，有限全射 $S' \to S$ が存在して，
底変換 $X_{S'} \to S'$ が切断を持つことになる．
$A$ が excellent であり，極大イデアルにおける剰余体が有限体の代数拡大であり，
$A'$ が $A$ の分数体の有限拡大における整閉包であるときに
$\Pic(A')$ が捩れ群であるならば，これは成り立つことが分かる．
\cite{MB1} を参照．
例えば $A = \mathbf{Z}$ または $A = \mathbf{F}_p[x]$，あるいはこれらの有限拡大の場合である．
しかし，極大素イデアルにおける剰余体が有限でありながら，
Picard 群が非捩れ元を持つ Dedekind 環 $A$ が存在することが分かる．
\cite{Goldman} を参照．そのような環のスペクトルについては結果は偽である．
演習 \ref{exercises-exercise-no-quasi-section} は，
これがうまくいかない幾何学的な例を与える．
\end{remark}

\begin{exercise}
\label{exercises-exercise-no-quasi-section}
% P12-ERR-0456: source existential-clause grammar.
どの $\alpha \in \mathbf{C}$ に対しても $t - \alpha$ で割り切れない
$f \in \mathbf{C}[x, t]$ であって，
$\Spec(\mathbf{C}[t])$ に有限全射で写る
$Z \subset \Spec(\mathbf{C}[x, t, 1/f])$ が存在しないものがあることを証明せよ．
（$f(x, t) = (xt - 2)(x - t + 3)$ でうまくいくと思う．
どの候補についても，求める性質を持つことを示すのはそれほど易しくないと思う．）
\end{exercise}

\begin{exercise}
\label{exercises-exercise-finite}
$A \to B$ を有限型環準同型とする．
$\Spec(B) \to \Spec(A)$ は，ある $n$ に対して
閉埋め込み $\Spec(B) \to \mathbf{P}^n_A$ を経由すると仮定する．
$A \to B$ が有限環準同型であること，
すなわち $B$ が $A$-加群として有限であることを証明せよ．
ヒント：$A$ が Noether 環ならば（そう仮定してよい），
任意の閉部分スキーム $Z \subset \mathbf{P}^n_A$ について，
$i \in \mathbf{Z}$ に対する $H^i(Z, \mathcal{O}_Z)$ が
有限 $A$-加群であることを用いて議論できる．
\end{exercise}

\begin{exercise}
\label{exercises-exercise-projective-finite}
$k$ を代数閉体とする．
$f : X \to Y$ を射影多様体の射であって，
すべての閉点 $y \in Y$ に対して $f^{-1}(\{y\})$ が有限であるものとする．
$f$ がスキームの射として有限であることを証明せよ．
ヒント：(a) 有限であることは局所的な性質である．
(b) 演習 \ref{exercises-exercise-finite} に帰着することを試みよ．
(c) 閉埋め込み $X \to \mathbf{P}^n_k$ を用いて，
$Y$ 上の閉埋め込み $X \to \mathbf{P}^n_Y$ を得よ．
\end{exercise}

% Frozen source: exercises.tex lines 3158--3326.
\section{接空間}
\label{exercises-section-tangent-space}

\begin{definition}
\label{exercises-definition-dual-numbers}
任意の環 $R$ に対して，{\it 双対数}の環を $R[\epsilon]$ と表す．
これは $R$-加群として基底 $1$，$\epsilon$ を持つ自由加群である．
環構造は $\epsilon^2 = 0$ とおくことで定まる．
\end{definition}

\begin{exercise}
\label{exercises-exercise-tangent-space-Zariski}
$f : X \to S$ をスキームの射とする．
$x \in X$ を点とし，$s = f(x)$ とする．
実線の可換図式
$$
\xymatrix{
\Spec(\kappa(x)) \ar[r] \ar[dr] \ar@/^1pc/[rr] &
\Spec(\kappa(x)[\epsilon]) \ar@{.>}[r] \ar[d]&
X \ar[d] \\
&
\Spec(\kappa(s)) \ar[r] &
S
}
$$
を考える．ここで曲がった矢印は $\Spec(\kappa(x))$ から $X$ への標準的な射である．
% P12-ERR-0457: source subject-verb agreement.
$\kappa(x) = \kappa(s)$ ならば，図式を可換にする点線の射の集合が
線型写像の集合
$$
\Hom_{\kappa(x)}(
\frac{\mathfrak m_x}{\mathfrak m_x^2 + \mathfrak m_s\mathcal{O}_{X, x}},
\kappa(x))
$$
と一対一に対応することを示せ．
言い換えると，そのような全単射を記述せよ．
（より一般に，$\kappa(x) \supset \kappa(s)$ が分離代数拡大の場合にも成り立つ．）
\end{exercise}

\begin{definition}
\label{exercises-definition-tangent-space}
$f : X \to S$ をスキームの射とする．
$x \in X$ とする．演習 \ref{exercises-exercise-tangent-space-Zariski}
の点線の射の集合を {\it $S$ 上の $X$ の接空間}と名付け，
$T_{X/S, x}$ と表す．この空間の元を
$X/S$ の $x$ における{\it 接ベクトル}と呼ぶ．
\end{definition}

\begin{exercise}
\label{exercises-exercise-simple-push-out}
任意の体 $K$ に対して，図式
$$
\xymatrix{
\Spec(K) \ar[r] \ar[d] & \Spec(K[\epsilon_1]) \ar[d] \\
\Spec(K[\epsilon_2]) \ar[r] &
\Spec(K[\epsilon_1, \epsilon_2]/(\epsilon_1\epsilon_2))
}
$$
がスキームの圏における押し出し図式であることを証明せよ．
（ここでも前と同様に $\epsilon_i^2 = 0$ である．）
\end{exercise}

\begin{exercise}
\label{exercises-exercise-tangent-space-vectors-space}
$f : X \to S$ をスキームの射とする．
$x \in X$ とする．演習 \ref{exercises-exercise-simple-push-out} と適切な射
$$
\Spec(K[\epsilon])
\longrightarrow
\Spec(K[\epsilon_1, \epsilon_2]/(\epsilon_1\epsilon_2)).
$$
を用いて，接ベクトルの加法を定義せよ．
同様に接ベクトルのスカラー倍を定義せよ（こちらの方が易しい）．
これらの構成によって $T_{X/S, x}$ が $\kappa(x)$-ベクトル空間となることを示せ．
\end{exercise}

\begin{exercise}
\label{exercises-exercise-compute-TS}
$k$ を体とする．構造射
$f : X = \mathbf{A}^1_k \to \Spec(k) = S$ を考える．
\begin{enumerate}
\item $x \in X$ を閉点とする．$T_{X/S, x}$ の次元はいくらか．
\item $\eta \in X$ を生成点とする．$T_{X/S, \eta}$ の次元はいくらか．
\item 今度は $X$ を $\Spec(\mathbf{Z})$ 上のスキームと考える．
$T_{X/\mathbf{Z}, x}$ と $T_{X/\mathbf{Z}, \eta}$ の次元はいくらか．
\end{enumerate}
\end{exercise}

\begin{remark}
\label{exercises-remark-tangent-space-relative}
演習 \ref{exercises-exercise-compute-TS} は，良い概念を得るには
$S$ 上の $X$ の接空間を考える必要がある理由を説明している．
\end{remark}

\begin{exercise}
\label{exercises-exercise-compute-TS-field}
スキームの射
$$
f : X = \Spec(\mathbf{F}_p(t))
\longrightarrow
\Spec(\mathbf{F}_p(t^p)) = S
$$
を考える．$X$ の唯一の点における $X/S$ の接空間を計算せよ．
奇妙ではないだろうか．$\mathbf{F}_p[t^p] \to \mathbf{F}_p[t]$
に対応するスキームの射をとると何が起こると思うか．
\end{exercise}

\begin{exercise}
\label{exercises-exercise-compute-TS-cusp}
$k$ を体とする．$X = \Spec(k[x, y]/(x^2 - y^3))$ としたとき，
点 $x = (0, 0)$ における $X/k$ の接空間を計算せよ．
\end{exercise}

\begin{exercise}
\label{exercises-exercise-map-tangent-spaces}
$f : X \to Y$ を $S$ 上のスキームの射とする．
$x \in X$ を点とし，$y = f(x)$ とおく．
自然な写像 $\kappa(y) \to \kappa(x)$ は全単射であると仮定する．
定義を用いて，$f$ が自然な線型写像
$$
\text{d}f : T_{X/S, x} \longrightarrow T_{Y/S, y}.
$$
を誘導することを示せ．
$\kappa(x) = \kappa(s)$ の場合に，
演習 \ref{exercises-exercise-tangent-space-Zariski} を通じて，
これが局所環上で起こることと対応することを確かめよ．
\end{exercise}

\begin{exercise}
\label{exercises-exercise-Jacobian}
$k$ を代数閉体とする．
% P12-ERR-0458: preserve the frozen single-symbol coordinate arguments.
\begin{eqnarray*}
f : \mathbf{A}_k^n & \longrightarrow & \mathbf{A}^m_k \\
(x_1, \ldots, x_n) & \longmapsto & (f_1(x_i), \ldots, f_m(x_i))
\end{eqnarray*}
を $k$ 上のスキームの射とする．これは $n$ 変数の
$m$ 個の多項式 $f_1, \ldots, f_m$ で与えられる．
% P12-ERR-0459: preserve the frozen Jacobian indexing/convention.
行列
$$
A = \left( \frac{\partial f_j}{\partial x_i} \right)
$$
を考える．$x \in \mathbf{A}^n_k$ を閉点とする．
$y =  f(x)$ とおく．接空間の写像
$T_{\mathbf{A}^n_k/k, x} \to T_{\mathbf{A}^m_k/k, y}$
は点 $x$ における行列 $A$ の値で与えられることを示せ．
\end{exercise}

% Frozen source: exercises.tex lines 3327--3443.
\section{準連接層}
\label{exercises-section-quasi-coherent}

\begin{definition}
\label{exercises-definition-quasi-coherent}
$X$ をスキームとする．
$\mathcal{O}_X$-加群の層 $\mathcal{F}$ が{\it 準連接}であるとは，
任意のアファイン開集合 $\Spec(R) = U \subset X$ に対して，
制限 $\mathcal{F}|_U$ が，ある $R$-加群 $M$ を用いて
$\widetilde M$ の形になることをいう．
\end{definition}

\noindent
% P12-ERR-0460: source demonstrative-number agreement.
この条件は $X$ のあるアファイン開被覆の各要素について調べれば十分である．
さらに多くの結果については，スキーム，節 \ref{schemes-section-quasi-coherent} を参照．

\begin{definition}
\label{exercises-definition-specialization}
$X$ を位相空間とする．$x, x' \in X$ とする．
$x \in \overline{\{x'\}}$ であるとき，またそのときに限り，
$x$ は $x'$ の{\it 特殊化}であるという．
\end{definition}

\begin{exercise}
\label{exercises-exercise-quasi-coherent-specialization-points}
$X$ をスキームとする．$x, x' \in X$ とする．
$\mathcal{F}$ を $\mathcal{O}_X$-加群の準連接層とする．
(a) $x$ は $x'$ の特殊化であり，
(b) $\mathcal{F}_{x'} \not = 0$ と仮定する．
$\mathcal{F}_x \not = 0$ を示せ．
\end{exercise}

\begin{exercise}
\label{exercises-exercise-O-module-specialization-points}
スキーム $X$，点 $x, x' \in X$，
および $\mathcal{O}_X$-加群の層 $\mathcal{F}$ であって，
(a) $x$ は $x'$ の特殊化であり，
(b) $\mathcal{F}_{x'} \not = 0$ かつ $\mathcal{F}_x = 0$
となるものの例を見つけよ．
\end{exercise}

\begin{definition}
\label{exercises-definition-Noetherian-scheme}
スキーム $X$ が{\it 局所 Noether}であるとは，
% P12-ERR-0544: preserve the missing explicit condition x in U.
各点 $x \in X$ に対してアファイン開集合 $\Spec(R) = U \subset X$
が存在し，$R$ が Noether 環となることをいう．
スキームが局所 Noether かつ準コンパクトであるとき，{\it Noether} であるという．
\end{definition}

\noindent
$X$ が局所 Noether ならば，$X$ の任意のアファイン開集合は
Noether 環のスペクトルである．
性質，補題 \ref{properties-lemma-locally-Noetherian} を参照．

\begin{definition}
\label{exercises-definition-coherent}
$X$ を局所 Noether スキームとする．
$\mathcal{F}$ を $\mathcal{O}_X$-加群の準連接層とする．
$\mathcal{F}$ が{\it 連接}であるとは，
% P12-ERR-0544: preserve the missing explicit condition x in U.
各点 $x \in X$ に対してアファイン開集合 $\Spec(R) = U \subset X$
が存在し，$\mathcal{F}|_U$ が，ある有限 $R$-加群 $M$ を用いて
$\widetilde M$ と同型になることをいう．
\end{definition}

\begin{exercise}
\label{exercises-exercise-extend-quasi-coherent}
$X = \Spec(R)$ をアファインスキームとする．
\begin{enumerate}
\item $f \in R$ とする．
$\mathcal{G}$ を開部分スキーム $D(f)$ 上の
$\mathcal{O}_{D(f)}$-加群の準連接層とする．
ある $\mathcal{O}_X$-加群の準連接層 $\mathcal{F}$ に対して
$\mathcal{G} = \mathcal{F}|_{D(f)}$ となることを示せ．
\item $I \subset R$ をイデアルとする．
$i : Z \to X$ を $I$ に対応する $X$ の閉部分スキームとする．
$\mathcal{G}$ を閉部分スキーム $Z$ 上の
$\mathcal{O}_Z$-加群の準連接層とする．
ある $\mathcal{O}_X$-加群の準連接層 $\mathcal{F}$ に対して
$\mathcal{G} = i^*\mathcal{F}$ となることを示せ．
（なぜこれは他愛もない問題なのか．）
\item $R$ が Noether 環であると仮定する．
$f \in R$ とする．
$\mathcal{G}$ を開部分スキーム $D(f)$ 上の
$\mathcal{O}_{D(f)}$-加群の連接層とする．
ある $\mathcal{O}_X$-加群の連接層 $\mathcal{F}$ に対して
$\mathcal{G} = \mathcal{F}|_{D(f)}$ となることを示せ．
\end{enumerate}
\end{exercise}

\begin{remark}
\label{exercises-remark-extend-off-open}
$U \to X$ が準コンパクトなはめ込みならば，
$U$ 上の任意の準連接層は $X$ 上の準連接層の制限である．
$X$ が Noether スキームで $U \subset X$ が開ならば，
$U$ 上の任意の連接層は $X$ 上の連接層の制限である．
もちろん，上の演習はより易しく，これらの一般的な事実を使うべきではない．
\end{remark}

% Frozen source: exercises.tex lines 3444--3598.
\section{Proj と射影スキーム}
\label{exercises-section-proj}

\begin{exercise}
\label{exercises-exercise-graded-ring-specified-result}
次を満たす次数付き環 $S$ の例を与えよ．
\begin{enumerate}
\item $\text{Proj}(S)$ はアファインで空でない．
\item $\text{Proj}(S)$ は整で空でないが，
どの $n\geq 0$ と環 $A$ に対しても ${\mathbf P}^n_A$ と同型でない．
\end{enumerate}
\end{exercise}

\begin{exercise}
\label{exercises-exercise-nonconstant-morphism-proj}
$\Spec({\mathbf C})$ 上のスキームの非定数射
${\mathbf P}^1_{\mathbf C} \to {\mathbf P}^5_{\mathbf C}$ の例を与えよ．
\end{exercise}

\begin{exercise}
\label{exercises-exercise-isomorphism-P1-conic}
スキームの同型
$$
{\mathbf P}^1_{\mathbf C} \to
\text{Proj}({\mathbf C}[X_0, X_1, X_2]/(X_0^2 + X_1^2 + X_2^2))
$$
の例を与えよ．
\end{exercise}

\begin{exercise}
\label{exercises-exercise-family-special-fibre-different}
スキームの射
$f : X \to {\mathbf A}^1_{\mathbf C} = \Spec({\mathbf C}[T])$
であって，$t \in {\mathbf A}^1_{\mathbf C}$ 上の $f$ の
（スキーム論的）ファイバー $X_t$ が，
(a) $t$ が $0$ と異なる閉点であるときには ${\mathbf P}^1_{\mathbf C}$ と同型であり，
(b) $t = 0$ のときには ${\mathbf P}^1_{\mathbf C}$ と同型でないものの例を与えよ．
$X_0$ をこの射の{\it 特殊ファイバー}と呼ぶことにする．
これは非常に多くの方法で実現できる．
以下の追加の制約を（それぞれ）満たす例を与えることを試みよ
（不可能な場合を除く）．
\begin{enumerate}
\item 特殊ファイバーを射影的にできるか．
\item 特殊ファイバーを既約かつ射影的にできるか．
\item 特殊ファイバーを整かつ射影的にできるか．
\item 特殊ファイバーを滑らかかつ射影的にできるか．
\item $f$ を平坦射にできるか．
これは，任意のアファイン開集合 $\Spec(A) \subset X$ に対して
誘導される環準同型 $\mathbf{C}[t] \to A$ が平坦であることを意味するだけであり，
この場合には，$t$ の任意の零でない多項式が $A$ 上の非零因子であることを意味する．
\item $f$ を平坦かつ射影的な射にできるか．
\item $f$ を平坦かつ射影的にし，特殊ファイバーを被約にできるか．
\item $f$ を平坦かつ射影的にし，特殊ファイバーを既約にできるか．
\item $f$ を平坦かつ射影的にし，特殊ファイバーを整にできるか．
\end{enumerate}
${\mathbf P}^1_{\mathbf C}$ を ${\mathbf C}$ 上の別の多様体に
置き換えると何が起こると思うか．
（上のどの変種を求めるかによっては，非常に難しくなる．）
\end{exercise}

\begin{exercise}
\label{exercises-exercise-affine-onto-projective-space}
$n \geq 1$ を任意の正の整数とする．
$X$ がアファインであるような全射 $X \to {\mathbf P}^n_{\mathbf C}$ の例を与えよ．
\end{exercise}

\begin{exercise}
\label{exercises-exercise-morphism-proj}
$\text{Proj}$ の写像．$R$ と $S$ を次数付き環とする．
環準同型
$$
\psi : R \to S
$$
と整数 $e \geq 1$ が与えられ，すべての $d \geq 0$ に対して
$\psi(R_d) \subset S_{de}$ が成り立つと仮定する．
（我々の規約では，$e = 1$ でない限りこれは次数付き環の準同型ではない．）
\begin{enumerate}
\item どの元 $\mathfrak p \in \text{Proj}(S)$ に対して，
$\text{Proj}(R)$ の対応する点が正しく定まるか．
言い換えると，$\psi$ が連続写像 $r_\psi : U \to \text{Proj}(R)$ を定めるような，
適切な開集合 $U \subset \text{Proj}(S)$ を求めよ．
\item $U \not = \text{Proj}(S)$ となる例を与えよ．
\item $U = \text{Proj}(S)$ となる例を与えよ．
\item （書き下さなくてよい．）
連続写像 $U \to \text{Proj}(R)$ には標準的に層の写像が付随し，
$r_\psi$ がスキームの射
$$
\text{Proj}(S) \supset U \longrightarrow \text{Proj}(R).
$$
となることを納得せよ．
\item $R = \bigoplus_{d \geq 0} S_{de}$
（$S_e$，$S_{2e}$，等々を次数 $1$，$2$，等々に持つ次数付き環として）で，
$\psi$ が包含写像である場合，この写像について何がいえるか．
\end{enumerate}
\end{exercise}

\noindent
{\bf 記法．} $R$ を上のような次数付き環とし，$n \geq 0$ を整数とする．
$X = \text{Proj}(R)$ とする．このとき，$X$ 上に一意的な
準連接 ${\mathcal O}_X$-加群 ${\mathcal O}_X(n)$ が存在し，
正の次数の任意の斉次元 $f \in R$ に対して，
% P12-ERR-0545: preserve the frozen untwisted O_X restriction.
${\mathcal O}_X |_{D_{+}(f)}$ は
$R_{(f)} = (R_f)_0$-加群 $(R_f)_n$ に付随する準連接層である．
% P12-ERR-0461: retain equality span, render the explanatory prose readably.
（$ = $ $R_f = R[1/f]$ の次数 $n$ の斉次元．）
\cite[page 116+]{H} を参照．自然な写像
$$
{\mathcal O}_X(n_1) \otimes_{{\mathcal O}_X} {\mathcal O}_X(n_2)
\longrightarrow
{\mathcal O}_X(n_1 + n_2)
$$
が存在することに注意せよ．

\begin{exercise}
\label{exercises-exercise-pathologies-proj}
$\text{Proj}$ における病理．
上のような $R$ で，次を満たすものの例を与えよ．
\begin{enumerate}
\item ${\mathcal O}_X(1)$ は可逆 ${\mathcal O}_X$-加群でない．
\item ${\mathcal O}_X(1)$ は可逆だが，自然な写像
${\mathcal O}_X(1) \otimes_{{\mathcal O}_X} {\mathcal O}_X(1) \to
{\mathcal O}_X(2)$ は同型でない．
\end{enumerate}
\end{exercise}

\begin{exercise}
\label{exercises-exercise-finitely-many-points-in-affine}
$S$ を次数付き環とする．$X = \text{Proj}(S)$ とする．
$X$ の任意の有限個の点からなる集合が，標準アファイン開集合に含まれることを示せ．
\end{exercise}

\begin{exercise}
\label{exercises-exercise-prepare-glueing}
$S$ を次数付き環とする．$X = \text{Proj}(S)$ とする．
$Z, Z' \subset X$ を二つの閉部分スキームとする．
$\varphi : Z \to Z'$ を同型とする．$Z \cap Z' = \emptyset$ と仮定する．
任意の $z \in Z$ に対してアファイン開集合 $U \subset X$ が存在し，
$z \in U$，$\varphi(z) \in U$，および
$\varphi(Z \cap U) = Z' \cap U$ となることを示せ．
（ヒント：演習 \ref{exercises-exercise-finitely-many-points-in-affine} と，
スキーム，補題 \ref{schemes-lemma-standard-open-two-affines}
に類似したものを用いよ．）
\end{exercise}

% Frozen source: exercises.tex lines 3599--3682.
\section{射影直線からの射}
\label{exercises-section-from-P1}

\noindent
この節では $\mathbf{P}^1$ から射影スキームへの射を調べる．

\begin{exercise}
\label{exercises-exercise-from-generic-point}
$k$ を体とする．$k[t] \subset k(t)$ を多項式環からその分数体への包含とする．
$X$ を $k$ 上有限型のスキームとする．
$k$ 上の任意の射
$$
\varphi : \Spec(k(t)) \longrightarrow X
$$
に対して，零でない $f \in k[t]$ と $k$ 上の射
$\psi : \Spec(k[t, 1/f]) \to X$ が存在し，
$\varphi$ が合成
$$
\Spec(k(t)) \longrightarrow \Spec(k[t, 1/f]) \longrightarrow X
$$
となることを示せ．
\end{exercise}

\begin{exercise}
\label{exercises-exercise-from-generic-point-Pn}
$k$ を体とする．$k[t] \subset k(t)$ を多項式環からその分数体への包含とする．
$k$ 上の任意の射
$$
\varphi : \Spec(k(t)) \longrightarrow \mathbf{P}^n_k
$$
に対して，$k$ 上の射 $\psi : \Spec(k[t]) \to \mathbf{P}^n_k$ が存在し，
$\varphi$ が合成
$$
\Spec(k(t)) \longrightarrow \Spec(k[t]) \longrightarrow \mathbf{P}^n_k
$$
となることを示せ．
ヒント：$\varphi$ の像は，ある $i$ に対して標準開集合 $D_+(T_i)$ に含まれる．
そこで「分母を払う」ことができると示せ．
\end{exercise}

\begin{exercise}
\label{exercises-exercise-from-generic-point-projective}
$k$ を体とする．$k[t] \subset k(t)$ を多項式環からその分数体への包含とする．
$X$ を $k$ 上の射影スキームとする．
$k$ 上の任意の射
$$
\varphi : \Spec(k(t)) \longrightarrow X
$$
に対して，$k$ 上の射 $\psi : \Spec(k[t]) \to X$ が存在し，
$\varphi$ が合成
$$
\Spec(k(t)) \longrightarrow \Spec(k[t]) \longrightarrow X
$$
となることを示せ．
ヒント：演習 \ref{exercises-exercise-from-generic-point-Pn} を用いよ．
\end{exercise}

\begin{exercise}
\label{exercises-exercise-from-generic-point-P1-to-projective}
$k$ を体とする．$X$ を $k$ 上の射影スキームとする．
$K$ を $\mathbf{P}^1_k$ の函数体とする（下のヒントを参照）．
$k$ 上の任意の射
$$
\varphi : \Spec(K) \longrightarrow X
$$
に対して，$k$ 上の射 $\psi : \mathbf{P}^1_k \to X$ が存在し，
$\varphi$ が合成
$$
\Spec(k(t)) \longrightarrow \mathbf{P}^1_k \longrightarrow X
$$
となることを示せ．
ヒント：アファイン開被覆
$\mathbf{P}^1_k = D_+(T_0) \cup D_+(T_1)$ の二つの部分それぞれに対して，
演習 \ref{exercises-exercise-from-generic-point-projective} を用いよ．
$D_+(T_0)$ は多項式環のスペクトルであり，
$K$ はこの多項式環の分数体であることを用いよ．
\end{exercise}

% Frozen source: exercises.tex lines 3683--3762.
\section{曲面から曲線への射}
\label{exercises-section-from-surfaces-to-curves}

\begin{exercise}
\label{exercises-exercise-points-projective-space}
$R$ を環とする．$R \to k$ を $R$ から体への写像とする．
$n \geq 0$ とする．
$$
\Mor_{\Spec(R)}(\Spec(k), \mathbf{P}^n_R)
=
(k^{n + 1} \setminus \{0\})/k^*
$$
を示せ．ここで $k^*$ は $k^{n + 1}$ にスカラー倍で作用する．
以後，元 $(x_0, \ldots, x_n) \in k^{n + 1} \setminus \{0\}$ の同値類に対応する
射 $\Spec(k) \to \mathbf{P}^n_k$ を $(x_0 : \ldots : x_n)$ と表す．
\end{exercise}

\begin{exercise}
\label{exercises-exercise-curve-projective-plane}
$k$ を体とする．$Z \subset \mathbf{P}^2_k$ を
既約かつ被約な閉部分スキームとする．
% P12-ERR-0462: source indefinite-article grammar.
(a) $Z$ は閉点であるか，(b) 次数 $> 0$ の斉次既約多項式
$F \in k[X_0, X_1, X_2]$ が存在して $Z = V_{+}(F)$ となるか，
または (c) $Z = \mathbf{P}^2_k$ となることを示せ．
（ヒント：標準アファイン開集合上で考えよ．）
\end{exercise}

\begin{exercise}
\label{exercises-exercise-bezout}
$k$ を体とする．$Z_1, Z_2 \subset \mathbf{P}^2_k$ を，
% P12-ERR-0463: preserve the frozen unindexed F and indexed F_i.
ある次数 $> 0$ の斉次既約多項式 $F_i \in k[X_0, X_1, X_2]$
に対して $V_{+}(F)$ の形をした既約閉部分スキームとする．
$Z_1 \cap Z_2$ が空でないことを示せ．
（ヒント：次元論を用いて，
$k[X_0, X_1, X_2]/(F_1, F_2)$ の $0$ における局所環の次元を評価せよ．）
\end{exercise}

\begin{exercise}
\label{exercises-exercise-no-nonconstant-morphism-proj}
$\Spec(\mathbf{C})$ 上のスキームの非定数射
$\mathbf{P}^2_{\mathbf{C}} \to \mathbf{P}^1_{\mathbf{C}}$
は存在しないことを示せ．
ここで{\it 定数射}とは，像が一点である射をいう．
（ヒント：射が定数でなければ $0$ と $\infty$ 上のファイバーを考え，
それらが交わらなければならないことを示して矛盾を得よ．）
\end{exercise}

\begin{exercise}
\label{exercises-exercise-nonconstant-morphism}
$k$ を体とする．$X \subset \mathbf{P}^3_k$ を，
一つの斉次方程式 $F \in k[X_0, X_1, X_2, X_3]$ で与えられる閉部分スキームとする．
言い換えると，講義で説明したように
$$
X = \text{Proj}(k[X_0, X_1, X_2, X_3]/(F)) \subset \mathbf{P}^3_k
$$
である．正の次数を持つある斉次多項式 $G, H \in k[X_0, X_1, X_2, X_3]$
に対して
$$
F = X_0 G + X_1 H
$$
となると仮定する．$X_0, X_1, G, H$ が共通零点を持たなければ，
$\Spec(k)$ 上のスキームの非定数射
$$
X \longrightarrow \mathbf{P}^1_k
$$
であって，体に値を持つ点
（演習 \ref{exercises-exercise-points-projective-space} を参照）において，
$x_0$ または $x_1$ が零でないときには
$(x_0 : x_1 : x_2 : x_3) \mapsto (x_0 : x_1)$
という形をとるものが存在することを示せ．
\end{exercise}

% Frozen source: exercises.tex lines 3763--3925.
\section{可逆層}
\label{exercises-section-invertible-sheaves}

\begin{definition}
\label{exercises-definition-invertible-sheaf}
$X$ を局所環付き空間とする．
$X$ 上の{\it 可逆 ${\mathcal O}_X$-加群}とは，
${\mathcal O}_X$-加群の層 ${\mathcal L}$ であって，
各点が開近傍 $U \subset X$ を持ち，
${\mathcal L}|_U$ が ${\mathcal O}_U$-加群として
${\mathcal O}_U$ と同型になるものをいう．
${\mathcal L}$ が ${\mathcal O}_X$-加群として
${\mathcal O}_X$ と同型であるとき，自明であるという．
\end{definition}

\begin{exercise}
\label{exercises-exercise-general-facts-invertible}
一般的な事実．
\begin{enumerate}
\item スキーム $X$ 上の可逆 ${\mathcal O}_X$-加群は準連接であることを示せ．
% P12-ERR-0464: source predicate omits "is".
\item $X\to Y$ を局所環付き空間の射とし，
${\mathcal L}$ を可逆 ${\mathcal O}_Y$-加群とする．
% P12-ERR-0465: preserve the source's unhyphenated module phrase.
$f^\ast {\mathcal L}$ が可逆 ${\mathcal O}_X$ 加群であることを示せ．
\end{enumerate}
\end{exercise}

\begin{exercise}
\label{exercises-exercise-invertible-algebra}
代数．
\begin{enumerate}
\item アファインスキーム $\Spec(A)$ 上の可逆 ${\mathcal O}_X$-加群は，
(i) 有限，(ii) 射影的，(iii) 階数 1 の局所自由であり，
したがって (iv) 平坦，(v) 有限表示であるような
$A$-加群 $M$ に対応することを示せ．
% P12-ERR-0466: source possessive grammar.
（前学期の講義や代数の本の結果は自由に引用してよい．）
% P12-ERR-0467: preserve the frozen manual item reference (a).
\item $A$ を整域とし，$M$ を (a) のような加群とする．
$M$ は $A$-加群として，任意の素イデアル ${\mathfrak p}$ に対して
$IA_{\mathfrak p}$ が単項となるようなイデアル $I \subset A$ と同型であることを示せ．
\end{enumerate}
\end{exercise}

\begin{definition}
\label{exercises-definition-invertible-module}
$R$ を環とする．{\it 可逆加群 $M$} とは，$R$-加群 $M$ であって，
$\widetilde M$ が $R$ のスペクトル上の可逆層であるものをいう．
$R$-加群として $M \cong R$ であるとき，$M$ は{\it 自明}であるという．
\end{definition}

\noindent
言い換えると，$M$ が可逆であることは，
平坦，有限表示，射影的，階数 1 の局所自由という条件をすべて満たすことと同値である．
（もちろん，階数 1 の局所自由であれば十分である．）

\begin{exercise}
\label{exercises-exercise-simple-examples-invertible}
簡単な例．
\begin{enumerate}
\item
\label{exercises-item-affine-line}
$k$ を体とする．$A = k[x]$ とする．
$X = \Spec(A)$ 上の可逆 ${\mathcal O}_X$-加群はすべて自明であることを示せ．
言い換えると，任意の可逆 $A$-加群は階数 1 の自由加群であることを示せ．
\item
\label{exercises-item-affine-line-with-0-and-1-identified}
$A$ を環
$$
A = \{ f\in k[x] \mid f(0) = f(1) \}.
$$
とする．$k = {\mathbf F}_2$ でない限り，
非自明な可逆 $A$-加群が存在することを示せ．
（ヒント：$\Spec(A)$ は ${\mathbf A}^1_k = \Spec(k[x])$
の $0$ と $1$ を同一視したものだと考えよ．）
\item
\label{exercises-item-affine-line-with-cusp}
環 $A = k[x^2, x^3] \subset k[x]$ について
(\ref{exercises-item-affine-line-with-0-and-1-identified}) と同じ問いを考えよ
（ただし今度は $k = {\mathbf F}_2$ の場合にも成り立つ）．
\end{enumerate}
\end{exercise}

\begin{exercise}
\label{exercises-exercise-higher-dimension-invertible}
高次元の場合．
\begin{enumerate}
\item 二次元アファイン空間上の任意の可逆層が自明であることを証明せよ．
より正確には，$k$ を体とし，${\mathbf A}^2_k = \Spec(k[x, y])$ とする．
${\mathbf A}^2_k$ 上の任意の可逆層が自明であることを示せ．
（ヒント：一つの方法は，対応する加群 $M$ を考え，
$M \otimes_{k[x, y]} k(x)[y]$ を調べ，
演習 \ref{exercises-exercise-simple-examples-invertible} (\ref{exercises-item-affine-line})
によってその生成元を見つけることである．その後も考える必要がある．
% P12-ERR-0468: source duplicated infinitive marker.
別の方法は，演習 \ref{exercises-exercise-invertible-algebra} と
多項式環のイデアルについて知っていること，
すなわち高さ一の素イデアルは既約多項式で生成されることを用いる方法である．
その後も考える必要がある．）
\item 二次元アファイン空間の任意の開部分スキーム上の可逆層が
すべて自明であることを証明せよ．
より正確には，$k$ を体とし，$U \subset {\mathbf A}^2_k$ を開部分スキームとする．
$U$ 上の任意の可逆層が自明であることを示せ．
ヒント：$U$ 上の任意の可逆層が ${\mathbf A}^2_k$ 上の可逆層に延長できることを示せ．
易しくはないが，\cite{H} に見つけることができる．
\item 体上の頂点を除いた錐の上で，非自明な可逆層の例を見つけよ．
より正確には，$k$ を体とし，$C = \Spec(k[x, y, z]/(xy-z^2))$ とする．
$U = C \setminus \{ (x, y, z) \}$ とする．
$U$ 上の非自明な可逆層を見つけよ．
ヒント：一つ見つけるだけよりも，$U$ 上の可逆層の同型類の群を計算する方が易しいかもしれない．
$U$ は，「扱いやすい」開集合
$\Spec(k[x, y, z, 1/x]/(xy-z^2))$ と
$\Spec(k[x, y, z, 1/y]/(xy-z^2))$
で被覆されることに注意せよ．
\end{enumerate}
\end{exercise}

\begin{definition}
\label{exercises-definition-picard-group}
$X$ を局所環付き空間とする．
{\it $X$ の Picard 群}とは，可逆 $\mathcal{O}_X$-加群の同型類の集合
$\Pic(X)$ に，テンソル積で与えられる加法を入れたものである．
加群，定義 \ref{modules-definition-pic} を参照．
環 $R$ に対しては $\Pic(R) = \Pic(\Spec(R))$ とおく．
\end{definition}

\begin{exercise}
\label{exercises-exercise-traverso}
$R$ を環とする．
\begin{enumerate}
\item $R$ が Noether 正規整域ならば
$\Pic(R) = \Pic(R[t])$ であることを示せ．
[ヒント：写像 $R \to R[t]$ の左逆となる写像
$R[t] \to R$，$t \mapsto 0$ が存在する．
したがって，$M/tM \cong R$ を満たす任意の可逆
$R[t]$-加群 $M$ が階数 $1$ の自由加群であることを示せば十分である．
$K$ を $R$ の分数体とする．
自明化 $K[t] \to M \otimes_{R[t]} K[t]$ を選ぶ．
これは演習 \ref{exercises-exercise-simple-examples-invertible} (\ref{exercises-item-affine-line})
により可能である．
これを調整して，上の $M/tM$ の自明化と一致するようにせよ．
実際にこれが $R[t]$ 上での $M$ の自明化であることを示せ
（ここで正規性が使われる）．]
\item $k$ を体とする．
$\Pic(k[x^2, x^3, t]) \not = \Pic(k[x^2, x^3])$ を示せ．
\end{enumerate}
\end{exercise}

\section{{\v C}ech コホモロジー}
\label{exercises-section-cech-cohomology}

\begin{exercise}
\label{exercises-exercise-cech-cohomology}
{\v C}ech コホモロジー．ここで $k$ は体とする．
\begin{enumerate}
\item $X$ を開被覆
${\mathcal U} : X = U_1 \cup U_2$ をもつスキームとする．ここで
$U_1 = \Spec(k[x])$，$U_2 =  \Spec(k[y])$，
$U_1 \cap U_2 = \Spec(k[z, 1/z])$ であり，開埋め込み
$U_1 \cap U_2 \to U_1$ および $U_1 \cap U_2 \to U_2$ はそれぞれ
$x \mapsto z$ および $y \mapsto z$ によって定まる
（ここは文字どおりの意味である）．
（講義でこのような $X$ が存在することを見た．これは原点を二重化した
アファイン直線である．）${\check H}^1({\mathcal U}, {\mathcal
O})$ を計算せよ．
たとえば，$k$-ベクトル空間としての基底を与えよ．
\item
${\check H}^1({\mathcal U}, {\mathcal O})$
の各元に対して，${\mathcal O}_X$-加群の層の完全列
$$
0 \to {\mathcal O}_X \to E \to {\mathcal O}_X \to 0
$$
であって，境界 $\delta(1) \in {\check H}^1({\mathcal U},
{\mathcal O})$ が与えられた元に等しくなるものを構成せよ．
（この主張を意味のあるものにすることも問題の一部である．
下も参照せよ．
このようなものが存在せざるを得ないことを抽象的に示してもよい．）
\end{enumerate}
\end{exercise}

\begin{definition}
\label{exercises-definition-delta}
（delta の定義．）任意の位相空間 $X$ 上のアーベル群の層の短完全列
$$
0 \to {\mathcal F}_1 \to {\mathcal F}_2 \to {\mathcal F}_3 \to 0
$$
が与えられたとする．境界写像
$\delta : H^0(X, {\mathcal F}_3) \to {\check H}^1(X, {\mathcal F}_1)$
を次のように定義する．元 $\tau \in H^0(X, {\mathcal F}_3)$ を取る．
各 $i$ に対して，$\tau$ の $U_i$ への制限を持ち上げる切断
$\tilde \tau_i \in {\mathcal F}_2$ が存在するような開被覆
${\mathcal U} : X = \bigcup_{i\in I} U_i$ を選ぶ．
このとき，対応
$$
(i_0, i_1) \longmapsto
\tilde \tau_{i_0}|_{U_{i_0i_1}} - \tilde \tau_{i_1}|_{U_{i_0i_1}}.
$$
を考える．これは明らかに {\v C}ech 複体
${\check C}^\ast({\mathcal U}, {\mathcal F}_2)$ の 1-コバウンダリである．
しかし（${\mathcal F}_1$ を ${\mathcal F}_2$ の部分層とみなすと），
右辺は常に $U_{i_0i_1}$ 上の ${\mathcal F}_1$ の切断であることに注意する．
% P12-ERR-0469: preserve the frozen F_2 coefficient and cochain wording.
したがって，この対応は複体
${\check C}^\ast({\mathcal U}, {\mathcal F}_2)$ の 1-コチェインを定める．
この 1-コチェインのコホモロジー類を，定義により
{\it $\delta(\tau)$} とする．
\end{definition}

\section{コホモロジー}
\label{exercises-section-cohomology}

\begin{exercise}
\label{exercises-exercise-cohomology-not-zero}
$X = \mathbf{R}$ に通常のユークリッド位相を入れる．
コホモロジーの形式的性質（関手性と長完全コホモロジー列）だけを用いて，
$H^1(X, \mathcal{F})$ が零でないような $X$ 上の層
$\mathcal{F}$ が存在することを示せ．
\end{exercise}

\begin{exercise}
\label{exercises-exercise-mayer-vietoris}
位相空間 $X = U \cup V$ が二つの開集合の和として書かれているとする．
このとき，長完全 Mayer--Vietoris 列
$$
0 \to
H^0(X, \mathcal{F}) \to
H^0(U, \mathcal{F}) \oplus H^0(V, \mathcal{F}) \to
H^0(U \cap V, \mathcal{F}) \to
H^1(X, \mathcal{F}) \to \ldots
$$
がある．Mayer--Vietoris 長完全列の構成に本質的な入射的層の性質は何か．
なぜその性質が成り立つのか．
\end{exercise}

\begin{exercise}
\label{exercises-exercise-cohomology-two-point-space}
$X$ を位相空間とする．
\begin{enumerate}
\item $X$ の点が $2$ 個以下ならば，$i > 0$ に対して
$H^i(X, \mathcal{F})$ が零であることを示せ．
\item $X$ の点が $3$ 個ならばどうか．
\end{enumerate}
\end{exercise}

\begin{exercise}
\label{exercises-exercise-cohomology-spec-local-ring}
$X$ を局所環のスペクトルとする．任意のアーベル群の層
$\mathcal{F}$ と $i > 0$ に対して
$H^i(X, \mathcal{F})$ が零であることを示せ．
\end{exercise}

\begin{exercise}
\label{exercises-exercise-affine-morphism}
$f : X \to Y$ をスキームのアファイン射とする．
任意の準連接 $\mathcal{O}_X$-加群 $\mathcal{F}$ に対して
$H^i(X, \mathcal{F}) = H^i(Y, f_*\mathcal{F})$
であることを証明せよ．
必要ならば $X$ と $Y$ にさらにいくつかの条件を課し，
分離スキームのアファイン開被覆と準連接層について，
{\v C}ech コホモロジーとコホモロジーが一致することを用いてよい．
\end{exercise}

\begin{exercise}
\label{exercises-exercise-affine-morphism-to-Pn}
$A$ を環とする．
$\mathbf{P}^n_A = \text{Proj}(A[T_0, \ldots, T_n])$
を $A$ 上の射影空間とする．
$\mathbf{A}^{n + 1}_A = \Spec(A[T_0, \ldots, T_n])$ とし，
$$
U = \bigcup\nolimits_{i = 0, \ldots, n} D(T_i) \subset \mathbf{A}^{n + 1}_A
$$
を閉埋め込み
$0 : \Spec(A) \to \mathbf{A}^{n + 1}_A$
の像の補集合とする．アファイン全射
$$
f : U \longrightarrow \mathbf{P}^n_A
$$
を構成し，
$f_*\mathcal{O}_U =
\bigoplus_{d \in \mathbf{Z}} \mathcal{O}_{\mathbf{P}^n_A}(d)$
であることを証明せよ．
より一般に，次数付き $A[T_0, \ldots, T_n]$-加群 $M$ に対して
$$
f_*(\widetilde{M}|_U) =
\bigoplus\nolimits_{d \in \mathbf{Z}} \widetilde{M(d)}
$$
であることを示せ．ここで左辺では，$\mathbf{A}^{n + 1}_A$ 上で
$M$ に付随する準連接層 $\widetilde{M}$ を用い，
右辺では次数付き加群 $M(d)$ に付随する準連接層
$\widetilde{M(d)}$ を用いる．
\end{exercise}

\begin{exercise}
\label{exercises-exercise-compute-Pn}
$A$ を環とし，
$\mathbf{P}^n_A = \text{Proj}(A[T_0, \ldots, T_n])$
を $A$ 上の射影空間とする．
$\mathbf{P}^n_A$ 上の構造層の Serre 捻り
$\mathcal{O}_{\mathbf{P}^n_A}(d)$ のコホモロジーを丁寧に計算せよ．
分離スキームのアファイン開被覆と準連接層について，
{\v C}ech コホモロジーとコホモロジーが一致すること，および
{\v C}ech コホモロジーを用いてよい．
\end{exercise}

\begin{exercise}
\label{exercises-exercise-cohomology-hypersurface}
$A$ を環とし，
$\mathbf{P}^n_A = \text{Proj}(A[T_0, \ldots, T_n])$
を $A$ 上の射影空間とする．
$F \in A[T_0, \ldots, T_n]$ を次数 $d$ の斉次式とする．
$X \subset \mathbf{P}^n_A$ を
$A[T_0, \ldots, T_n]$ の次数付きイデアル $(F)$
に対応する閉部分スキームとする．
$H^i(X, \mathcal{O}_X)$ について何がいえるか．
\end{exercise}

\begin{exercise}
\label{exercises-exercise-characterize-finite-product-fields}
$R$ を，任意の左完全関手
$F : \text{Mod}_R \to \textit{Ab}$ に対して
$i > 0$ ならば $R^iF = 0$ となる環とする．
$R$ が有限個の体の直積であることを示せ．
\end{exercise}

\section{さらなるコホモロジー}
\label{exercises-section-more-cohomology}

\begin{exercise}
\label{exercises-exercise-cohomology-coordinate-axes}
$k$ を体とする．$X \subset \mathbf{P}^n_k$ を
「座標十字」とする．すなわち，$X$ は斉次方程式
$$
T_i T_j = 0\text{，ただし }i > j > 0
$$
で定義されるものとする．ここで通常どおり
$\mathbf{P}^n_k = \text{Proj}(k[T_0, \ldots, T_n])$ と書く．
言い換えると，$X$ は多項式環の商
$k[T_0, \ldots, T_n]/(T_iT_j; i > j > 0)$
に対応する閉部分スキームである．
すべての $i$ に対して $H^i(X, \mathcal{O}_X)$ を計算せよ．
ヒント：{\v C}ech コホモロジーを用いよ．
\end{exercise}

\begin{exercise}
\label{exercises-exercise-compute-cohomology-punctured}
$A$ を環とする．$I = (f_1, \ldots, f_t)$ を
$A$ の有限生成イデアルとする．
$U \subset \Spec(A)$ を $V(I)$ の補集合とする．
任意の $A$-加群 $M$ に対して，
そのコホモロジー群が $H^n(U, \widetilde{M})$ を与えるような
$A$-加群の複体を
（$A$，$f_1, \ldots, f_t$，$M$ を用いて）書き下せ．
\end{exercise}

\begin{exercise}
\label{exercises-exercise-compute-cohomology-affine-space-punctured}
$k$ を体とする．$U \subset \mathbf{A}^d_k$ を
$\mathbf{A}^d_k$ の閉点 $0$ の補集合とする．
すべての $n$ に対して $H^n(U, \mathcal{O}_U)$ を計算せよ．
\end{exercise}

\begin{exercise}
\label{exercises-exercise-find-curve-genus-one-hundred}
$k$ を体とする．$k$ 上射影的で次元 $1$ のスキーム $X$ であって，
$H^0(X, \mathcal{O}_X) = k$ かつ
$\dim_k H^1(X, \mathcal{O}_X) = 100$
を満たすものを具体的に見つけよ．
\end{exercise}

\begin{exercise}
\label{exercises-exercise-degree-2-cover}
$f : X \to Y$ を次数 $2$ の有限局所自由射とする．
$X$ と $Y$ は整スキームであり，$2$ は $Y$ の構造層で可逆，
すなわち $2 \in \Gamma(Y, \mathcal{O}_Y)$ は可逆であると仮定する．
$\mathcal{O}_Y$-加群の写像
$$
f^\sharp : \mathcal{O}_Y \longrightarrow f_*\mathcal{O}_X
$$
が左逆をもつこと，すなわち
$\tau \circ f^\sharp = \text{id}$ を満たす
$\mathcal{O}_Y$-加群の写像
$\tau : f_*\mathcal{O}_X \to \mathcal{O}_Y$
が存在することを示せ．
$H^n(Y, \mathcal{O}_Y) \to H^n(X, \mathcal{O}_X)$
が単射であることを導け\footnote{次数 $2$ の整スキーム間の有限局所自由射
$X \to Y$ であって，写像
$H^1(Y, \mathcal{O}_Y) \to H^1(X, \mathcal{O}_X)$
が単射でないものは実際に存在する．}．
\end{exercise}

\begin{exercise}
\label{exercises-exercise-pic}
$X$ をスキーム（または局所環付き空間）とする．
対応 $U \mapsto \mathcal{O}_X(U)^*$ は
$\mathcal{O}_X^*$ と書かれる群の層を定める．
$X$ の Picard 群
（定義 \ref{exercises-definition-picard-group}）が
$H^1(X, \mathcal{O}_X^*)$ に等しい理由を簡潔に説明せよ．
\end{exercise}

\begin{exercise}
\label{exercises-exercise-pic-nontrivial}
$\Pic(X)$ が非自明であるアファインスキーム $X$ の例を与えよ．
演習 \ref{exercises-exercise-pic} を用いて，そのような任意の $X$ について
$H^1(X, -)$ は零関手でないことを導け．
\end{exercise}

\begin{exercise}
\label{exercises-exercise-kill-cohomology-complement}
$A$ を環とする．$I = (f_1, \ldots, f_t)$ を $A$ の有限生成イデアルとする．
$U \subset \Spec(A)$ を $V(I)$ の補集合とする．
準連接 $\mathcal{O}_{\Spec(A)}$-加群 $\mathcal{F}$ と，
$p > 0$ を満たす $\xi \in H^p(U, \mathcal{F})$ が与えられたとき，
すべての $i = 1, \ldots, t$ に対して $f_i^n \xi = 0$ となる
$n > 0$ が存在することを示せ．
ヒント：演習 \ref{exercises-exercise-compute-cohomology-punctured} で得た複体を
用いるのが一つの方法である．
\end{exercise}

\begin{exercise}
\label{exercises-exercise-h0-complement}
$A$ を環とする．$I = (f_1, \ldots, f_t)$ を
$A$ の有限生成イデアルとする．
$U \subset \Spec(A)$ を $V(I)$ の補集合とする．
$M$ を $I$-捩れが零である $A$-加群，すなわち
$0 = \Ker((f_1, \ldots, f_t) : M \to M^{\oplus t})$
を満たすものとする．標準同型
$$
H^0(U, \widetilde{M}) = \colim \Hom_A(I^n, M).
$$
があることを示せ．
注意：これは自明ではない．
\end{exercise}

\begin{exercise}
\label{exercises-exercise-Noetherian}
$A$ を Noether 環とする．$I$ を $A$ のイデアルとする．
$M$ を $A$-加群とする．$M[I^\infty]$ を，次で定義される
$I$-冪捩れ元の集合とする：
$$
M[I^\infty] = \{x \in M \mid
\text{ある }n \geq 1\text{ が存在して }I^nx = 0\}
$$
$M' = M/M[I^\infty]$ とおく．このとき $M'$ の $I$-冪捩れは零である．
$$
\colim \Hom_A(I^n, M) = \colim \Hom_A(I^n, M').
$$
であることを示せ．
% P12-ERR-0470: preserve the duplicated final hint number (3).
注意：これは自明ではない．ヒント：(1) $M$ が有限の場合に帰着することを試みよ，
(2) $M$ が有限で $N = M[I^\infty]$ のとき，
$\Ext^1_A(I^n, N)$ の任意の元は，ある $m > 0$ に対する
$\Ext^1_A(I^{n + m}, N)$ への写像で零になることを示せ，
(3) $\Hom_A(I^n, N)$ について (2) と同じことを示せ，
(3) $M$ が有限の場合に長完全列
$$
0 \to \Hom_A(I^n, M[I^\infty]) \to \Hom_A(I^n, M) \to \Hom_A(I^n, M')
\to \Ext^1_A(I^n, M[I^\infty])
$$
を考え，$I^{n + m}$ に対する列と比較して結論を導け．
\end{exercise}

\section{コホモロジー再訪}
\label{exercises-section-more-more-cohomology}


\begin{exercise}
\label{exercises-exercise-nonkillable}
体 $k$，$k$ 上の曲線 $X$，可逆 $\mathcal{O}_X$-加群
$\mathcal{L}$，およびコホモロジー類
$\xi \in H^1(X, \mathcal{L})$ であって，次の性質をもつものの例を作れ：
スキームの任意の全射有限射 $\pi : Y \to X$ に対して，元 $\xi$ の引き戻しは
$H^1(Y, \pi^*\mathcal{L})$ の零でない元である．
ヒント：閉部分スキーム $Z \subset X$ が $1$ 個より多くの閉点からなり，
短完全列
$0 \to \mathcal{L} \to \mathcal{O}_X \to i_*\mathcal{O}_Z \to 0$
が存在するように $X$，$k$，$\mathcal{L}$ を構成せよ．
% P12-ERR-0471: preserve the frozen pullback-as-verb source locus.
その後，引き戻したときに何が起こるかを調べよ．
\end{exercise}

\begin{exercise}
\label{exercises-exercise-cohomology-cycle-curves}
$k$ を代数閉体とする．
$X$ を射影的な $1$ 次元スキームとする．
$X$ が曲線の閉路を含むと仮定する．すなわち，
$n \geq 2$，互いに異なる $1$ 次元整閉部分スキーム
$C_1, \ldots, C_n$，および互いに異なる閉点
$x_1, \ldots, x_n \in X$ が存在して，$x_n \in C_n \cap C_1$，
かつ $i = 1, \ldots, n - 1$ に対して
$x_i \in C_i \cap C_{i + 1}$ であるとする．
$H^1(X, \mathcal{O}_X)$ が零でないことを証明せよ．
ヒント：$\mathcal{F}$ を写像
% P12-ERR-0472: preserve the frozen unindexed direct sum.
$\mathcal{O}_X \to \bigoplus \mathcal{O}_{C_i}$
の像とし，$\kappa(x_i) = k$ および
$H^0(C_i, \mathcal{O}_{C_i}) = k$ を用いて
$H^1(X, \mathcal{F})$ が零でないことを示せ．
さらに，Grothendieck の定理により $H^2(X, -) = 0$ であることを用いよ．
\end{exercise}

\begin{exercise}
\label{exercises-exercise-surface}
$X$ を代数閉体 $k$ 上の射影曲面とする．
$H^1(Z, \mathcal{O}_Z)$ が零でないような真の閉部分スキーム
$Z \subset X$ が存在することを証明せよ．
ヒント：演習 \ref{exercises-exercise-cohomology-cycle-curves} を用いよ．
\end{exercise}

\begin{exercise}
\label{exercises-exercise-surface-better}
$X$ を代数閉体 $k$ 上の射影曲面とする．
任意の $n \geq 0$ に対して，
$\dim_k H^1(Z, \mathcal{O}_Z) > n$
となる真の閉部分スキーム $Z \subset X$ が存在することを示せ．
演習 \ref{exercises-exercise-surface} および
\ref{exercises-exercise-cohomology-cycle-curves} の議論をどのように修正すればよいかだけを説明し，
細部をすべて与える必要はない．
\end{exercise}

\begin{exercise}
\label{exercises-exercise-surface-h2-nonzero}
$X$ を代数閉体 $k$ 上の射影曲面とする．
$H^2(X, \mathcal{F})$ が零でないような連接 $\mathcal{O}_X$-加群
$\mathcal{F}$ が存在することを証明せよ．
ヒント：演習 \ref{exercises-exercise-surface-better} の結果と，
巧妙に選んだ完全列を用いよ．
\end{exercise}

\begin{exercise}
\label{exercises-exercise-kunneth}
$X$ と $Y$ を体 $k$ 上のスキームとする
（必要ならば $X$ と $Y$ はよいスキーム，たとえば qcqs または
$k$ 上射影的であると仮定してよい）．
$X \times Y = X \times_{\Spec(k)} Y$ とおき，射影を
$p : X \times Y \to X$ および $q : X \times Y \to Y$ とする．
準連接 $\mathcal{O}_X$-加群 $\mathcal{F}$ と
準連接 $\mathcal{O}_Y$-加群 $\mathcal{G}$ に対して
$$
H^n(X \times Y,
p^*\mathcal{F} \otimes_{\mathcal{O}_{X \times Y}} q^*\mathcal{G}) =
\bigoplus\nolimits_{a + b = n}
H^a(X, \mathcal{F}) \otimes_k H^b(Y, \mathcal{G})
$$
を証明せよ．または，次元を取ったときにこれが成り立つことだけを示してもよい．
「簡潔な」解答には追加点を与える．
\end{exercise}

\begin{exercise}
\label{exercises-exercise-Oab}
$k$ を体とする．
% P12-ERR-0473: preserve the frozen malformed first projective-line factor.
$X = \mathbf{P}|^1 \times \mathbf{P}^1$ を，
$k$ 上の射影直線とそれ自身との積とし，射影を
$p : X \to \mathbf{P}^1_k$ および $q : X \to \mathbf{P}^1_k$ とする．
$$
\mathcal{O}(a, b) = p^*\mathcal{O}_{\mathbf{P}^1_k}(a)
\otimes_{\mathcal{O}_X} q^*\mathcal{O}_{\mathbf{P}^1_k}(b)
$$
とおく．すべての $i, a, b$ に対して
$H^i(X, \mathcal{O}(a, b))$ の次元を計算せよ．
ヒント：演習 \ref{exercises-exercise-kunneth} を用いよ．
\end{exercise}

\section{コホモロジーと Hilbert 多項式}
\label{exercises-section-cohomology-hilbert-polynomials}

\begin{situation}
\label{exercises-situation-hilbert-polynomial}
$k$ を体とする．$X = \mathbf{P}^n_k$ を
$n$ 次元射影空間とする．$\mathcal{F}$ を連接
$\mathcal{O}_X$-加群とする．
$$
\chi(X, \mathcal{F}) =
\sum\nolimits_{i = 0}^n (-1)^i \dim_k H^i(X, \mathcal{F})
$$
であったことを想起せよ．
$\mathcal{F}$ の {\it Hilbert 多項式}とは，関数
$$
t \longmapsto \chi(X, \mathcal{F}(t))
$$
のことであった．また，
$\mathcal{F}(t) = \mathcal{F} \otimes_{\mathcal{O}_X} \mathcal{O}_X(t)$
であり，ここで $\mathcal{O}_X(t)$ は，構成，定義
\ref{constructions-definition-twist} における構造層の $t$ 次捻りであった．
多様体，\ref{varieties-subsection-hilbert} 節で，
Hilbert 多項式が $t$ の多項式であることを証明した．
\end{situation}

\begin{exercise}
\label{exercises-exercise-hilbert-pol-easy}
設定 \ref{exercises-situation-hilbert-polynomial} のもとで考える．
\begin{enumerate}
\item $P(t)$ が $\mathcal{F}$ の Hilbert 多項式ならば，
$\mathcal{F}(-13)$ の Hilbert 多項式は何か．
\item $P_i$ が $\mathcal{F}_i$ の Hilbert 多項式ならば，
$\mathcal{F}_1 \oplus \mathcal{F}_2$ の Hilbert 多項式は何か．
\item $P_i$ が $\mathcal{F}_i$ の Hilbert 多項式であり，
$\mathcal{F}$ が全射
$\mathcal{F}_1 \to \mathcal{F}_2$ の核であるとき，
$\mathcal{F}$ の Hilbert 多項式は何か．
\end{enumerate}
\end{exercise}

\begin{exercise}
\label{exercises-exercise-find-given-hilbert-pol-dim-1}
設定 \ref{exercises-situation-hilbert-polynomial} において $n \geq 1$ と仮定する．
Hilbert 多項式が $t - 101$ である連接層を見つけよ．
\end{exercise}

\begin{exercise}
\label{exercises-exercise-find-given-hilbert-pol-dim-2}
設定 \ref{exercises-situation-hilbert-polynomial} において $n \geq 2$ と仮定する．
Hilbert 多項式が
$t^2/2 + t/2 - 1$
である連接層を見つけよ．
（これは少し難しい．そのような連接層が存在することだけを示せば十分である．）
\end{exercise}

\begin{exercise}
\label{exercises-exercise-bound-genus-in-degree}
設定 \ref{exercises-situation-hilbert-polynomial} において $n \geq 2$ かつ
$k$ は代数閉体であると仮定する．
$C \subset X$ を次元 $1$ の整閉部分スキームとする．
言い換えると，$C$ は射影曲線である．
$X$ 上の連接層とみなした $\mathcal{O}_C$ の Hilbert 多項式を
$d t + e$ とする．
\begin{enumerate}
\item $e$ の上界を与えよ．
（ヒント：$\mathcal{O}_C(t)$ は次数 $0$ と $1$ にしか
コホモロジーをもたないことを用い，$H^0(C, \mathcal{O}_C)$ を調べよ．）
\item $C$ と横断的に交わる $\mathcal{O}_X(1)$ の大域切断 $s$，
すなわち，互いに異なる閉点 $c_1, \ldots, c_r \in C$ と
短完全列
$$
0 \to \mathcal{O}_C \xrightarrow{s} \mathcal{O}_C(1) \to
\bigoplus\nolimits_{i = 1, \ldots, r} k_{c_i} \to 0
$$
が存在するようなものを選べ．ここで $k_{c_i}$ は
$c_i$ における値が $k$ である摩天楼層である．
（このような $s$ は存在する．ここではそのことを用いてよい．）
$r = d$ であることを示せ．
（ヒント：この列を捻って，何が得られるか調べよ．）
\item 短完全列を捻ると，任意の $t$ に対して $k$-線形写像
$\varphi_t : \Gamma(C, \mathcal{O}_C(t)) \to \bigoplus_{i = 1, \ldots, d} k$
が得られる．
% P12-ERR-0474: preserve the frozen extraneous-if source locus.
$t \geq d - 1$ ならばこの写像が全射であることを示せ．
\item $d$ を用いて $e$ の下界を与えよ．
（ヒント：(3) の結果を用いて
$H^1(C, \mathcal{O}_C(d - 2)) = 0$ であることを示し，消滅を用いよ．）
\end{enumerate}
\end{exercise}

\begin{exercise}
\label{exercises-exercise-three-quadrics-in-plane}
設定 \ref{exercises-situation-hilbert-polynomial} において $n = 2$ と仮定する．
$s_1, s_2, s_3 \in \Gamma(X, \mathcal{O}_X(2))$
を三つの二次方程式とする．連接層
$$
\mathcal{F} = \Coker\left(\mathcal{O}_X(-2)^{\oplus 3}
\xrightarrow{s_1, s_2, s_3} \mathcal{O}_X\right)
$$
を考える．このような $\mathcal{F}$ の取り得る Hilbert 多項式を列挙せよ．
（射影平面内の二次曲線の交わりを視覚化してみよ．）
\end{exercise}

\section{曲線}
\label{exercises-section-curves}

\begin{exercise}
\label{exercises-exercise-closed-subset-vanshing}
$k$ を代数閉体とする．$X$ を $k$ 上の射影曲線とする．
$\mathcal{L}$ を可逆 $\mathcal{O}_X$-加群とする．
$s_0, \ldots, s_n \in H^0(X, \mathcal{L})$
を $\mathcal{L}$ の大域切断とする．
閉点 $((\lambda_0 : \ldots : \lambda_n), x)$ が $Z$ に属することと，
切断 $\lambda_0 s_0 + \ldots + \lambda_n s_n$ が $x$ で消えることとが
同値になるような，自然な閉部分スキーム
$$
Z \subset
\mathbf{P}^n \times X
$$
が存在することを証明せよ．
（ヒント：$Z$ をアファイン局所的に記述せよ．）
\end{exercise}

\begin{exercise}
\label{exercises-exercise-diagonal-cartier}
$k$ を代数閉体とする．$X$ を $k$ 上の滑らかな曲線とする．
$r \geq 1$ とする．閉点が，ある $i$ に対して
$x = x_i$ を満たす組 $(x, x_1, \ldots, x_r)$ であるような閉部分集合
$$
D \subset X \times X^r = X^{r + 1}
$$
が可逆なイデアル層をもつことを示せ．
言い換えると，$D$ が有効 Cartier 因子であることを示せ．
ヒント：$r = 1$ に帰着し，
% P12-ERR-0475: preserve the frozen singular/plural source locus.
$X$ が滑らかな曲線であることを用いて対角線について何かをいえ
（この点については書籍を参照せよ）．
\end{exercise}

\begin{exercise}
\label{exercises-exercise-one-divisor-in-another}
$k$ を代数閉体とする．$X$ を $k$ 上の滑らかな射影曲線とする．
$T$ を $k$ 上有限型のスキームとし，
$$
D_1 \subset X \times T
\quad\text{かつ}\quad
D_2 \subset X \times T
$$
を二つの有効 Cartier 因子とする．
$t \in T$ に対して，ファイバー $D_{i, t} \subset X_t$ は稠密でない
（すなわち，曲線 $X_t$ の一般点を含まない）と仮定する．
閉点 $t \in T$ が $Z$ に属することと，$D_1$，$D_2$ の
スキーム論的ファイバー $D_{1, t}$，$D_{2, t}$ が
$X_t$ の閉部分スキームとして
$$
D_{1, t} \subset D_{2, t}
$$
を満たすこととが同値になるような，標準的な閉部分スキーム
$Z \subset T$ が存在することを証明せよ．
ヒント：必要なら $T$ を縮小して，$T = \Spec(A)$ はアファインであり，
$D_i \subset U \times T$ となるアファイン開集合 $U \subset X$ が
存在すると仮定してよいことを示せ．次に
$M_1 = \Gamma(D_1, \mathcal{O}_{D_1})$ が有限局所自由
$A$-加群であることを示せ
（ここでは非自明な代数が必要になるので，友人に尋ねよ）．
$T$ を縮小して，$M_1$ は自由 $A$-加群であると仮定してよい．
そこで
$$
\Gamma(U \times T, \mathcal{I}_{D_2}) \to M_1 = A^{\oplus N}
$$
を調べ，この写像の像に含まれるベクトルの係数が生成するイデアルで
切り出される閉部分スキームとして $Z$ を定義せよ．
これでうまくいく理由を説明せよ
（おそらく，もう少し可換環論が必要になる）．
\end{exercise}

\begin{exercise}
\label{exercises-exercise-closed-subset-vanshing-r}
$k$ を代数閉体とする．$X$ を $k$ 上の滑らかな射影曲線とする．
$\mathcal{L}$ を可逆 $\mathcal{O}_X$-加群とする．
$s_0, \ldots, s_n \in H^0(X, \mathcal{L})$
を $\mathcal{L}$ の大域切断とする．$r \geq 1$ とする．
閉点
$((\lambda_0 : \ldots : \lambda_n), x_1, \ldots, x_r)$
が $Z$ に属することと，切断
$s_\lambda = \lambda_0 s_0 + \ldots + \lambda_n s_n$ が
因子 $D = x_1 + \ldots + x_r$ 上で消えること，
すなわち切断 $s_\lambda$ が $\mathcal{L}(-D)$ に属することとが
同値になるような，自然な閉部分スキーム
$$
Z \subset
\mathbf{P}^n \times X \times \ldots \times X = \mathbf{P}^n \times X^r
$$
が存在することを証明せよ．
% P12-ERR-0476: preserve the frozen then/the typo locus.
ヒント：演習 \ref{exercises-exercise-diagonal-cartier} と
\ref{exercises-exercise-one-divisor-in-another} の結果を組み合わせれば
従うことを説明せよ．
\end{exercise}

\begin{exercise}
\label{exercises-exercise-effective}
$k$ を代数閉体とする．$X$ を $k$ 上の滑らかな射影曲線とする．
$\mathcal{L}$ を可逆 $\mathcal{O}_X$-加群とする．
$X^r$ の閉点 $(x_1, \ldots, x_r)$ が $Z$ に属することと，
$\mathcal{L}(-x_1 - \ldots -x_r)$ が零でない大域切断をもつこととが
同値になるような，自然な閉部分集合
$$
Z \subset X^r
$$
が存在することを示せ．
ヒント：演習 \ref{exercises-exercise-closed-subset-vanshing-r} を用いよ．
\end{exercise}

\begin{exercise}
\label{exercises-exercise-difference-effective}
$k$ を代数閉体とする．$X$ を $k$ 上の滑らかな射影曲線とする．
$r \geq s$ を整数とする．
$X^r \times X^s$ の閉点
$(x_1, \ldots, x_r, y_1, \ldots, y_s)$ が $Z$ に属することと，
$x_1 + \ldots + x_r - y_1 - \ldots - y_s$ が有効因子に線形同値であることとが
同値になるような，自然な閉部分集合
$$
Z \subset X^r \times X^s
$$
が存在することを示せ．
% P12-ERR-0477: preserve the frozen nonvanshing typo locus.
ヒント：次数 $r$ の任意の有効因子 $D$ に対して
$\mathcal{L}(-D)$ が零でない切断をもつほど次数の高い補助的な
可逆加群 $\mathcal{L}$ を選べ．
次に演習 \ref{exercises-exercise-effective} の結果を二度用いよ．
\end{exercise}

\begin{exercise}
\label{exercises-exercise-genus-7}
好みの代数閉体 $k$ を選べ．
種数 $7$ のある曲線上に存在し得るすべての $\mathfrak g^r_d$ を，
できる限り決定せよ．その際，
\begin{enumerate}
\item どの場合に $\mathfrak g^r_d$ が不動点をもたないか，および
\item どの場合に $\mathfrak g^r_d$ が $\mathbf{P}^r$ への閉埋め込みを
与えるか
\end{enumerate}
も決定せよ．曲線が「一般」であると仮定した場合にも同じ問題を解け
（一般という概念は自分で定めてよい．こちらは上の問題より易しいかもしれない）．
曲線が超楕円的であると仮定した場合にも同じ問題を解け．
曲線が trigonal（かつ超楕円的でない）であると仮定した場合にも同じ問題を解け．
以下同様である．
\end{exercise}

\section{モジュライ}
\label{exercises-section-moduli}

\noindent
この節では，代数幾何学的対象のモジュライに対する素朴な取り扱いをいくつか考える．

\medskip\noindent
$k$ を代数閉体とする．$M$ を $k$ 上のモジュライ問題とする．
これが正確に何を意味するかはここでは定義しないが，
各演習で意図は明らかであろう．
以下を理解するには，$k$ 上の $M$ の対象，$k$ 上の
$M$ の対象間の同型，および多様体上の
% P12-ERR-0478: preserve the frozen singular-object source locus.
$M$ の対象の族が何であるかを知れば十分である．
このとき，次の条件が成り立つ場合，{\it $M$ のモジュライ数}は
$d \geq 0$ であるという：
\begin{enumerate}
\item 有限個の族 $X_i \to V_i$，$i = 1, \ldots, n$ が存在し，
$k$ 上の $M$ の任意の対象は，それらのいずれかのファイバーと同型であり，
かつ $\max \dim(V_i) = d$ である．
\item これをより小さい $d$ で行うことはできない．
\end{enumerate}
これは実際には，文献にあるよりよい概念の非常に粗い近似にすぎない．

\begin{exercise}
\label{exercises-exercise-number-moduli-lines-conics}
$k$ を代数閉体とする．$d \geq 1$ および $n \geq 1$ とする．
$P^n$ 内の次数 $d$ の超曲面のモジュライを次のように定める：
\begin{enumerate}
\item 対象は次数 $d$ の超曲面
$X \subset \mathbf{P}^n_k$ である．
\item 二つの対象 $X \subset \mathbf{P}^n_k$ と
$Y \subset \mathbf{P}^n_k$ の間の同型は，
% P12-ERR-0479: preserve the frozen PGL_n formula.
$g(X) = Y$ を満たす元 $g \in \text{PGL}_n(k)$ である．
\item 多様体 $V$ 上の超曲面の族は閉部分スキーム
$X \subset \mathbf{P}^n_V$ であって，任意の $v \in V$ に対して，
% P12-ERR-0480: preserve the frozen singular/plural source locus.
$X \to V$ のスキーム論的ファイバー $X_v$ が
$\mathbf{P}^n_v$ 内の超曲面であるものである．
\end{enumerate}
次のモジュライ数を計算せよ
（可能ならば．後になるほど難しくなる）：
\begin{enumerate}
\item $n = 1$ で $d$ が任意の場合，
\item $n = 2$ かつ $d = 1$ の場合，
\item $n = 2$ かつ $d = 2$ の場合，
\item $n \geq 1$ かつ $d = 2$ の場合，
\item $n = 2$ かつ $d = 3$ の場合，
\item $n = 3$ かつ $d = 3$ の場合，および
\item $n = 2$ かつ $d = 4$ の場合．
\end{enumerate}
\end{exercise}

\begin{exercise}
\label{exercises-exercise-number-moduli-hyperelliptic}
$k$ を代数閉体とする．$g \geq 2$ とする．
種数 $g$ の超楕円曲線のモジュライを次のように定める：
\begin{enumerate}
\item 対象は種数 $g$ の滑らかな射影超楕円曲線 $X$ である．
\item 二つの対象 $X$ と $Y$ の間の同型は，
$k$ 上のスキームの同型 $X \to Y$ である．
\item 多様体 $V$ 上の種数 $g$ の超楕円曲線の族は，
% P12-ERR-0481: preserve the frozen X-to-Y family morphism.
固有平坦\footnote{この仮定を外しても問題の答えは変わらない．}射
$X \to Y$ であって，$X \to V$ のすべてのスキーム論的ファイバーが
種数 $g$ の滑らかな射影超楕円曲線であるものである．
\end{enumerate}
モジュライ数が $2g - 1$ であることを示せ．
\end{exercise}

\section{大域 Ext}
\label{exercises-section-global-exts}

\begin{exercise}
\label{exercises-exercise-ext-line-plane-p3}
$k$ を体とする．$X = \mathbf{P}^3_k$ とする．
$L \subset X$ と $P \subset X$ を，それぞれ $1$ 本および $2$ 本の
一次方程式で切り出される閉部分スキームとみなした直線および平面とする．
すべての $i$ に対して
$$
\Ext^i_X(\mathcal{O}_L, \mathcal{O}_P)
$$
の次元を計算せよ．$L$ が $P$ に含まれる場合と，
$L$ が $P$ に含まれない場合の両方を必ず扱うこと．
\end{exercise}

\begin{exercise}
\label{exercises-exercise-ext-point}
$k$ を体とする．$X = \mathbf{P}^n_k$ とする．
$Z \subset X$ を，閉部分スキームとみなした閉 $k$-有理点とする．
たとえば斉次座標 $(1 : 0 : \ldots : 0)$ をもつ点である．
すべての $i$ に対して
$$
\Ext^i_X(\mathcal{O}_Z, \mathcal{O}_Z)
$$
の次元を計算せよ．
\end{exercise}

\begin{exercise}
\label{exercises-exercise-ext-pairings}
$X$ を環付き空間とする．$\mathcal{O}_X$-加群
$\mathcal{F}, \mathcal{G}, \mathcal{H}$ に対して
カップ積写像
$$
\Ext^i_X(\mathcal{G}, \mathcal{H})
\times
\Ext^j_X(\mathcal{F}, \mathcal{G})
\longrightarrow
\Ext^{i + j}_X(\mathcal{F}, \mathcal{H})
$$
を定義せよ．
（ヒント：これはきわめて一般的な構成である．）
\end{exercise}

\begin{exercise}
\label{exercises-exercise-move-out-locally-free}
$X$ を環付き空間とする．$\mathcal{E}$ を有限局所自由
$\mathcal{O}_X$-加群とし，その双対を
$\mathcal{E}^\vee = \SheafHom_{\mathcal{O}_X}(\mathcal{E}, \mathcal{O}_X)$
とする．次の主張を証明せよ：
\begin{enumerate}
\item $\SheafExt^i_{\mathcal{O}_X}(
\mathcal{F} \otimes_{\mathcal{O}_X} \mathcal{E}, \mathcal{G}) =
\SheafExt^i_{\mathcal{O}_X}(
\mathcal{F},
\mathcal{E}^\vee \otimes_{\mathcal{O}_X} \mathcal{G}) =
\SheafExt^i_{\mathcal{O}_X}(
\mathcal{F}, \mathcal{G}) \otimes_{\mathcal{O}_X} \mathcal{E}^\vee$，および
\item $\Ext^i_X(
\mathcal{F} \otimes_{\mathcal{O}_X} \mathcal{E}, \mathcal{G}) =
\Ext^i_X(\mathcal{F},
\mathcal{E}^\vee \otimes_{\mathcal{O}_X} \mathcal{G})$．
\end{enumerate}
ここで $\mathcal{F}$ と $\mathcal{G}$ は
$\mathcal{O}_X$-加群である．次を導け：
$$
\Ext^i_X(\mathcal{E}, \mathcal{G}) =
H^i(X, \mathcal{E}^\vee \otimes_{\mathcal{O}_X} \mathcal{G})
$$
\end{exercise}

\begin{exercise}
\label{exercises-exercise-trace-on-ext}
$X$ を環付き空間とする．$\mathcal{E}$ を有限局所自由
$\mathcal{O}_X$-加群とする．すべての $i$ に対してトレース写像
$$
\Ext^i_X(\mathcal{E}, \mathcal{E}) \to H^i(X, \mathcal{O}_X)
$$
を構成せよ．任意の $\mathcal{O}_X$-加群 $\mathcal{F}$ に対する
トレース写像
$$
\Ext^i_X(\mathcal{E}, \mathcal{E} \otimes_{\mathcal{O}_X} \mathcal{F})
\to H^i(X, \mathcal{F})
$$
へ一般化せよ．
\end{exercise}

\begin{exercise}
\label{exercises-exercise-duality}
$k$ を体とする．$X = \mathbf{P}^d_k$ とする．
$\omega_{X/k} = \mathcal{O}_X(-d - 1)$ とおく．
有限局所自由加群 $\mathcal{E}$，$\mathcal{F}$ に対して，
Ext 上のカップ積と Ext 上のトレース写像を組み合わせた写像
$$
\Ext^i_X(\mathcal{E}, \mathcal{F} \otimes_{\mathcal{O}_X} \omega_{X/k})
\times
\Ext^{d - i}_X(\mathcal{F}, \mathcal{E})
\to
\Ext_X^d(\mathcal{F}, \mathcal{F} \otimes_{\mathcal{O}_X} \omega_{X/k})
\to
H^d(X, \omega_{X/k}) = k
$$
が非退化対を与えることを証明せよ．
ヒント：この設定で双対性を改めて証明してもよいし，
層のコホモロジーに帰着して，講義で証明した Serre 双対定理を適用してもよい．
\end{exercise}

\section{因子}
\label{exercises-section-divisors}

\noindent
便宜のため，関係するすべての定義をここ一箇所にまとめる．

\begin{definition}
\label{exercises-definition-divisor}
以下，$S$ を任意のスキームとし，$X$ を Noether 整スキームとする．
\begin{enumerate}
\item $X$ 上の {\it Weil 因子}とは，素因子 $Z_i$ の整数係数の
形式的線形結合 $\Sigma n_i[Z_i]$ のことである．
\item {\it 素因子}とは，整であり，一般点 $\xi \in Z$ において
${\mathcal O}_{X, \xi}$ の次元が $1$ となる閉部分スキーム
$Z \subset X$ のことである．上のような $\xi \in Z \subset X$ に対して，
記号 ${\mathcal O}_{X, Z} = {\mathcal O}_{X, \xi}$ を用いる．
${\mathcal O}_{X, Z} \subset K(X)$ は $X$ の関数体の部分環であることに注意せよ．
\item {\it 有理関数 $f \in K(X)^\ast$ に付随する Weil 因子}とは，
和 $\Sigma v_Z(f)[Z]$ のことである．ここで $v_Z(f)$ を次のように定義する：
\begin{enumerate}
\item $f \in {\mathcal O}_{X, Z}^\ast$ ならば $v_Z(f) = 0$ とする．
\item $f \in {\mathcal O}_{X, Z}$ ならば
$$
v_Z(f) = \text{length}_{{\mathcal O}_{X, Z}}({\mathcal O}_{X, Z}/(f)).
$$
とする．
\item $a, b \in {\mathcal O}_{X, Z}$ を用いて $f = \frac{a}{b}$ と書けるならば
$$
v_Z(f) = \text{length}_{{\mathcal O}_{X, Z}}({\mathcal O}_{X, Z}/(a)) -
\text{length}_{{\mathcal O}_{X, Z}}({\mathcal O}_{X, Z}/(b)).
$$
とする．
\end{enumerate}
\item スキーム $S$ 上の {\it 有効 Cartier 因子}とは，任意の点
$d\in D$ が $S$ 内にアファイン開近傍
% P12-ERR-0482: preserve the frozen duplicated ambient-space locus.
$\Spec(A) = U \subset S$ をもち，そこで
$D \cap U = \Spec(A/(f))$ が非零因子 $f \in A$ によって定まるような
閉部分スキーム $D \subset S$ のことである．
\item Noether 整スキーム $X$ の有効 Cartier 因子 $D \subset X$ に
{\it 付随する Weil 因子 $[D]$}とは，和 $\Sigma v_Z(D)[Z]$ のことである．
ここで $v_Z(D)$ を次のように定義する：
\begin{enumerate}
\item $Z$ の一般点 $\xi$ が $D$ に属さないならば
$v_Z(D) = 0$ とする．
\item $Z$ の一般点 $\xi$ が $D$ に属するならば
$$
v_Z(D) = \text{length}_{{\mathcal O}_{X, Z}}({\mathcal O}_{X, Z}/(f))
$$
とする．ここで $f \in {\mathcal O}_{X, Z} = {\mathcal O}_{X, \xi}$ は，
$\xi$ のアファイン近傍で $D$ を定義する非零因子
（上の (4) の意味でのもの）である．
\end{enumerate}
\item $S$ をスキームとする．{\it 全商環の層 ${\mathcal K}_S$}とは，
次のように定義される前層 ${\mathcal K}'$ を層化して得られる
${\mathcal O}_S$-代数の層のことである．
開集合 $U \subset S$ に対して
${\mathcal K}'(U) = S_U^{-1}{\mathcal O}_S(U)$ とおく．
ここで $S_U \subset {\mathcal O}_S(U)$ は，切断
$f \in {\mathcal O}_S(U)$ であって，任意の $u\in U$ に対して
${\mathcal O}_{S, u}$ における $f$ の芽が非零因子であるものからなる
乗法的部分集合である．
特に，$S_U$ の元はすべて非零因子である．
したがって ${\mathcal O}_S$ は ${\mathcal K}_S$ の部分層であり，
短完全列
$$
0 \to {\mathcal O}_S^\ast \to {\mathcal K}_S^\ast \to
{\mathcal K}_S^\ast/{\mathcal O}_S^\ast \to 0.
$$
を得る．
\item スキーム $S$ 上の {\it Cartier 因子}とは，商層
${\mathcal K}_S^\ast/{\mathcal O}_S^\ast$ の大域切断のことである．
\item Noether 整スキーム $X$ 上の Cartier 因子
$\tau \in \Gamma(X, {\mathcal K}_X^\ast/{\mathcal O}_X^\ast)$ に
{\it 付随する Weil 因子}とは，和 $\Sigma v_Z(\tau)[Z]$ のことである．
% P12-ERR-0483: preserve the frozen extraneous-as source locus.
ここで $v_Z(\tau)$ は次の手順で定義される：
\begin{enumerate}
\item $Z$ の一般点 $\xi$ における $\tau$ の芽が零である，
言い換えると，茎 $({\mathcal K}^\ast/{\mathcal O}^\ast)_\xi$ における
$\tau$ の像が「零」であるならば，$v_Z(\tau) = 0$ とする．
\item $\tau|_U$ が切断 $f \in {\mathcal K}(U)$ の像であり，
しかも $a, b \in A$ を用いて $f = a/b$ と書けるような
アファイン開近傍 $\Spec(A) = U \subset X$ を見つけよ．このとき
$$
v_Z(f) = \text{length}_{{\mathcal O}_{X, Z}}({\mathcal O}_{X, Z}/(a)) -
\text{length}_{{\mathcal O}_{X, Z}}({\mathcal O}_{X, Z}/(b)).
$$
とおく．
\end{enumerate}
\end{enumerate}
\end{definition}

\begin{remarks}
\label{exercises-remarks-divisors}
自明な注意をいくつか述べる．
\begin{enumerate}
\item Noether 整スキーム $X$ 上では，層 ${\mathcal K}_X$ は
関数体 $K(X)$ を値にもつ定値層である．
\item 上の定義を意味のあるものにするには，1 次元 Noether 局所整域
${\mathcal O}$ の任意の零でない元の対 $(a, b)$ に対して
$$
\text{length}_{\mathcal O}({\mathcal O}/(ab)) =
\text{length}_{\mathcal O}({\mathcal O}/(a)) +
\text{length}_{\mathcal O}({\mathcal O}/(b))
$$
を示す必要がある．これは講義で行う．
\end{enumerate}
\end{remarks}

\begin{exercise}
\label{exercises-exercise-effective-cartier-cartier}
（任意のスキーム上で．）
有効 Cartier 因子に Cartier 因子を対応させる方法を記述せよ．
\end{exercise}

\begin{exercise}
\label{exercises-exercise-rational-function-cartier}
（整スキーム上で．）
有理関数 $f$ に Cartier 因子 $D$ を対応させ，
$D$ と $f$ に付随する Weil 因子が一致するようにする方法を記述せよ．
（これはたわいない．）
\end{exercise}

\begin{exercise}
\label{exercises-exercise-weil-not-cartier}
いかなる Cartier 因子にも付随しない，多様体上の Weil 因子の例を与えよ．
\end{exercise}

\begin{exercise}
\label{exercises-exercise-weil-Q-cartier}
いかなる Cartier 因子にも付随しない多様体上の Weil 因子 $D$ であって，
ある $n > 1$ に対して $nD$ が Cartier 因子に付随する Weil 因子となるものの
例を与えよ．
\end{exercise}

\begin{exercise}
\label{exercises-exercise-weil-not-Q-cartier}
いかなる Cartier 因子にも付随せず，かつ任意の $n > 1$ に対して
$nD$ も Cartier 因子に付随する Weil 因子ではないような，
多様体上の Weil 因子 $D$ の例を与えよ．
（ヒント：錐，たとえば $\mathbf{A}^4_k$ 内の
$X : xy - zw = 0$ を考えよ．
$D = [x = 0, z = 0]$ が条件を満たすことを示してみよ．）
\end{exercise}

\begin{exercise}
\label{exercises-exercise-cartier-not-difference-effective-cartier}
体上有限型の分離スキーム $X$ について：
二つの有効 Cartier 因子の差ではない Cartier 因子の例を与えよ．
ヒント：空でない有効 Cartier 因子を一つももたない $X$，
たとえば \cite[III Exercise 5.9]{H} で構成されるスキームを見つけよ．
$X$ が多様体である例さえ存在する．すなわち，演習
\ref{exercises-exercise-nonprojective} の多様体である．
\end{exercise}

\begin{exercise}
\label{exercises-exercise-nonprojective}
非射影的固有多様体の例．
$k$ を体とする．$L \subset \mathbf{P}^3_k$ を直線，
$C \subset \mathbf{P}^3_k$ を非特異二次曲線とする．
$C \cap L = \emptyset$ と仮定する．
同型 $\varphi : L \to C$ を選ぶ．
$C$ と $L$ を $\varphi$ によって貼り合わせて得られる
$k$-多様体を $X$ とする．言い換えると，全射固有双有理射
$$
\pi : \mathbf{P}^3_k \longrightarrow X
$$
と開集合 $U \subset X$ が存在して，
$\pi : \pi^{-1}(U) \to U$ は同型であり，
$\pi^{-1}(U) = \mathbf{P}^3_k \setminus (L \cup C)$ であり，
$\pi|_L = \pi|_C \circ \varphi$ を満たす．
（これらの条件だけではまだ $X$ を一意に定めない．
一意に定めるには，$Z = X \setminus U$ の点に沿う $X$ の構造層を
指定する必要がある．）
$X$ が存在し，固有多様体であるが射影的でないことを示せ．
（ヒント：存在については演習 \ref{exercises-exercise-prepare-glueing} の結果を用いよ．
非射影性については $\Pic(\mathbf{P}^3_k) = \mathbf{Z}$ を用いて，
$X$ が豊富可逆層をもつことはできないと示せ．）
\end{exercise}

\section{微分}
\label{exercises-section-differentials}

\noindent
{\bf 定義と結果．} K\"ahler 微分．
\begin{enumerate}
\item $R \to A$ を環準同型とする．
{\it $R$ 上の $A$ の K\"ahler 微分加群}を
$\Omega_{A/R}$ と書く．これは元 $\text{d}a$，$a \in A$ で生成され，
次の関係を満たす：
$$
\text{d}(a_1 + a_2) = \text{d}a_1 + \text{d}a_2,\quad
\text{d}(a_1a_2) = a_1\text{d}a_2 + a_2\text{d}a_1,\quad
\text{d}r = 0
$$
標準的な普遍 $R$-導分
$\text{d} : A \to \Omega_{A/R}$ は
$a\mapsto \text{d}a$ と写す．
\item イデアル $I$ を定義する短完全列
$$
0 \to I \to A \otimes_R A \to A \to 0
$$
を考える．$a$ を $a \otimes 1 - 1 \otimes a$ の類に写す標準的な導分
$\text{d} : A \to I/I^2$ がある．
% P12-ERR-0546: preserve the frozen module-of-derivations wording.
これは $R$ 上の $A$ の導分の加群の別の表示である．言い換えると
$$
(I/I^2, \text{d}) \cong (\Omega_{A/R}, \text{d}).
$$
\item $S_R$ の像が $S_A$ に含まれるような乗法的部分集合
$S_R \subset R$ と $S_A \subset A$ に対して
$$
\Omega_{S_A^{-1}A / S_R^{-1}R} =
S_A^{-1}\Omega_{A/R}.
$$
が成り立つ．
\item $A$ が有限表示 $R$-代数ならば，
$\Omega_{A/R}$ は有限表示 $A$-加群である．
したがってこの場合，$\Omega_{A/R}$ の {\it Fitting} イデアルが定義される．
\item $f : X \to S$ をスキームの射とする．
準連接 ${\mathcal O}_X$-加群の層 $\Omega_{X/S}$ と
% P12-ERR-0484: preserve the frozen O_S-linear wording.
${\mathcal O}_S$-線形導分
$$
\text{d} : {\mathcal O}_X \longrightarrow \Omega_{X/S}
$$
が存在して，$f(U) \subset V$ を満たす任意のアファイン開集合
$\Spec(A) = U \subset X$，$\Spec(R) = V \subset S$ に対して
$$
\Gamma(\Spec(A), \Omega_{X/S}) = \Omega_{A/R}
$$
であり，これは $\text{d}$ と両立する．
\end{enumerate}

\begin{exercise}
\label{exercises-exercise-dual-numbers}
$k[\epsilon]$ を体 $k$ 上の双対数の環，すなわち
$\epsilon^2 = 0$ とする．
\begin{enumerate}
\item 環準同型
$$
R = k[\epsilon] \to A = k[x, \epsilon]/(\epsilon x)
$$
を考える．$\Omega_{A/R}$ の Fitting イデアルが
（第零 Fitting イデアルから始めて）
$$
(\epsilon), A, A, \ldots
$$
であることを示せ．
\item 写像
$R = k[t] \to
A = k[x, y, t]/(x(y-t)(y-1), x(x-t))$
を考える．簡単のため $k$ の標数は零と仮定して，
$A$ における $\Omega_{A/R}$ の Fitting イデアルが
$$
x(2x-t)(2y-t-1)A, \ (x, y, t)\cap (x, y-1, t), \ A, \ A, \ldots
$$
であることを示せ．したがって，
% P12-ERR-0485: preserve the frozen 0-the ordinal source locus.
第 $0$ Fitting イデアルは $A$ の一つの元で切り出され，
第 $1$ Fitting イデアルは $\Spec(A)$ の二つの閉点を定義し，
その他はすべて自明である．
\item 写像 $R = k[t] \to A = k[x, y, t]/(xy-t^n)$ を考える．
$\Omega_{A/R}$ の Fitting イデアルを計算せよ．
\end{enumerate}
\end{exercise}

\begin{remark}
\label{exercises-remark-fitting-omega-not-sings}
% P12-ERR-0486: preserve the frozen alphabetic references to numeric items.
$X$ が $S$ 上相対次元 $k$ をもつとき，
$\Omega_{X/S}$ の第 $k$ Fitting イデアルは，
射 $X \to S$ の特異スキームを定義するために普通用いられる．
しかし (a) が示すように，族のファイバー次元が「一定」でない場合，
たとえば平坦でない場合には，これを行う際に注意が必要である．
(b) が示すように，平坦性もこれがうまくいくことを保証しない
（実際，これは平坦な族である）．
(c) のような「よい場合」には，曲線の族について，
第 $0$ Fitting イデアルは零であり，
第 $1$ Fitting イデアルは特異点集合を
（スキーム論的に）定義すると期待される．
\end{remark}

\begin{exercise}
\label{exercises-exercise-formally-smooth}
$R$ を環とし，
$$
A = R[x_1, \ldots, x_n]/(f_1, \ldots, f_n).
$$
とする．$A$ は，変数と同じ個数の方程式で表示されていることに注意せよ．
したがって偏導関数の行列
$$
( \partial f_i / \partial x_j )
$$
は $n \times n$ 行列，すなわち正方行列である．
その行列式が $A$ の元として可逆であると仮定する．
これは，$n$ 個の生成元と $n$ 個の関係式をもつこの場合に
$\Omega_{A/R} = (0)$ であるという条件にほかならない．
$\pi : B' \to B$ を $R$-代数の全射とし，
その核 $J$ は（$B'$ のイデアルとして）平方零であるとする．
$\varphi : A \to B$ を $R$-代数の準同型とする．
$\varphi = \pi \circ \varphi'$ を満たす $R$-代数の準同型
$\varphi' : A \to B'$ が一意に存在することを示せ．
\end{exercise}

\begin{exercise}
\label{exercises-exercise-formally-smooth-one-equation}
$A = R[x, y]/(f)$ の場合へ，演習
\ref{exercises-exercise-formally-smooth} の結果を一般化せよ．
\end{exercise}

\begin{exercise}
\label{exercises-exercise-Jacobian-criterion}
$k$ を体とし，$f_1, \ldots, f_c \in k[x_1, \ldots, x_n]$ とし，
$A = k[x_1, \ldots, x_n]/(f_1, \ldots, f_c)$ とする．
$f_j(0, \ldots, 0) = 0$ と仮定する．これは
$\mathfrak m = (x_1, \ldots, x_n)A$ が極大イデアルであることを意味する．
行列
$$
(\partial f_j/ \partial x_i)|_{(x_1, \ldots, x_n) = (0, \ldots, 0)}
$$
の階数が $c$ ならば，局所環 $A_\mathfrak m$ が正則であることを証明せよ．
この場合の $A_\mathfrak m$ の次元は何か．
$A_\mathfrak m$ は正則であるが階数が $c$ 未満となる例を与えて，
逆が偽であることを示せ．その例における
$A_\mathfrak m$ の次元は何か．
\end{exercise}

\section{スキーム論，2007 年秋学期末試験}
\label{exercises-section-final-exam-fall-2007}

\noindent
% P12-ERR-0487: preserve the frozen Fall-heading/Spring-prose mismatch.
以下は，Columbia 大学で 2007 年春に行われたスキーム論の講義の
期末試験問題である．

\begin{exercise}[定義]
\label{exercises-exercise-definitions}
次の概念を定義せよ．
\begin{enumerate}
\item $X$ は {\it スキーム}である．
\item スキームの射 $f : X \to Y$ は {\it 有限}である．
% P12-ERR-0488: preserve the frozen plural-subject grammar locus.
\item スキームの射 $f : X \to Y$ は {\it 有限型}である．
\item スキーム $X$ は {\it Noether} である．
\item スキーム $X$ 上の ${\mathcal O}_X$-加群 ${\mathcal L}$ は
{\it 可逆}である．
\item 代数閉体上の非特異射影曲線の {\it 種数}．
\end{enumerate}
\end{exercise}

\begin{exercise}
\label{exercises-exercise-kill-global-sections}
$X = \Spec({\mathbf Z}[x, y])$ とし，${\mathcal F}$ を準連接
${\mathcal O}_X$-加群とする．
${\mathcal F}$ の標準アファイン開集合 $D(x)$ への制限が零であると仮定する．
\begin{enumerate}
\item ${\mathcal F}$ の任意の大域切断 $s$ は
$x$ のある冪で消されること，すなわち，ある $n\in {\mathbf N}$ に対して
$x^ns = 0$ となることを示せ．
\item ${\mathcal F}$ が準連接であると仮定しなければ，
同じ主張は成り立つと思うか．
\end{enumerate}
\end{exercise}

\begin{exercise}
\label{exercises-exercise-empty-fibre-empty}
$X \to \Spec(R)$ を固有射とし，$R$ を剰余体 $k$ をもつ
離散付値環とする．
$X \times_{\Spec(R)} \Spec(k)$ が空スキームであると仮定する．
$X$ が空スキームであることを示せ．
\end{exercise}

\begin{exercise}
\label{exercises-exercise-curve-p1-p1}
複素数体 ${\mathbf C}$ 上の射影多様体
\footnote{射影埋め込みは
$((X_0, X_1), (Y_0, Y_1))\mapsto (X_0Y_0, X_0Y_1, X_1Y_0, X_1Y_1)$，
言い換えると $(x, y)\mapsto (1, y, x, xy)$ である．}
$$
{\mathbf P}^1 \times {\mathbf P}^1 = {\mathbf P}^1_{{\mathbf C}}
\times_{\Spec({\mathbf C})} {\mathbf P}^1_{\mathbf C}
$$
を考える．これは，${\mathbf P}^1_{\mathbf C}$ の標準アファイン部分の
組に対応する四つのアファイン部分で被覆される．
たとえば，第一因子に斉次座標 $X_0, X_1$，
第二因子に $Y_0, Y_1$ を用いるとする．
$x = X_1/X_0$，$y = Y_1/Y_0$ とおく．このとき四つの
アファイン開部分は次の環のスペクトルである：
$$
{\mathbf C}[x, y], \quad
{\mathbf C}[x^{-1}, y], \quad
{\mathbf C}[x, y^{-1}], \quad
{\mathbf C}[x^{-1}, y^{-1}].
$$
$X \subset {\mathbf P}^1 \times {\mathbf P}^1$ を，
第一のアファイン部分内の方程式
$$
y^3(x^4 + 1) = x^4 -1.
$$
で与えられる閉部分集合の閉包である閉部分スキームとする．
\begin{enumerate}
\item $X$ が四つのアファイン開部分のうち第一と最後の和に含まれることを示せ．
\item $X$ が非特異射影曲線であることを示せ．
% P12-ERR-0489: preserve the frozen pr_2/first-factor mismatch.
\item 射 $pr_2 : X \to {\mathbf P}^1$（第一因子への射影）を考える．
第一のアファイン部分では，これは写像 $(x, y) \mapsto x$ である．
これが次数 $3$ をもつ理由を簡潔に説明せよ．
\item 写像 $pr_2 : X \to {\mathbf P}^1$ の分岐点と分岐指数を計算せよ．
\item $X$ の種数を計算せよ．
\end{enumerate}
\end{exercise}

\begin{exercise}
\label{exercises-exercise-finite-type-over-Z}
$X \to \Spec({\mathbf Z})$ を有限型の射とする．
$X \times_{\Spec({\mathbf Z})} \Spec({\mathbf F}_p)$ が空でないような
素数 $p$ が無限個存在すると仮定する．
\begin{enumerate}
\item $X \times_{\Spec({\mathbf Z})}\Spec(\mathbf{Q})$
が空でないことを示せ．
\item 条件「有限型」を条件「局所有限型」に置き換えた場合にも，
同じ主張は成り立つと思うか．
\end{enumerate}
\end{exercise}

\section{スキーム論，2009 年春学期末試験}
\label{exercises-section-final-exam-spring-2009}

\noindent
以下は，Columbia 大学で 2009 年春に行われたスキーム論の講義の
期末試験問題である．

\begin{exercise}
\label{exercises-exercise-Noetherian-coherent}
$X$ を Noether スキームとする．
$\mathcal{F}$ を $X$ 上の連接層とする．
$x \in X$ を一点とする．
$\text{Supp}(\mathcal{F}) = \{ x \}$ と仮定する．
\begin{enumerate}
\item $x$ が $X$ の閉点であることを示せ．
\item $H^0(X, \mathcal{F})$ が零でないことを示せ．
\item $\mathcal{F}$ が大域切断で生成されることを示せ．
\item $p > 0$ に対して $H^p(X, \mathcal{F}) = 0$ であることを示せ．
\end{enumerate}
\end{exercise}

\begin{remark}
\label{exercises-remark-invertible-projective-space}
$k$ を体とする．
$\mathbf{P}^2_k = \text{Proj}(k[X_0, X_1, X_2])$ とおく．
$\mathbf{P}^2_k$ 上の任意の可逆層は，ある $n \in \mathbf{Z}$ に対する
$\mathcal{O}_{\mathbf{P}^2_k}(n)$ と同型である．
次を思い出そう：
$$
\Gamma(\mathbf{P}^2_k, \mathcal{O}_{\mathbf{P}^2_k}(n)) =
k[X_0, X_1, X_2]_n
$$
は多項式環の次数 $n$ の部分である．
$\mathbf{P}^2_k$ 上の準連接層 $\mathcal{F}$ に対して，通常どおり
$\mathcal{F}(n) =
\mathcal{F}
\otimes_{\mathcal{O}_{\mathbf{P}^2_k}}
\mathcal{O}_{\mathbf{P}^2_k}(n)$
とおく．
\end{remark}

\begin{exercise}
\label{exercises-exercise-nonsplit-vectorbundle}
$k$ を体とする．
$\mathcal{E}$ を $\mathbf{P}^2_k$ 上のベクトル束，すなわち有限局所自由
$\mathcal{O}_{\mathbf{P}^2_k}$-加群とする．
% P12-ERR-0490: preserve the frozen missing-of source locus.
$\mathcal{E}$ が可逆 $\mathcal{O}_{\mathbf{P}^2_k}$-加群の直和と同型であるとき，
$\mathcal{E}$ は {\it 分裂する}という．
\begin{enumerate}
\item $\mathcal{E}$ が分裂することと $\mathcal{E}(n)$ が分裂することが
同値であることを示せ．
\item $\mathcal{E}$ が分裂するならば，すべての $n \in \mathbf{Z}$ に対して
$H^1({\mathbf{P}^2_k}, \mathcal{E}(n)) = 0$
であることを示せ．
\item 線形形式
$L_0, L_1, L_2 \in \Gamma(\mathbf{P}^2_k, \mathcal{O}_{\mathbf{P}^2_k}(1))$
によって与えられる
$$
\varphi :
\mathcal{O}_{\mathbf{P}^2_k}
\longrightarrow
\mathcal{O}_{\mathbf{P}^2_k}(1) \oplus
\mathcal{O}_{\mathbf{P}^2_k}(1) \oplus
\mathcal{O}_{\mathbf{P}^2_k}(1)
$$
を考える．ある $i$ に対して $L_i \not = 0$ と仮定する．
$\varphi$ の余核がベクトル束となるための
$L_0, L_1, L_2$ に関する条件は何か．
また，それはなぜか．
% P12-ERR-0491: preserve the frozen Given/Give imperative locus.
\item そのような $\varphi$ の例を一つ与えよ．
\item $\Coker(\varphi)$ がベクトル束であるならば，
それは分裂しないことを示せ．
\end{enumerate}
\end{exercise}

\begin{remark}
\label{exercises-remark-recall-dimension-theory}
次の次元論の事実を自由に用いてよい
（必要なら，さらに事実を加えてよい）．
\begin{enumerate}
\item スキームの次元は，既約閉部分集合の鎖の長さの上限である．
\item 体上有限型なスキームの次元は，そのアファイン開集合の
次元の最大値である．
\item Noether スキームの次元は，その既約成分の次元の最大値である．
\item アファインスキームの次元は，対応する環の次元と一致する．
\item $k$ を体とし，$A$ を有限型 $k$-代数とする．
% P12-ERR-0492: preserve the frozen unbound-x source locus.
$A$ が整域で $x \not = 0$ ならば，
$\dim(A) = \dim(A/xA) + 1$ である．
\end{enumerate}
\end{remark}

\begin{exercise}
\label{exercises-exercise-irreducible-fibres-same-dimension-irreducible}
$k$ を体とする．
$X$ を $k$ 上の射影的な被約スキームとする．
$f : X \to \mathbf{P}^1_k$ を $k$ 上のスキームの射とする．
すべての点 $t \in \mathbf{P}^1_k$ に対し，
ファイバー $X_t = f^{-1}(t)$ が次元 $d$ の既約空間となるような
整数 $d \geq 0$ が存在すると仮定する
（既約空間は空でないことを思い出そう）．
\begin{enumerate}
\item $\dim(X) = d + 1$ であることを示せ．
\item $X_0 \subset X$ を次元 $d + 1$ の $X$ の既約成分とする．
すべての $t \in \mathbf{P}^1_k$ に対して
ファイバー $X_{0, t}$ の次元が $d$ であることを証明せよ．
\item 上の結果から $X_t$ と $X_{0, t}$ について何が結論できるか．
\item $X$ が既約であることを示せ．
\end{enumerate}
\end{exercise}

\begin{remark}
\label{exercises-remark-chi}
体 $k$ 上の射影スキーム $X$ と $X$ 上の連接層
$\mathcal{F}$ に対して
$$
\chi(X, \mathcal{F}) =
\sum\nolimits_{i \geq 0} (-1)^i\dim_k H^i(X, \mathcal{F}).
$$
とおく．
\end{remark}

\begin{exercise}
\label{exercises-exercise-complete-intersection}
$k$ を体とする．
$\mathbf{P}^3_k = \text{Proj}(k[X_0, X_1, X_2, X_3])$ と書く．
$C \subset \mathbf{P}^3_k$ を
{\it 型 $(5, 6)$ の完全交差曲線}とする．
これは，
$F \in k[X_0, X_1, X_2, X_3]_5$ と
$G \in k[X_0, X_1, X_2, X_3]_6$ が存在して，
$$
C = \text{Proj}(k[X_0, X_1, X_2, X_3]/(F, G))
$$
が次元 $1$ の多様体となることを意味する
（多様体であることから被約かつ既約であるが，
必要なら $C$ が非特異であると仮定してよい）．
$i : C \to \mathbf{P}^3_k$ を対応する閉埋め込みとする．
完全交差であることから，さらに層の完全列
$$
\xymatrix{
0 \ar[r] &
\mathcal{O}_{\mathbf{P}^3_k}(-11)
\ar[r]^-{
\left(
\begin{matrix}
-G \\
F
\end{matrix}
\right)
} &
\mathcal{O}_{\mathbf{P}^3_k}(-5) \oplus \mathcal{O}_{\mathbf{P}^3_k}(-6)
\ar[r]^-{(F, G)} &
\mathcal{O}_{\mathbf{P}^3_k} \ar[r] &
i_*\mathcal{O}_C \ar[r] &
0
}
$$
も得られる．これらの事実を用いて，次を行え：
\begin{enumerate}
\item 任意の $n \in \mathbf{Z}$ に対して
$\chi(C, i^*\mathcal{O}_{\mathbf{P}^3_k}(n))$ を計算せよ；
\item $H^1(C, \mathcal{O}_C)$ の次元を計算せよ．
\end{enumerate}
\end{exercise}

\begin{exercise}
\label{exercises-exercise-glueing}
$k$ を体とする．
環
\begin{align*}
A & = k[x, y]/(xy) \\
B & = k[u, v]/(uv) \\
C & = k[t, t^{-1}] \times k[s, s^{-1}]
\end{align*}
と，$k$-代数準同型
$$
\begin{matrix}
A \longrightarrow C, &
x \mapsto (t, 0), &
y \mapsto (0, s) \\
B \longrightarrow C, &
u \mapsto (t^{-1}, 0), &
v \mapsto (0, s^{-1})
\end{matrix}
$$
を考える．これらの写像は同型
$A_{x + y} \to C$ および $B_{u + v} \to C$ を誘導する．
したがって，写像 $A \to C$ および $B \to C$ は，
$\Spec(C)$ を $\Spec(A)$ および $\Spec(B)$ の開部分集合と同一視する．
$\Spec(A)$ と $\Spec(B)$ を $\Spec(C)$ に沿って貼り合わせて得られる
スキームを $X$ とする：
$$
X = \Spec(A) \amalg_{\Spec(C)} \Spec(B).
$$
講義で見たように，このようなスキームは存在し，そこには
アファイン開集合 $\Spec(A) \subset X$ と
$\Spec(B) \subset X$ があり，それらの共通部分はちょうど
$\Spec(C)$ であり，上の写像を用いて各々の開部分と同一視される．
\begin{enumerate}
\item なぜ $X$ は分離的か．
\item なぜ $X$ は $k$ 上有限型か．
\item $H^1(X, \mathcal{O}_X)$ を計算せよ．または，その次元はいくつか．
\item $X$ を記述する，より幾何学的な方法は何か．
\end{enumerate}
\end{exercise}

\section{スキーム論，2010 年秋学期末試験}
\label{exercises-section-final-exam-fall-2010}

\noindent
以下は，Columbia 大学で 2010 年秋に行われたスキーム論の講義の
期末試験問題である．

\begin{exercise}[定義]
\label{exercises-exercise-definitions-fall-2010}
次の概念を定義せよ．
\begin{enumerate}
\item 分離スキーム，
\item スキームの準コンパクト射，
\item スキームのアファイン射，
\item 環の乗法的部分集合，
\item Noether スキーム，
\item 多様体．
\end{enumerate}
\end{exercise}

\begin{exercise}
\label{exercises-exercise-prime-avoidance}
素イデアル回避．
\begin{enumerate}
\item $A$ を環とする．$I \subset A$ をイデアルとし，
$\mathfrak q_1$，$\mathfrak q_2$ を
$I \not \subset \mathfrak q_i$ を満たす素イデアルとする．
$I \not \subset \mathfrak q_1 \cup \mathfrak q_2$
であることを示せ．
\item (1) の幾何学的解釈は何か．
\item $X = \text{Proj}(S)$ とし，$S$ はある次数付き環とする．
$x_1, x_2 \in X$ とする．$x_1$ と $x_2$ の両方を含む標準開集合
$D_{+}(F)$ が存在することを示せ．
\end{enumerate}
\end{exercise}

\begin{exercise}
\label{exercises-exercise-compose-affine}
アファイン射の合成がアファインであるのはなぜか．
\end{exercise}

\begin{exercise}[例]
\label{exercises-exercise-examples}
次の例を与えよ：
\begin{enumerate}
\item 体 $k$ 上の可約な射影スキーム．
\item 100 個の点をもつスキーム．
\item スキームのアファインでない射．
\end{enumerate}
\end{exercise}

\begin{exercise}
\label{exercises-exercise-chevalley-hilbert-nullstellensatz}
Chevalley の定理と Hilbert の零点定理．
\begin{enumerate}
\item $\mathfrak p \subset \mathbf{Z}[x_1, \ldots, x_n]$
を極大イデアルとする．Chevalley の定理は
$\mathfrak p \cap \mathbf{Z}$ について何を意味するか．
\item さらに，Hilbert の零点定理は $\kappa(\mathfrak p)$
について何を意味するか．
\end{enumerate}
\end{exercise}

\begin{exercise}
\label{exercises-exercise-P0}
$A$ を環とする．$X$ の次数を $1$ とする次数付き $A$-代数として
$S = A[X]$ とおく．
$A$ 上のスキームとして
$\text{Proj}(S) \cong \Spec(A)$ であることを示せ．
\end{exercise}

\begin{exercise}
\label{exercises-exercise-finite-is-projective}
$A \to B$ を有限な環準同型とする．
$\Spec(B)$ が $\Spec(A)$ 上の H-projective スキームであることを示せ．
\end{exercise}

\begin{exercise}
\label{exercises-exercise-not-geometrically-irreducible}
体 $k$ 上のスキーム $X$ であって，$X$ は既約であるが，
ある有限拡大 $k'/k$ に対する基底変換
$X_{k'} = X \times_{\Spec(k)} \Spec(k')$ は連結かつ可約となるものの
例を与えよ．
\end{exercise}

\section{スキーム論，2011 年春学期末試験}
\label{exercises-section-final-exam-spring-2011}

\noindent
以下は，Columbia 大学で 2011 年春に行われたスキーム論の講義の
期末試験問題である．

\begin{exercise}[定義]
\label{exercises-exercise-definitions-spring-2011}
斜体で示した概念を定義せよ．
\begin{enumerate}
\item {\it 分離的}なスキーム，
\item スキームの {\it 普遍閉}射，
\item 共通の体に含まれる局所環 $A, B$ に対する
{\it $A$ は $B$ を支配する}，
\item スキーム $X$ の {\it 次元}，
\item スキーム $X$ の既約閉部分スキーム $Y$ の {\it 余次元}．
\end{enumerate}
\end{exercise}

\begin{exercise}[結果]
\label{exercises-exercise-results-spring-2011}
講義で論じた事実と形式的に同値な命題を述べよ．
\begin{enumerate}
\item たとえば，多様体の射 $X \to Y$ に対する固有性の付値判定法．
\item 体 $k$ 上の多様体 $X$ に対する
$\dim(X)$ と $X$ の函数体 $k(X)$ の関係．
\item 空欄を埋めよ：$k$ 上の非特異射影曲線と非定数射の圏は，
$\ldots\ldots\ldots$ と反同値である．
\item Noether 正規化．
\item Jacobian 判定法．
\end{enumerate}
\end{exercise}

\begin{exercise}
\label{exercises-exercise-genus-plane-curve}
$k$ を体とする．
$F \in k[X_0, X_1, X_2]$ を次数 $d$ の斉次形式とする．
$C = V_{+}(F) \subset \mathbf{P}^2_k$ が $k$ 上の滑らかな曲線であると仮定する．
% P12-ERR-0493: preserve the frozen missing-by declaration locus.
$i : C \to \mathbf{P}^2_k$ を対応する閉埋め込みとする．
\begin{enumerate}
\item $\mathbf{P}^2_k$ 上の連接層の短完全列
$$
0 \to \mathcal{O}_{\mathbf{P}^2_k}(-d)
\to \mathcal{O}_{\mathbf{P}^2_k} \to i_*\mathcal{O}_C \to 0
$$
が存在することを示せ：写像が何かを述べ，なぜ完全かを簡潔に説明せよ．
\item $H^0(C, \mathcal{O}_C) = k$ であることを結論せよ．
\item $C$ の種数を計算せよ．
\item 今度は $P = (0 : 0 : 1)$ が $C$ 上にないと仮定する．
$(a_0 : a_1 : a_2) \mapsto (a_0 : a_1)$ で与えられる
$\pi : C \to \mathbf{P}^1_k$ の次数が $d$ であることを証明せよ．
\item $k$ は代数閉であり，すべての分岐指数
（「$e_i$」）は $1$ または $2$ であり，$k$ の標数は
$2$ でないと仮定する．
$\pi : C \to \mathbf{P}^1_k$ は分岐点をいくつもつか．
\item $F$ を用いて，$\pi$ の分岐点集合を定める方程式系は
何であると思うか．
\item $K_C$ を推測できるか．
\end{enumerate}
\end{exercise}

\begin{exercise}
\label{exercises-exercise-Pic-triangle}
$k$ を体とする．$X$ を $k$ 上の「三角形」とする．
すなわち，$\mathbf{A}^1_k$ の三つのコピーを，
第一のコピーの $0$ と第二のコピーの $1$，
第二のコピーの $0$ と第三のコピーの $1$，
第三のコピーの $0$ と第一のコピーの $1$ をそれぞれ同一視して
互いに貼り合わせることにより $X$ を得る．
$X$ は $\Spec(k[x, y]/(xy(x + y + 1)))$ と同型になる；
これは自由に用いてよい．
$X$ の Picard 群を計算せよ．
\end{exercise}

\begin{exercise}
\label{exercises-exercise-birational-morphism-curves-ample}
$k$ を体とする．$\pi : X \to Y$ を曲線の有限双有理射とし，
$X$ は $k$ 上の非特異射影曲線とする．
講義の内容から，$Y$ は固有曲線であり，
$\pi$ は $Y$ の正規化射である．
また，稠密開集合 $V \subset Y$ であって，
$U = \pi^{-1}(V)$ は $X$ の稠密開集合であり
$\pi : U \to V$ は同型となるものが存在することも講義で見た．
\begin{enumerate}
\item 有効 Cartier 因子 $D \subset X$ であって，
$D \subset U$ かつ $\mathcal{O}_X(D)$ は $X$ 上豊富となるものが
存在することを示せ．
\item $D$ を (1) のようなものとする．
$E = \pi(D)$ が $Y$ 上の有効 Cartier 因子であることを示せ．
\item 次の理由を簡潔に述べよ：
\begin{enumerate}
\item 写像 $\mathcal{O}_Y \to \pi_*\mathcal{O}_X$ は
$Y \setminus V$ に台をもつ連接な余核 $Q$ をもつ；
\item すべての $n$ に対し，余核が $Q$ と同型となる対応する写像
$\mathcal{O}_Y(nE) \to \pi_*\mathcal{O}_X(nD)$ が存在する．
\end{enumerate}
\item
$\dim_k H^0(X, \mathcal{O}_X(nD)) - \dim_k H^0(Y, \mathcal{O}_Y(nE))$
が有界であること（何によってか）を示し，
$n$ が十分大きいとき可逆層 $\mathcal{O}_Y(nE)$ は多数の切断をもつことを
結論せよ（なぜか）．
\end{enumerate}
\end{exercise}

\section{スキーム論，2011 年秋学期末試験}
\label{exercises-section-final-exam-fall-2011}

\noindent
% P12-ERR-0547: preserve the frozen Schemes-heading/Commutative-Algebra-prose mismatch.
以下は，Columbia 大学で 2011 年秋に行われた可換代数の講義の
期末試験問題である．


\begin{exercise}[定義]
\label{exercises-exercise-definitions-fall-2011}
斜体で示した概念を定義せよ．
\begin{enumerate}
\item {\it Noether} 環，
\item {\it Noether} スキーム，
\item {\it 有限}な環準同型，
\item スキームの {\it 有限}射，
\item 環の {\it 次元}．
\end{enumerate}
\end{exercise}

\begin{exercise}[結果]
\label{exercises-exercise-results-fall-2011}
講義で論じた事実と形式的に同値な命題を述べよ．
\begin{enumerate}
\item Zariski の主定理．
\item Noether 正規化．
\item 中国剰余定理．
\item 有限な環準同型に対する上昇定理．
\end{enumerate}
\end{exercise}

\begin{exercise}
\label{exercises-exercise-dimension-of-ring}
$(A, \mathfrak m, \kappa)$ を，剰余体の標数が $2$ でない
Noether 局所環とする．
$\mathfrak m$ は三つの元 $x, y, z$ で生成され，
$A$ において $x^2 + y^2 + z^2 = 0$ であると仮定する．
\begin{enumerate}
\item $\dim(A)$ の可能な値は何か．
\item 各値が可能であることを示す例を与えよ．
\item $\dim(A) = 2$ ならば $A$ は整域であることを示せ．
（ヒント：
$\bigoplus_{n \geq 0} \mathfrak m^n/\mathfrak m^{n + 1}$
を調べよ．）
\end{enumerate}
\end{exercise}

\begin{exercise}
\label{exercises-exercise-localization}
$A$ を環とする．
$S \subset T \subset A$ を乗法的部分集合とする．
$$
\{\mathfrak q \mid \mathfrak q \cap S = \emptyset\} =
\{\mathfrak q \mid \mathfrak q \cap T = \emptyset\}.
$$
と仮定する．$S^{-1}A \to T^{-1}A$ が同型であることを示せ．
\end{exercise}

\begin{exercise}
\label{exercises-exercise-locus-of-rank-1}
$k$ を代数閉体とする．
$$
V_0 = \{ A \in \text{Mat}(3 \times 3, k) \mid \text{rank}(A) = 1\}
\subset \text{Mat}(3 \times 3, k) = k^9.
$$
とおく．
\begin{enumerate}
\item $V_0$ が，(Zariski) 局所閉部分集合
$V \subset \mathbf{A}^9_k$ の閉点全体であることを示せ．
\item $V$ は既約か．
\item $\dim(V)$ は何か．
\end{enumerate}
\end{exercise}

\begin{exercise}
\label{exercises-exercise-not-complete-intersection}
$\mathbf{C}[x, y]$ のイデアル $(x^2, xy, y^2)$ が
$2$ 個の元では生成できないことを証明せよ．
\end{exercise}

\begin{exercise}
\label{exercises-exercise-finite-projection}
$f \in \mathbf{C}[x, y]$ を非定数多項式とする．
ある $\alpha, \beta \in \mathbf{C}$ に対して
$\mathbf{C}$-代数準同型
$$
\mathbf{C}[t] \longrightarrow \mathbf{C}[x, y]/(f),\quad
t \longmapsto \alpha x + \beta y
$$
が有限であることを示せ．
\end{exercise}

\begin{exercise}
\label{exercises-exercise-union-of-two-affines}
有限個の点 $p_1, \ldots, p_n \in \mathbf{C}^2$ が与えられたとき，
スキーム $\mathbf{A}^2_\mathbf{C} \setminus \{p_1, \ldots, p_n\}$ が
二つのアファイン開集合の和であることを示せ．
\end{exercise}

\begin{exercise}
\label{exercises-exercise-surjection-A1-P1}
スキームの全射
$\mathbf{A}^1_\mathbf{C} \to \mathbf{P}^1_\mathbf{C}$
が存在することを示せ．
（全射とは，単に台となる点集合上で全射であることを意味する．）
\end{exercise}

\begin{exercise}
\label{exercises-exercise-almost-surjective}
$k$ を代数閉体とする．$A \subset B$ を，ともに有限型
$k$-代数である整域の拡大とする．
次の手順を用いて，$\Spec(B) \to \Spec(A)$ の像が
$\Spec(A)$ の空でない開部分集合を含むことを証明せよ：
\begin{enumerate}
\item $A \to B$ も有限である場合に証明せよ．
\item $B$ の分数体が $A$ の分数体の有限拡大である場合に証明せよ．
\item 主張を前の場合に帰着せよ．
\end{enumerate}
\end{exercise}

\section{スキーム論，2013 年秋学期末試験}
\label{exercises-section-final-exam-fall-2013}

\noindent
% P12-ERR-0548: preserve the frozen Schemes-heading/Commutative-Algebra-prose mismatch.
以下は，Columbia 大学で 2013 年秋に行われた可換代数の講義の
期末試験問題である．

\begin{exercise}[定義]
\label{exercises-exercise-definitions-fall-2013}
斜体で示した概念を定義せよ．
\begin{enumerate}
\item 環の {\it 根基イデアル}，
\item {\it 有限型}の環準同型，
\item {\it ヴェイユ微分}，
\item {\it スキーム}．
\end{enumerate}
\end{exercise}

\begin{exercise}[結果]
\label{exercises-exercise-results-fall-2013}
講義で論じた事実と形式的に同値な命題を述べよ．
\begin{enumerate}
\item 次数付き加群の Hilbert 多項式に関する結果．
\item Noether 局所環 $(R, \mathfrak m)$ と
$\bigoplus_{n \geq 0} \mathfrak m^n/\mathfrak m^{n + 1}$ の次元．
\item Riemann--Roch の定理．
\item Clifford の定理．
\item Chevalley の定理．
\end{enumerate}
\end{exercise}

\begin{exercise}
\label{exercises-exercise-surjective-after-localization}
$A \to B$ を環準同型とし，$S \subset A$ を乗法的部分集合とする．
$A \to B$ は有限型であり，$S^{-1}A \to S^{-1}B$ は
全射であると仮定する．
$A_f \to B_f$ が全射となるような $f \in S$ が存在することを示せ．
\end{exercise}

\begin{exercise}
\label{exercises-exercise-injective-local-ring-map-not-surjective}
局所環の単射な局所準同型 $A \to B$ であって，
$\Spec(B) \to \Spec(A)$ が全射でないものの例を与えよ．
\end{exercise}

\begin{situation}[平面曲線の記法]
\label{exercises-situation-curve-in-the-plane}
$k$ を代数閉体とする．
$F(X_0, X_1, X_2) \in k[X_0, X_1, X_2]$ を
次数 $d$ の既約斉次多項式とする．
$$
D = V(F) \subset \mathbf{P}^2
$$
を，$F$ の消滅で与えられる射影平面曲線とする．
$x = X_1/X_0$，$y = X_2/X_0$ とおき，
$f(x, y) = X_0^{-d}F(X_0, X_1, X_2) = F(1, x, y)$ とおく．
% P12-ERR-0494: preserve the frozen missing-by declaration locus.
$K$ を整域 $k[x, y]/(f)$ の分数体とする．
$C$ を $K$ に対応する抽象曲線とする．
講義で見たように，全射 $C \to D$ が存在し，
これは $D$ の非特異点軌跡の上で全単射であり，
$D$ が非特異ならば同型である．
$f_x = \partial f/\partial x$，$f_y = \partial f/\partial y$ とおく．
% P12-ERR-0549: preserve the frozen repeated missing-by denotation loci.
最後に，講義で論じた $\Omega_{K/k}$ の元
$\omega = \text{d}x/f_y = - \text{d}y/f_x$ を考える．
$K_C$ を $\omega$ の零点と極の因子とする．
\end{situation}

\begin{exercise}
\label{exercises-exercise-node-in-the-plane}
状況 \ref{exercises-situation-curve-in-the-plane} において
$d \geq 3$ とし，曲線 $D$ は特異点をちょうど一つ，
すなわち $P = (1 : 0 : 0)$ のみをもつと仮定する．
さらに，$P = (0, 0)$ の近傍で展開
$$
f(x, y) = xy + h.o.t
$$
をもつと仮定する．このとき，$C$ は $P$ の上に二点
$v$ と $w$ をもち，それらは
$$
v(x) = 1, v(y) > 1
\quad\text{and}\quad
w(x) > 1, w(y) = 1
$$
によって特徴付けられる．
\begin{enumerate}
\item $\Omega_{K/k}$ の元
$\omega = \text{d}x/f_y = - \text{d}y/f_x$ が，
$v$ と $w$ の両方で一位の極をもつことを示せ
（$\omega$ の非特異点における振舞いは，講義で論じたとおりである）．
\item 講義では，$C \to D$ によって $C$ へ引き戻した因子
$X_0 = 0$ において，$\omega$ が位数 $d - 3$ で消えることを示した．
(1) の情報と合わせて，$C$ 上の $\omega$ の零点と極の因子の次数は何か．
\item 曲線 $C$ の種数は何か．
\end{enumerate}
\end{exercise}

\begin{exercise}
\label{exercises-exercise-smooth-plane-curve-linear-system}
状況 \ref{exercises-situation-curve-in-the-plane} において $d = 5$ とし，
曲線 $C = D$ は非特異であると仮定する．
講義では，$C$ の種数が $6$ であり，線形系 $K_C$ が
$$
L(K_C) = \{h\omega \mid h \in k[x, y],\ \deg(h) \leq 2\}
$$
で与えられることを示した．ここで $\deg$ は全次数を表す
\footnote{$d - 3 = 5 - 3 = 2$ なので $\leq 2$ を得る．}．
$P_1, P_2, P_3, P_4, P_5 \in D$ を，アファイン開集合
$X_0 \not = 0$ に属する相異なる点とする．
% P12-ERR-0550: preserve the frozen missing-by divisor denotation locus.
$\sum P_i = P_1 + P_2 + P_3 + P_4 + P_5$ を
$C$ の対応する因子とする．
\begin{enumerate}
\item $L(K_C - \sum P_i)$ を多項式を用いて記述せよ．
\item $l(\sum P_i)$ の可能性は何か．
\end{enumerate}
\end{exercise}

\begin{exercise}
\label{exercises-exercise-rational-curve-high-degree}
状況 \ref{exercises-situation-curve-in-the-plane} のような $F$ であって，
$d = 100$ かつ $C$ の種数が $0$ となるものを書き下せ．
\end{exercise}

\begin{exercise}
\label{exercises-exercise-high-degree-curve-quadratic-equation}
$k$ を代数閉体とする．
% P12-ERR-0495: preserve the frozen missing-article source locus.
$K/k$ を超越次数 $1$ の有限生成体拡大とする．
$C$ を $K$ に対応する抽象曲線とする．
$V \subset K$ を一つの $g^r_d$ とし，
$\Phi : C \to \mathbf{P}^r$ を対応する射とする．
$V$ が完全線形系であり，$d$ が $C$ の種数に比べて十分大きければ，
$C$ の像が二次超曲面
\footnote{二次超曲面とは次数 $2$ の超曲面，すなわち
$\mathbf{P}^r$ における次数 $2$ の斉次多項式の零点集合である．}
に含まれることを示せ．
（追加点：必要な次数の良い評価を与えよ．）
\end{exercise}

\begin{exercise}
\label{exercises-exercise-surjective-map-a2-p2}
状況 \ref{exercises-situation-curve-in-the-plane} の記法を用いる．
$U \subset \mathbf{P}^2_k$ を，補集合が $D$ である開部分スキームとする．
$k$-代数 $A = \mathcal{O}_{\mathbf{P}^2_k}(U)$ を記述せよ．
$k$-代数としての $A$ の生成元の個数に上界を与えよ．
\end{exercise}

\section{スキーム論，2014 年春学期末試験}
\label{exercises-section-final-exam-spring-2014}

\noindent
以下は，Columbia 大学で 2014 年春に行われたスキーム論の講義の
期末試験問題である．

\begin{exercise}[定義]
\label{exercises-exercise-definitions-spring-2014}
$(X, \mathcal{O}_X)$ をスキームとする．
斜体で示した概念を定義せよ．
\begin{enumerate}
\item {\it 点 $x$ における $X$ の局所環}，
\item $\mathcal{O}_X$-加群の {\it 準連接}層，
% P12-ERR-0496: preserve the frozen unclosed-parenthesis source locus.
\item $\mathcal{O}_X$-加群の {\it 連接}層
（$X$ は局所 Noether であると仮定してよい），
\item $X$ の {\it アファイン開集合}，
\item {\it スキームの有限射} $X \to Y$．
\end{enumerate}
\end{exercise}

\begin{exercise}[定理]
\label{exercises-exercise-results-spring-2014}
各項目に関係する，講義で論じた非自明な事実を正確に述べよ．
\begin{enumerate}
\item 多様体の多重種数の双有理不変性について，
\item アファイン射であることが局所的性質であること，
\item 射のスキーム論的ファイバーの位相について，
\item 固有性の付値判定法．
\end{enumerate}
\end{exercise}

\begin{exercise}
\label{exercises-exercise-miss-curve}
$\mathbf{C}$ を複素数体とし，$X = \mathbf{A}^2_\mathbf{C}$ とおく．
{\it 直線}とは，ある $a, b, c \in \mathbf{C}$ で
$(a, b) \not = (0, 0)$ を満たすものに対し，
一つの一次方程式 $ax + by + c = 0$ で定義される
$X$ の閉部分スキームを意味するものとする．
{\it 曲線}とは，次元 $1$ の既約な（したがって空でない）閉部分スキーム
$C \subset X$ を意味するものとする．
{\it 二次曲線}とは，ある
$a, b, c, d, e, f \in \mathbf{C}$ で
$(a, b, c) \not = (0, 0, 0)$ を満たすものに対し，
一つの二次方程式
$ax^2 + bxy + cy^2 + dx + ey + f = 0$
で定義される曲線を意味するものとする．
\begin{enumerate}
\item すべての直線と $C$ の共通部分が空でない曲線 $C$ を見つけよ．
\item すべての直線およびすべての二次曲線と $C$ の共通部分が空でない
曲線 $C$ を見つけよ．
\item すべての曲線 $C$ に対し，
$C \cap C' = \emptyset$ となる別の曲線が存在することを示せ．
\end{enumerate}
\end{exercise}

\begin{exercise}
\label{exercises-exercise-normal-bundle-exceptional-curve}
$k$ を体とする．$b : X \to \mathbf{A}^2_k$ を
アファイン平面の原点におけるブローアップとする．
言い換えると，$\mathbf{A}^2_k = \Spec(k[x, y])$ ならば，
$X = \text{Proj}(\bigoplus_{n \geq 0} \mathfrak m^n)$
である．ここで $\mathfrak m = (x, y) \subset k[x, y]$ とする．
次の主張を証明せよ：
\begin{enumerate}
\item 原点上の $b$ のスキーム論的ファイバー $E$ は
$\mathbf{P}^1_k$ と同型である；
\item $E$ は $X$ 上の有効 Cartier 因子である；
\item $\mathcal{O}_X(-E)$ の $E$ への制限は次数 $1$ の直線束である．
\end{enumerate}
（$\mathcal{O}_X(-E)$ は $X$ における $E$ のイデアル層であることを
思い出そう．）
\end{exercise}

\begin{exercise}
\label{exercises-exercise-surjective-map-affine-variety-projective-variety}
$k$ を体とし，$X$ を $k$ 上の射影多様体とする．
$k$ 上のアファイン多様体 $U$ と，多様体の全射
$U \to X$ が存在することを示せ．
\end{exercise}

\begin{exercise}
\label{exercises-exercise-vandermonde}
$k$ を標数 $p > 0$ の体とし，これは $2,3$ と異なるとする．
$\mathbf{P}^n_k$ の閉部分スキーム $X$ を
$$
\sum\nolimits_{i = 0, \ldots, n} X_i = 0,\quad
\sum\nolimits_{i = 0, \ldots, n} X_i^2 = 0,\quad
\sum\nolimits_{i = 0, \ldots, n} X_i^3 = 0
$$
によって定義する．
% P12-ERR-0497: preserve the frozen closed-subsheme/variety terminology locus.
どの組 $(n, p)$ に対してこの閉部分スキームは特異か．
\end{exercise}

\section{可換代数，2016 年秋学期末試験}
\label{exercises-section-final-exam-fall-2016}

\noindent
以下は，Columbia 大学で 2016 年秋に行われた可換代数の講義の
期末試験問題である．

\begin{exercise}[定義]
\label{exercises-exercise-definitions-fall-2016}
$R$ を環とする．斜体で示した概念を定義せよ．
\begin{enumerate}
\item 素イデアル $\mathfrak p$ における $R$ の {\it 局所環}，
\item {\it 有限} $R$-加群，
\item {\it 有限表示} $R$-加群，
\item $R$ は {\it 正則}である，
\item $R$ は {\it 鎖状}である，
\item $R$ は {\it Cohen--Macaulay} である．
\end{enumerate}
\end{exercise}

\begin{exercise}[定理]
\label{exercises-exercise-results-fall-2016}
各項目に関係する，講義で論じた非自明な事実を正確に述べよ．
\begin{enumerate}
\item 正則環，
\item Cohen--Macaulay 加群の随伴素イデアル，
\item 体上有限型な整域の次元，
\item Chevalley の定理．
\end{enumerate}
\end{exercise}

\begin{exercise}
\label{exercises-exercise-finite-injective}
$A \to B$ を次の条件を満たす環準同型とする：
\begin{enumerate}
\item $A$ は極大イデアル $\mathfrak m$ をもつ局所環である；
\item $A \to B$ は有限な
\footnote{これは，$B$ が $A$-加群として有限であることを意味する．}
環準同型である；
\item $A \to B$ は単射である
（$A$ を $B$ の部分環とみなす）．
\end{enumerate}
$\mathfrak m = A \cap \mathfrak q$ となる素イデアル
$\mathfrak q \subset B$ が存在することを示せ．
\end{exercise}

\begin{exercise}
\label{exercises-exercise-find-components}
$k$ を体とし，$R = k[x, y, z, w]$ とおく．
イデアル $I = (xy, xz, xw)$ を考える．
$V(I) \subset \Spec(R)$ の既約成分は何か．また，それらの次元は何か．
\end{exercise}

\begin{exercise}
\label{exercises-exercise-no-nonconstant-morphism}
$k$ を体とし，$A = k[x, x^{-1}]$ および $B = k[y]$ とおく．
任意の $k$-代数準同型 $\varphi : A \to B$ は
$x$ を定数へ写すことを示せ．
\end{exercise}

\begin{exercise}
\label{exercises-exercise-regular-over-Z}
環 $R = \mathbf{Z}[x, y]/(xy - 7)$ を考える．
$R$ が正則であることを証明せよ．
\end{exercise}

\noindent
Noether 局所環 $(R, \mathfrak m, \kappa)$ が与えられたとき，
$n \geq 0$ に対して
$\varphi_R(n) = \dim_\kappa(\mathfrak m^n/\mathfrak m^{n + 1})$
とおく．

\begin{exercise}
\label{exercises-exercise-hilbert-function-allowed}
すべての $n \geq 0$ に対して
$\varphi_R(n) = n + 1$ となる Noether 局所環 $R$ は存在するか．
\end{exercise}

\begin{exercise}
\label{exercises-exercise-hilbert-function-prohibited}
$R$ を Noether 局所環とする．
$\varphi_R(0) = 1$，$\varphi_R(1) = 3$，
$\varphi_R(2) = 5$ と仮定する．
$\varphi_R(3) \leq 7$ であることを示せ．
\end{exercise}

\section{スキーム論，2017 年春学期末試験}
\label{exercises-section-final-exam-spring-2017}

\noindent
以下は，Columbia 大学で 2017 年春に行われたスキーム論の講義の
期末試験問題である．

\begin{exercise}[定義]
\label{exercises-exercise-definitions-spring-2017}
$f : X \to Y$ をスキームの射とする．
斜体で示した概念を簡潔に定義せよ．
\begin{enumerate}
\item $y \in Y$ における $f$ の {\it スキーム論的ファイバー}，
\item $f$ は {\it 有限射}である，
\item {\it 準連接} $\mathcal{O}_X$-加群，
% P12-ERR-0498: preserve the frozen missing-article variety locus.
\item $X$ は {\it 多様体}である，
\item $f$ は {\it 滑らかな射}である，
\item $f$ は {\it 固有射}である．
\end{enumerate}
\end{exercise}

\begin{exercise}[定理]
\label{exercises-exercise-results-spring-2017}
各項目に関係する，講義で論じた非自明な事実を
正確かつ簡潔に述べよ．
\begin{enumerate}
\item 準連接層の順像，
\item 射影多様体上の連接層のコホモロジー，
\item 体上の射影スキームに対する Serre 双対，
\item Riemann--Hurwitz の公式．
\end{enumerate}
\end{exercise}

\begin{exercise}
\label{exercises-exercise-compute-degree}
$k$ を代数閉体とする．
$\ell > 100$ を，$k$ の標数とは異なる素数とする．
$X$ を，$\mathbf{A}^2_k$ における方程式
$$
y^\ell = x(x - 1)^3
$$
で与えられるアファイン曲線の非特異射影モデルとする．
次の問いに答えよ：
\begin{enumerate}
\item $X$ の種数は何か．
\item $X$ のゴナリティ
\footnote{ゴナリティとは，$X$ から $\mathbf{P}^1_k$ への
非定数射の次数の最小値である．}
の上界を一つ与えよ．
\end{enumerate}
\end{exercise}

\begin{exercise}
\label{exercises-exercise-count-components}
$k$ を代数閉体とする．$X$ を $k$ 上の被約射影スキームとし，
そのすべての既約成分は同じ次元 $1$ をもつとする．
$\omega_{X/k}$ を相対双対化加群とする．
$\dim_k H^1(X, \omega_{X/k}) > 1$ ならば
$X$ は非連結であることを示せ．
\end{exercise}

\begin{exercise}
\label{exercises-exercise-sections-L-L-inverse}
スキーム $X$ と非自明な可逆 $\mathcal{O}_X$-加群
$\mathcal{L}$ であって，
$H^0(X, \mathcal{L})$ と $H^0(X, \mathcal{L}^{\otimes -1})$ の
両方が零でないものの例を与えよ．
\end{exercise}

\begin{exercise}
\label{exercises-exercise-in-product}
$k$ を代数閉体とし，$g \geq 3$ とする．
$X$ と $X'$ を，それぞれ種数 $g$ と $g + 1$ の
$k$ 上の滑らかな射影曲線とする．
$Y \subset X \times X'$ を曲線とし，射影
$Y \to X$ と $Y \to X'$ は非定数であるとする．
$Y$ の非特異射影モデルの種数が $\geq 2g + 1$ であることを証明せよ．
\end{exercise}

\begin{exercise}
\label{exercises-exercise-finitely-many}
$k$ を有限体とし，$g > 1$ とする．
次の主張の証明の概略を述べよ：
$k$ 上の種数 $g$ の滑らかな射影曲線の同型類は有限個しかない．
（答えようと試みるだけでも得点を与える．）
\end{exercise}

\section{可換代数，2017 年秋学期末試験}
\label{exercises-section-final-exam-fall-2017}

\noindent
以下は，Columbia 大学で 2017 年秋に行われた可換代数の講義の
期末試験問題である．

\begin{exercise}[定義]
\label{exercises-exercise-definitions-fall-2017}
斜体で示した概念を簡潔に定義せよ．
\begin{enumerate}
\item 関手 $F : \mathcal{A} \to \mathcal{B}$ の {\it 左随伴}，
\item 体の拡大 $L/K$ の {\it 超越次数}，
\item 古典的アファイン多様体 $X \subset k^n$ 上の {\it 正則函数}，
\item 位相空間上の {\it 層}，
\item {\it 局所環}，
\item スキームの射 $f : X \to Y$ が {\it アファイン}であること．
\end{enumerate}
\end{exercise}

\begin{exercise}[定理]
\label{exercises-exercise-results-fall-2017}
各項目に関係する，講義で論じた非自明な事実を
正確かつ簡潔に述べよ
（複数ある場合には，そのうち一つを選べばよい）．
\begin{enumerate}
\item Yoneda の補題，
\item Mayer--Vietoris，
\item 次元とコホモロジー，
\item Hilbert 多項式，
\item 射影空間に対する双対性．
\end{enumerate}
\end{exercise}

\begin{exercise}
\label{exercises-exercise-compute-dimension}
$k$ を代数閉体とする．Zariski 位相と座標
$x_1, x_2, x_3, x_4, x_5$ をもつ $k^5$ の閉部分集合
$X$ を，方程式
$$
x_1^2 - x_4 = 0,\quad
x_2^5 - x_5 = 0,\quad
x_3^2 + x_3 + x_4 + x_5 = 0
$$
によって定める．$X$ の次元は何か．また，それはなぜか．
\end{exercise}

\begin{exercise}
\label{exercises-exercise-can-there-be}
$k$ を体とする．$X = \mathbf{P}^1_k$ を
$k$ 上の次元 $1$ の射影空間とする．
$\mathcal{E}$ を有限局所自由 $\mathcal{O}_X$-加群とする．
% P12-ERR-0551: preserve the frozen missing-by Serre-twist declaration locus.
$d \in \mathbf{Z}$ に対して，
$\mathcal{E}(d) = \mathcal{E} \otimes_{\mathcal{O}_X} \mathcal{O}_X(d)$
を $\mathcal{E}$ の第 $d$ Serre ねじりとし，
$h^i(X, \mathcal{E}(d)) = \dim_k H^i(X, \mathcal{E}(d))$
とおく．
\begin{enumerate}
\item
$h^0(X, \mathcal{E}) = 5$ かつ $h^0(X, \mathcal{E}(1)) = 4$
を満たす $\mathcal{E}$ が存在しないのはなぜか．
\item
$h^1(X, \mathcal{E}(1)) = 5$ かつ $h^1(X, \mathcal{E}) = 4$
を満たす $\mathcal{E}$ が存在しないのはなぜか．
\item どの $a \in \mathbf{Z}$ に対して，$X$ 上のベクトル束
$\mathcal{E}$ であって
$$
\begin{matrix}
h^0(X, \mathcal{E})\phantom{(1)} = 1 &
h^1(X, \mathcal{E})\phantom{(1)} = 1 \\
h^0(X, \mathcal{E}(1)) = 2 &
h^1(X, \mathcal{E}(1)) = 0 \\
h^0(X, \mathcal{E}(2)) = 4 &
h^1(X, \mathcal{E}(2)) = a
\end{matrix}
$$
を満たすものが存在しうるか．
\end{enumerate}
部分的な解答も歓迎し，推奨する．
\end{exercise}

\begin{exercise}
\label{exercises-exercise-banana}
$X$ を，二つの閉部分集合 $Y$ と $Z$ の和
$X = Y \cup Z$ である位相空間とし，
それらの共通部分を $W = Y \cap Z$ と書く．
% P12-ERR-0552: preserve the frozen compressed "Denote ... the inclusion maps" locus.
包含写像を $i : Y \to X$，$j : Z \to X$，
$k : W \to X$ と書く．
\begin{enumerate}
\item 層の短完全列
$$
0 \to \underline{\mathbf{Z}}_X \to
i_*(\underline{\mathbf{Z}}_Y) \oplus
j_*(\underline{\mathbf{Z}}_Z) \to
k_*(\underline{\mathbf{Z}}_W) \to 0
$$
が存在することを示せ．ここで $\underline{\mathbf{Z}}_X$ は
$X$ 上の値 $\mathbf{Z}$ の定数層を表し，他も同様である．
\item $\mathbf{Z}$-係数の $X$，$Y$，$Z$，$W$ の
コホモロジー群の関係について何が結論できるか．
\end{enumerate}
\end{exercise}

\begin{exercise}
\label{exercises-exercise-cohomology-infinite-punctured}
$k$ を体とする．$A = k[x_1, x_2, x_3, \ldots]$ を
無限個の変数の多項式環とする．
% P12-ERR-0553: preserve the frozen missing-by maximal-ideal declaration locus.
$\mathfrak m$ を，すべての変数で生成される $A$ の極大イデアルとする．
$X = \Spec(A)$，$U = X \setminus \{\mathfrak m\}$ とおく．
\begin{enumerate}
\item $H^1(U, \mathcal{O}_U) = 0$ であることを示せ．
ヒント：{\v C}ech コホモロジーによる計算．
\item $i \geq 1$ に対する $H^i(U, \mathcal{O}_U)$ は
何であると予想するか．
\end{enumerate}
\end{exercise}

\begin{exercise}
\label{exercises-exercise-principal}
$A$ を局所環とし，$a \in A$ を非零因子とする．
$I, J \subset A$ を $IJ = (a)$ を満たすイデアルとする．
イデアル $I$ が単項，すなわち一つの元で生成されることを示せ
（その元は非零因子となる）．
\end{exercise}

\section{スキーム論，2018 年春学期末試験}
\label{exercises-section-final-exam-spring-2018}

\noindent
以下は，Columbia 大学で 2018 年春に行われたスキーム論の講義の
期末試験問題である．

\begin{exercise}[定義]
\label{exercises-exercise-definitions-spring-2018}
斜体で示した概念を簡潔に定義せよ．
$k$ を代数閉体とする．
$X$ を $k$ 上の射影曲線とする．
\begin{enumerate}
\item $k$ 上の {\it 滑らかな}代数，
\item $X$ 上の可逆 $\mathcal{O}_X$-加群の {\it 次数}，
\item $X$ の {\it 種数}，
\item $X$ の {\it Weil 因子類群}，
\item $X$ は {\it 超楕円的}である，および
\item $k$ 上の滑らかな射影曲面上の二曲線の {\it 交点数}．
\end{enumerate}
\end{exercise}

\begin{exercise}[定理]
\label{exercises-exercise-results-spring-2018}
各項目に関係する，講義で論じた非自明な事実を
正確かつ簡潔に述べよ
（複数ある場合には，そのうち一つを選べばよい）．
\begin{enumerate}
\item Riemann--Hurwitz の定理，
\item Clifford の定理，
\item 滑らかな射影曲面の間の射の分解，
\item Hodge の指数定理，および
\item 有限体上の曲線に対する Riemann 予想．
\end{enumerate}
\end{exercise}

\begin{exercise}
\label{exercises-exercise-hyperelliptic}
$k$ を代数閉体とする．$X \subset \mathbf{P}^3_k$ を
次数 $d$，種数 $\geq 2$ の滑らかな曲線とする．
$X$ は平面に含まれず，かつ $\mathbf{P}^3_k$ 内に
$X$ と $d - 2$ 個の点で交わる直線 $\ell$ が存在すると仮定する．
$X$ が超楕円的であることを示せ．
\end{exercise}

\begin{exercise}
\label{exercises-exercise-singular}
$k$ を代数閉体とする．$X$ を，相異なる特異点
$p_1, \ldots, p_n$ をもつ射影曲線とする．
$X$ の正規化の種数が高々
$-n + \dim_k H^1(X, \mathcal{O}_X)$ である理由を説明せよ．
\end{exercise}

\begin{exercise}
\label{exercises-exercise-how-many-blowups}
$k$ を体とする．$X = \Spec(k[x, y])$ をアファイン $2$ 次元空間とする．
$$
I = (x^3, x^2y, xy^2, y^3) \subset k[x, y].
$$
とおく．
$Y \subset X$ を $I$ に対応する閉部分スキームとする．
$b : X' \to X$ をイデアル $(x, y)$ のブローアップ，すなわち
原点におけるアファイン空間のブローアップとする．
\begin{enumerate}
\item スキーム論的逆像 $b^{-1}Y \subset X'$ が
有効 Cartier 因子であることを示せ．
% P12-ERR-0499: preserve the frozen "Given an example" imperative locus.
\item $I \subset J \subset (x, y)$ を満たすイデアル
$J \subset k[x, y]$ であって，$J$ に対応する閉部分スキームを
$Z \subset X$ とすると，スキーム論的逆像 $b^{-1}Z$ が
有効 Cartier 因子とならないものの例を与えよ．
\end{enumerate}
\end{exercise}

\begin{exercise}
\label{exercises-exercise-rational-surface}
$k$ を代数閉体とする．次の型の曲面を考える：
\begin{enumerate}
\item $C_1$ と $C_2$ が滑らかな射影曲線である
$S = C_1 \times C_2$，
\item $C_1$ と $C_2$ が滑らかな射影曲線で，かつ $C_1$ の種数が
$> 0$ である $S = C_1 \times C_2$，
\item 次数 $4$ の超曲面 $S \subset \mathbf{P}^3_k$，および
\item 次数 $4$ の滑らかな超曲面 $S \subset \mathbf{P}^3_k$．
\end{enumerate}
各型について，その型の曲面の類が有理曲面を含むか否かと，その理由を
簡潔に述べよ．
\end{exercise}

\begin{exercise}
\label{exercises-exercise-self-square-line}
$k$ を代数閉体とする．$S \subset \mathbf{P}^3_k$ を
次数 $d$ の滑らかな超曲面とする．$S$ が直線 $\ell$ を含むと仮定する．
% P12-ERR-0500: preserve the frozen nonstandard "self square" source phrase.
$S$ 上の因子とみた $\ell$ の自己交点数はいくつか．
\end{exercise}

\section{可換代数，2019 年秋学期末試験}
\label{exercises-section-final-exam-fall-2019}

\noindent
以下は，Columbia 大学で 2019 年秋に行われた可換代数の講義の
期末試験問題である．

\begin{exercise}[定義]
\label{exercises-exercise-definitions-fall-2019}
斜体で示した概念を簡潔に定義せよ．
\begin{enumerate}
\item Noether 位相空間の {\it 構成可能部分集合}，
\item 素イデアル $\mathfrak p$ における $R$-加群 $M$ の {\it 局所化}，
\item Noether 局所環 $(A, \mathfrak m, \kappa)$ 上の加群の {\it 長さ}，
\item 環 $R$ 上の {\it 射影加群}，および
\item Noether 局所環 $(A, \mathfrak m, \kappa)$ 上の
{\it Cohen--Macaulay} 加群．
\end{enumerate}
\end{exercise}

\begin{exercise}[定理]
\label{exercises-exercise-results-fall-2019}
各項目に関係する，講義で論じた非自明な事実を
正確かつ簡潔に述べよ
（複数ある場合には，そのうち一つを選べばよい）．
\begin{enumerate}
\item 構成可能集合の像，
\item Hilbert の零点定理，
\item 体上有限型な代数の次元，
\item Noether の正規化，および
\item 正則局所環．
\end{enumerate}
\end{exercise}

\noindent
環 $R$ とイデアル $I \subset R$ に対して，$V(I)$ は
$I \subset \mathfrak p$ を満たす $\mathfrak p \in \Spec(R)$ の集合を
表すことを思い出そう．

\begin{exercise}[素イデアルの構成]
\label{exercises-exercise-infinitely-many-primes}
$V(\mathfrak p)$ が $(x, y)$ と $(x - 1, y - 1)$ を含むような
相異なる素イデアル $\mathfrak p \subset \mathbf{C}[x, y]$ を
無限個構成せよ．
\end{exercise}

\begin{exercise}[素イデアルなし]
\label{exercises-exercise-no-prime}
$R = \mathbf{C}[x, y, z]/(xy)$ とする．
$V(\mathfrak p)$ が $(x, y - 1, z - 5)$ と $(x - 1, y, z - 7)$ を
含むような素イデアル $\mathfrak p \subset R$ は存在しないことを
簡潔に論ぜよ．
\end{exercise}

\begin{exercise}[Frobenius]
\label{exercises-exercise-frobenius}
$p$ を素数とする（式を簡単にするため $p = 2$ と仮定してよい）．
環 $R$ を取り，$R$ において $p = 0$ と仮定する．
\begin{enumerate}
\item 写像 $F : R \to R$，$x \mapsto x^p$ が環準同型であることを示せ．
\item $\Spec(F) : \Spec(R) \to \Spec(R)$ が恒等写像であることを示せ．
\end{enumerate}
\end{exercise}

\noindent
位相空間の点の {\it 特殊化} $x \leadsto y$ とは，単に $y$ が
$x$ の閉包に属することを意味することを思い出そう．
$x$ および $y$ と異なる $z$ で，$x \leadsto z$ かつ
$z \leadsto y$ となるものが存在しないとき，$x \leadsto y$ を
{\it 直前特殊化}という．

\begin{exercise}[次元]
\label{exercises-exercise-dimension}
$5$ 個の相異なる点 $x, y, z, u, v$ を含む sober 位相空間 $X$ があり，
その特殊化は次のとおりであるとする：
$$
\xymatrix{
x \ar[d] \ar[r] & u & v \ar[l] \ar[dl] \\
y \ar[r] & z
}
$$
このような $X$ がもちうる最小の次元はいくつか．
$X$ が体上有限型な代数のスペクトルであり，$x \leadsto u$ が
直前特殊化であるならば，特殊化 $v \leadsto z$ について何がいえるか．
\end{exercise}

\begin{exercise}[Tor の計算]
\label{exercises-exercise-tor-computation}
$R = \mathbf{C}[x, y, z]$ とする．
$M = R/(x, z)$ および $N = R/(y, z)$ とおく．
どの $i \in \mathbf{Z}$ に対して $\text{Tor}_i^R(M, N)$ は零でないか．
\end{exercise}

\begin{exercise}
\label{exercises-exercise-depth-goes-up}
$A \to B$ を Noether 局所環の平坦な局所準同型とする．
$A$ の深さが $k$ ならば，$B$ の深さは $k$ 以上であることを示せ．
\end{exercise}

\section{代数幾何，2020 年春学期末試験}
\label{exercises-section-final-exam-spring-2020}

\noindent
以下は，Columbia 大学で 2020 年春に行われた代数幾何の講義の
期末試験問題である．

\begin{exercise}[定義]
\label{exercises-exercise-definitions-spring-2020}
斜体で示した概念を簡潔に定義せよ．
\begin{enumerate}
\item {\it スキーム}，
\item {\it スキームの射}，
\item スキーム上の {\it 準連接加群}，
\item 体 $k$ 上の {\it 多様体}，
\item 体 $k$ 上の {\it 曲線}，
\item スキームの {\it 有限射}，
\item 位相空間 $X$ 上の Abel 群の層 $\mathcal{F}$ の
{\it コホモロジー}，
\item 体 $k$ 上固有で次元 $d$ のスキーム $X$ 上の
{\it 双対化層}，および
\item 多様体 $X$ から多様体 $Y$ への {\it 有理写像}．
\end{enumerate}
\end{exercise}

\begin{exercise}[定理]
\label{exercises-exercise-results-spring-2020}
各項目に関係する，講義で論じた非自明な事実を
正確かつ簡潔に述べよ
（複数ある場合には，そのうち一つを選べばよい）．
\begin{enumerate}
\item 次元 $d$ の Noether 位相空間 $X$ 上の Abel 層のコホモロジー，
\item 体 $k$ 上の滑らかな多様体の微分層 $\Omega^1_{X/k}$，
\item 体 $k$ 上の滑らかな射影多様体 $X$ の双対化層 $\omega_X$，
\item 代数閉体 $k$ 上の滑らかで固有な種数 $0$ の曲線，および
\item 次数 $d$ の平面曲線の種数．
\end{enumerate}
\end{exercise}

\begin{exercise}
\label{exercises-exercise-coh-P1-infty-doubled}
$k$ を体とする．$X$ を $k$ 上のスキームとする．
$X = X_1 \cup X_2$ は開被覆であり，$X_1$ と $X_2$ はともに
$\mathbf{P}^1_k$ と同型で，$X_1 \cap X_2$ は
$\mathbf{A}^1_k$ と同型であると仮定する
（このようなスキームは存在する．例えば $\mathbf{P}^1_k$ の
$\infty$ を二重化すればよい）．
$\dim_k H^1(X, \mathcal{O}_X)$ が無限であることを示せ．
\end{exercise}

\begin{exercise}
\label{exercises-exercise-no-morphism}
$k$ を代数閉体とする．$Y$ を種数 $10$ の滑らかな射影曲線とする．
同型でない非定数射 $f : X \to Y$ が存在するような滑らかな射影曲線
$X$ の種数について，良い下界を求めよ．
\end{exercise}

\begin{exercise}
\label{exercises-exercise-element-order-n-in-pic}
$k$ を標数 $0$ の代数閉体とする．
$$
X : T_0^d + T_1^d - T_2^d = 0 \subset \mathbf{P}^2_k
$$
を次数 $d \geq 3$ の Fermat 曲線とする．
$X$ 上の閉点 $p = [1 : 0 : 1]$ と $q = [0 : 1 : 1]$ を考え，
$D = [p] - [q]$ とおく．
\begin{enumerate}
\item $D$ が Weil 因子類群において非自明であることを示せ．
\item $d D$ が Weil 因子類群において自明であることを示せ．
（ヒント：$d[p]$ と $d[q]$ がともに $X$ と平面内のある直線との
交わりであることを示してみよ．）
\end{enumerate}
\end{exercise}

\begin{exercise}
\label{exercises-exercise-number-equations}
$k$ を代数閉体とする．$2$ 次 Veronese 埋め込み
$$
\varphi : \mathbf{P}^2 \longrightarrow \mathbf{P}^5
$$
を考える．講義の内容と記法では，これは射
$$
\varphi = \varphi_{\mathcal{O}_{\mathbf{P}^2}(2)} :
\mathbf{P}^2
\longrightarrow
\mathbf{P}(\Gamma(\mathbf{P}^2, \mathcal{O}_{\mathbf{P}^2}(2)))
$$
である．斉次座標では，これは
$$
[a_0 : a_1 : a_2] \longmapsto
[a_0^2 : a_0a_1 : a_0a_2 : a_1^2 : a_1a_2 : a_2^2]
$$
で与えられる．これは閉埋め込みである（この事実はそのまま用いてよい）．
$I \subset k[T_0, \ldots, T_5]$ を
$\varphi(\mathbf{P}^2)$ の斉次イデアルとする．すなわち，斉次部分
$I_d$ の元は，閉部分スキーム $\varphi(\mathbf{P}^2)$ 上で
零に制限される次数 $d$ の斉次多項式 $F(T_0, \ldots, T_5)$ である．
$\dim_k I_d$ を $d$ の函数として計算せよ．
\end{exercise}

\begin{exercise}
\label{exercises-exercise-dualizing-sheaf-rank-1}
$k$ を代数閉体とする．$X$ を，$k$ 上の次元 $d$ の固有スキームで，
双対化加群 $\omega_X$ をもつものとする．
次の情報が与えられている：
\begin{enumerate}
\item 任意の $i$ と任意の連接 $\mathcal{O}_X$-加群
$\mathcal{F}$ に対して，
$\text{Ext}^i_X(\mathcal{F}, \omega_X) \times
H^{d - i}(X, \mathcal{F}) \to H^d(X, \omega_X) \xrightarrow{t} k$
は非退化である，および
\item $\omega_X$ はある階数 $r$ の有限局所自由加群である．
\end{enumerate}
$r = 1$ であることを示せ
（ヒント：閉点に台をもつ適切な加群 $\mathcal{F}$ を取ったときに
何が起こるか調べよ）．
\end{exercise}

\section{可換代数，2021 年秋学期末試験}
\label{exercises-section-final-exam-fall-2021}

\noindent
以下は，Columbia 大学で 2021 年秋に行われた可換代数の講義の
期末試験問題である．

\begin{exercise}[定義]
\label{exercises-exercise-definitions-fall-2021}
斜体で示した概念を簡潔に定義せよ．
\begin{enumerate}
\item 環 $A$ の {\it 乗法的部分集合}，
\item {\it Artin 環} $A$，
\item 位相空間としての環 $A$ の {\it スペクトル}，
\item {\it 平坦な環準同型} $A \to B$，
\item $A$ の素イデアル $\mathfrak p$ の {\it 高さ}，および
\item 環 $A$ 上の関手 {\it $\text{Tor}^A_i(-, -)$}．
\end{enumerate}
\end{exercise}

\begin{exercise}[定理]
\label{exercises-exercise-results-fall-2021}
各項目に関係する，講義で論じた非自明な事実を
正確かつ簡潔に述べよ
（複数ある場合には，そのうち一つを選べばよい）．
\begin{enumerate}
\item Artin 環，
\item 平坦性と素イデアル，
\item Noether 局所環 $(A, \mathfrak m)$ に対する
$A/\mathfrak m^n$ の長さ，
\item 普遍鎖状 Noether 環に対する次元公式，
\item Noether 局所環の完備化，および
\item Artin 局所環に対する Matlis 双対性．
\end{enumerate}
\end{exercise}

\begin{exercise}[単元]
\label{exercises-exercise-simple-units}
Abel 群として，$\mathbf{Z}[x, 1/x]$ の単元群はどのような構造を
もつか．説明は不要である．
\end{exercise}

\begin{exercise}[イデアル]
\label{exercises-exercise-ideals}
$A = \mathbf{F}_2[x, y]/(x^2, xy, y^2)$ とする．
% P12-ERR-0554: preserve the frozen missing-by notation clause.
$\overline{x}$ と $\overline{y}$ をそれぞれ $A$ における
$x$ と $y$ の像とする．$A$ のイデアルを列挙せよ．説明は不要である．
\end{exercise}

\begin{exercise}[Tor と Ext]
\label{exercises-exercise-compute-tor-ext}
$(A, \mathfrak m, \kappa)$ を Noether 局所環とする．
$\varphi(n) = \dim_\kappa \mathfrak m^n/\mathfrak m^{n + 1}$ とおく．
\begin{enumerate}
\item $\text{Tor}_1^A(A/\mathfrak m^n, \kappa)$ が
$\kappa$-ベクトル空間として次元 $\varphi(n)$ をもつことを示せ．
\item $\text{Ext}^1_A(A/\mathfrak m^n, \kappa)$ が
$\kappa$-ベクトル空間として次元 $\varphi(n)$ をもつことを示せ．
\end{enumerate}
\end{exercise}

\begin{exercise}[二つのベクトル]
\label{exercises-exercise-two-vectors}
% P12-ERR-0501: preserve the frozen b3 generator token.
$A = \mathbf{Z}[a_1, a_2, a_3, b_1, b_2, b3]$ とする．
$A^{\oplus 3}$ において $a = (a_1, a_2, a_3)$ および
$b = (b_1, b_2, b_3)$ とおく．集合
$$
Z = \{\mathfrak p \in \Spec(A) \mid a, b
\text{ は次のベクトル空間で線形従属な像をもつ：}
\kappa(\mathfrak p)^{\oplus 3}\}
$$
を考える．
\begin{enumerate}
% P12-ERR-0502: preserve the frozen extraneous-article source locus.
\item $Z$ が $\Spec(A)$ の閉部分集合であることを証明せよ．
\item $Z$ の次元 $\dim(Z)$ はいくつか．
\item $\mathbf{Z}$ を体で置き換えると，$\dim(Z)$ はどうなるか．
\end{enumerate}
\end{exercise}

\begin{exercise}[入射的加群]
\label{exercises-exercise-injective-artinian-local}
$(A, \mathfrak m, \kappa)$ を Artin 局所環とする．
$A$ は $A$-加群として入射的であると仮定する．
% P12-ERR-0503: preserve the frozen malformed as/has source phrase.
$\Hom_A(\kappa, A)$ が $\kappa$-ベクトル空間として次元 $1$ を
もつことを示せ．
\end{exercise}

\section{代数幾何，2022 年春学期末試験}
\label{exercises-section-final-exam-spring-2022}

\noindent
以下は，Columbia 大学で 2022 年春に行われた代数幾何の講義の
期末試験問題である．

\begin{exercise}[定義]
\label{exercises-exercise-definitions-spring-2022}
斜体で示した概念を簡潔に定義せよ．
\begin{enumerate}
\item {\it スキーム}，
\item スキーム $X$ 上の {\it 準連接加群}，
\item スキームの {\it 平坦射} $X \to Y$，
\item スキームの {\it 有限射} $X \to Y$，
\item 基底スキーム $S$ 上の {\it 群スキーム} $G$，
\item 基底スキーム $S$ 上の {\it 多様体族}，
\item 体 $k$ 上の多様体 $X$ の閉点 $x$ の {\it 次数}，
\item $\mathbf{P}^n(\mathbf{Q})$ の点
$p = (a_0 : \ldots : a_n)$ の通常の {\it 対数的高さ}，および
\item {\it $C_i$ 体}．
\end{enumerate}
\end{exercise}

\begin{exercise}[定理]
\label{exercises-exercise-results-spring-2022}
各項目に関係する，講義で論じた非自明な事実を
正確かつ簡潔に述べよ
（複数ある場合には，そのうち一つを選べばよい）．
\begin{enumerate}
\item スキーム $X$ からアファインスキーム $\Spec(A)$ への射，
\item アファインスキーム $X$ 上の準連接加群 $\mathcal{F}$ の
コホモロジー，
\item $k$ が体であるときの $\mathbf{P}^1_k$ の Picard 群，
\item $S$ が Noether であるときの平坦かつ固有な射
$X \to S$ のファイバーの次元，
\item スキーム $S$ 上の $\mathbf{G}_m$-同変加群，および
\item 交わりに関する Bezout の定理
（望むなら特別な場合に限ってよい）．
\end{enumerate}
\end{exercise}

\begin{exercise}[三次超曲面]
\label{exercises-exercise-cubics}
$F \in \mathbf{C}[T_0, \ldots, T_n]$ を次数 $3$ の斉次式とする．
$3$ 個のベクトル $x, y, z \in \mathbf{C}^{n + 1}$ が与えられたとき，
条件
$$
(*)\quad
F(\lambda x + \mu y + \nu z) = 0 \text{ が次の多項式環で成り立つ：} \mathbf{C}[\lambda, \mu, \nu]
$$
を考える．
\begin{enumerate}
\item $x, y, z$ の選び方全体の空間の次元はいくつか．
% P12-ERR-0555: preserve the frozen subject-verb agreement locus.
\item 条件 (*) は $x$，$y$，$z$ の座標に何本の方程式を課すか．
\item (*) が成り立つ三つ組 $x, y, z$ 全体の空間の期待次元はいくつか．
\item $x, y, z$ が線形従属である三つ組全体の空間の次元はいくつか．
\item $\mathbf{P}^n$ 内の次数 $3$ の超曲面には，$n \geq a$ ならば
次元 $2$ の線形部分空間が存在すると期待されることを結論せよ．
ここで $a$ は自ら求めよ．
\end{enumerate}
\end{exercise}

\begin{exercise}[高さ]
\label{exercises-exercise-heights}
$K$ を体とする．講義で論じた $2$ 個の公理を満たす函数の族
$h_n : \mathbf{P}^n(K) \to \mathbf{R}$，$n \geq 0$ を取る．
$X$ を $K$ 上の射影多様体とする．
$\mathcal{L}$ を可逆 $\mathcal{O}_X$-加群とし，講義では
付随する高さ函数 $h_\mathcal{L} : X(K) \to \mathbf{R}$ を
構成したことを思い出そう．
$\alpha : X \to X$ を $K$ 上の $X$ の自己同型とする．
\begin{enumerate}
\item $P \mapsto h_\mathcal{L}(\alpha(P))$ は函数
$h_{\alpha^*\mathcal{L}}$ と有界な量しか違わないことを証明せよ．
（ヒント：$\mathcal{L} = \varphi^*\mathcal{O}_{\mathbf{P}^n}(1)$ を
満たす射 $\varphi : X \to \mathbf{P}^n$ が存在するとき，構成により
$h_\mathcal{L}(P) = h_n(\varphi(P))$ であることを思い出し，
これを使って式をいろいろ変形せよ．一般には $\mathcal{L}$ を
この形の二つの差として書け．）
\item $h_\mathcal{L}(P) - h_\mathcal{L}(\alpha(P))$ が
$X(K)$ 上で非有界であると仮定する．
$\mathcal{N} = \mathcal{L} \otimes \alpha^*\mathcal{L}^{\otimes -1}$
に対する $h_\mathcal{N}$ が $X(K)$ 上で非有界であることを示せ．
\item $X$ は楕円曲線であり，$\mathcal{L}$ は $X$ 上の
対称かつ豊富な可逆加群で，$h_\mathcal{L}$ が $X(K)$ 上で
非有界であると仮定する．次数 $0$ の可逆加群 $\mathcal{N}$ であって
$h_\mathcal{N}$ が非有界となるものが存在することを示せ．
（ヒント：$X$ は次元 $1$ の Abel 多様体であることを思い出そう．
したがって講義の結果により，$h_\mathcal{L}$ は定数差を除いて
二次的である．適切な点 $P_0 \in X(K)$ を選べ．
$\alpha : X \to X$ を $P_0$ による平行移動とする．
$P \mapsto h_\mathcal{L}(P) - h_\mathcal{L}(P + P_0)$ を考え，
上で証明した結果を適用せよ．）
\end{enumerate}
\end{exercise}

\begin{exercise}[モノ射]
\label{exercises-exercise-monomorphisms}
$f : X \to Y$ をスキームの圏におけるモノ射とする．
すなわち，任意の二つのスキームの射 $a, b : T \to X$ に対し，
$f \circ a = f \circ b$ ならば $a = b$ であるとする．
$f$ が点上で単射であることを示せ．
% P12-ERR-0504: preserve the frozen "Does you argument" source locus.
この議論から他に何か分かるか．
\end{exercise}

\begin{exercise}[不動点]
\label{exercises-exercise-fix-points}
$k$ を代数閉体とする．
\begin{enumerate}
\item $G = \mathbf{G}_{m, k}$ とする．$G$ が $k$ 上の射影多様体
$X$ に作用するならば，その作用は不動点をもつことを示せ．
すなわち，すべての $g \in G(k)$ に対して $a(g, x) = x$ となる点
$x \in X(k)$ が存在することを証明せよ．
\item $G = (\mathbf{G}_{m, k})^n$ が乗法群の $n \geq 1$ 個の
直積である場合にも同じことを示せ．
\item 滑らかな射影多様体 $X$ 上の連結群スキーム $G$ の作用で，
不動点をもたないものの例を与えよ．
\end{enumerate}
\end{exercise}

\section{代数幾何，2025 年春学期末試験}
\label{exercises-section-final-exam-spring-2025}

\noindent
以下は，Columbia 大学で 2025 年春に行われた代数幾何の講義の
期末試験問題である．

\begin{exercise}[定義]
\label{exercises-exercise-definitions-spring-2025}
斜体で示した概念を簡潔に定義せよ．
\begin{enumerate}
\item {\it スキーム} $X$，
\item スキーム $X$ 上の {\it 準連接加群}，
\item スキーム $X$ の {\it Picard 群}，
% P12-ERR-0505: preserve the frozen "give a morphism" source locus.
\item スキームの射 $f : X \to Y$ と $\mathcal{O}_X$-加群
$\mathcal{F}$ が与えられたときの {\it 順像} $f_*\mathcal{F}$，
\item スキームの {\it 閉埋め込み}，および
\item 与えられた体 $k$ 上の {\it 多様体}．
\end{enumerate}
\end{exercise}

\begin{exercise}[定理]
\label{exercises-exercise-results-spring-2025}
% P12-ERR-0506: preserve the frozen "succintly" spelling locus.
各項目に関係する，講義で論じた非自明な事実を一つ
正確かつ簡潔に述べよ
（複数ある場合には，そのうち一つを選べばよい）．
\begin{enumerate}
\item スキーム $X$ の位相空間，
\item スキームの圏上の関手に対する表現可能性判定法，
\item 代数閉体 $k$ 上の滑らかな射影曲線 $X$ の Picard 関手，
\item 代数閉体 $k$ 上の滑らかな射影曲線 $X$ の Picard 群の捩れ，および
\item 代数閉体 $k$ 上の滑らかな射影多様体 $X$，$Y$，$Z$ の積
$X \times Y \times Z$ の Picard 群に関する定理．
\end{enumerate}
\end{exercise}

\begin{exercise}
\label{exercises-exercise-affine-line-infinite-nr-points}
$k$ を体とする．$\Spec(k[x])$ が無限個の点をもつ理由を説明せよ．
\end{exercise}

\begin{exercise}
\label{exercises-exercise-hypersurface-variety}
$k$ を体とし，$x_1, \ldots, x_n$ を変数とする．
$f \in k[x_1, \ldots, x_n]$ を非定数多項式とする．
$X = \Spec(k[x_1, \ldots, x_n]/(f))$ とおく．
$X$ が $k$ 上の多様体となるためには，$f$ はどのような性質を
もつべきか．
\end{exercise}

\begin{exercise}
\label{exercises-exercise-find-singular-point}
$\mathbf{C}$ 上の多様体とみた
$\mathbf{C}[x, y]/(x^5 - 5xy^4 + 20y - 16)$ のスペクトル上の
特異点を一つ求めよ（すべて求める必要はない）．
\end{exercise}

\begin{exercise}
\label{exercises-exercise-kernel-restriction-pic-curve}
$X$ を代数閉体 $k$ 上の滑らかな射影曲線とする．
$x \in X(k)$ を $X$ 上の閉点とし，$U = X \setminus \{x\}$ とおく．
制限写像 $\Pic(X) \to \Pic(U)$ の核について何がいえるか．
\end{exercise}

\begin{exercise}
\label{exercises-exercise-sections-degree-2g-minus-4}
$X$ を，種数 $g \geq 100$ の，代数閉体 $k$ 上の滑らかな射影曲線とする．
% P12-ERR-0507: preserve the frozen "for L be" source locus.
$X$ 上の次数 $2g - 4$ の可逆 $\mathcal{O}_X$-加群
$\mathcal{L}$ に対し，
$h^0(\mathcal{L}) = \dim_k H^0(X, \mathcal{L})$ はどのような値を
取りうるか．できるだけ完全に答えよ．
\end{exercise}
